\documentclass[11pt]{article}
\usepackage{amsmath,amssymb,amsthm,hyperref}
% Reuse cached script fonts in the desktop compiler, including nested indices.
\DeclareMathSizes{10}{10}{8}{8}
\DeclareMathSizes{10.95}{10.95}{8}{8}
\setlength{\textwidth}{6.5in}
\setlength{\oddsidemargin}{0in}
\setlength{\evensidemargin}{0in}
\setlength{\textheight}{9in}
\setlength{\topmargin}{-0.5in}
\newtheorem{theorem}{Theorem}[section]
\newtheorem{lemma}[theorem]{Lemma}
\newtheorem{corollary}[theorem]{Corollary}
\newtheorem{proposition}[theorem]{Proposition}
\theoremstyle{remark}
\newtheorem{remark}[theorem]{Remark}
\newcommand{\one}{\mathbf1}
\newcommand{\R}{\mathbb R}
\newcommand{\E}{\mathbb E}
\newcommand{\supp}{\operatorname{supp}}
\newcommand{\diag}{\operatorname{diag}}
\newcommand{\e}{\mathrm e}
\newcommand{\phiG}{\varphi}
\title{Signed compensation and relative Hamiltonian enumeration\\
in arithmetic graphs}
\author{Research draft}
\date{November 14, 2026}
\begin{document}
\maketitle
\begin{abstract}
We develop two Hamiltonian counting results for stochastic kernels
with a fixed spectral gap, allowing zero entries.
For diffuse symmetric kernels, including balanced bipartite kernels,
Hamilton cycles and every admissible prescribed-endpoint spanning
path have total weight with logarithm
\[
 -N+O_{\delta,C}(N\alpha^{1/8}\log(2N)),
 \qquad \alpha=\max P_{ij}.
\]
A sufficient condition is
\(\alpha\leq C(\log(2N))^{-30,000,000}\).
The finite ingredients are simple exact-regular spectral rounding,
observable-preserving sampling, complementary-cut control, and
the Brada\v{c}--Janzer spanning-path theorem.

For denser kernels we obtain relative asymptotics by tracking
growing entry bounds in the Gaussian permanent expansion and
retaining the signed determinant--permanent identity.
For bipartite cross blocks at logarithmic density, an explicit
minor expansion and signed occupancy derivatives give a relative
comparison whenever
\(\det((I+BB^\top)/(I-BB^\top))\leq m^{2-\varepsilon}\).
In particular a fixed gap and \(A=o(\log m)\) suffice.
The Hamilton-cycle weight is
\[
 (1+o(1))\frac{(m!)^2}{2m\,m^{2m}}
                  =(1+o(1))\pi e^{-2m}.
\]
Cauchy--Binet cancels the complete nontrivial spectral determinant,
including mismatched row and column supports.
An aperiodic relative formula holds when \(A^2=o(\log N)\).
For any fixed number of directed cyclic layers, we retain the
complete product determinant at the larger range \(A=o(\log m)\).
The two-layer symmetric case is exactly the cancellation case.

Two arithmetic consumers are independent of RH.
Prime-sum graphs are Hamiltonian at every sufficiently large even
order and have the full relative count
\[
 h(G_N)=\pi e^{2\mathcal H_N-N}(1+O(N^{-3/5})),
\]
where \(\mathcal H_N\) is their exact proper fractional-edge entropy.
The arithmetic input is the existing signed prime operator transfer;
no prime-pair asymptotic is assumed.
We also obtain complete fixed-congruence classification at the
exponential scale and a residue-profile relative-entropy law.
For coprime circles through \(\mathbb Z/n\mathbb Z\), we determine
the full relative count uniformly over all growing moduli:
\[
 h(X_n)=(1+O(n^{-1/10}))(\varphi(n)/e)^n
 \begin{cases}
 \pi,& n\ {\rm even},\\
 \mathfrak C(n)\sqrt{\pi/(2n)},& n\ {\rm odd}.
 \end{cases}
\]
The odd arithmetic factor has an explicit uniformly convergent
prime-product expansion.
Filz's all-orders finite remainder is not settled.
These are ordinary paper proofs; Lean certification and independent
verification remain separate.
\end{abstract}

\section{Statements and the complete consumer}
Let \(N=2m\), with fixed parts \(A,B\) of size \(m\).
All matrices in the rounding statements act on Euclidean space.
Write \(s=\one_A-\one_B\) and
\[
 V_0=\{\one,s\}^{\perp},\qquad L=\log(2N).
\]
A stochastic bipartite kernel \(P\) is a real symmetric nonnegative
matrix with zero within-part blocks and \(P\one=\one\).
Thus \(Ps=-s\).
Its nontrivial gap is at least \(\delta\in(0,1)\) if
\(\|P|_{V_0}\|\leq1-\delta\).
Put \(\alpha=\max_{ij}P_{ij}\).
Throughout, \(\e(t)=\exp(2\pi\mathrm{i}t)\).
Here a support graph has an edge precisely at a positive off-diagonal
entry; it has no loops.

\begin{theorem}[Simple regular spectral rounding]\label{thm:rounding}
For each fixed \(\delta\in(0,1)\), there are constants
\(c_\delta,C_\delta>0\) such that the following holds.
If \(P\) has gap at least \(\delta\), \(1\leq d\leq m\)
is an integer, and
\[
 d\alpha L^2\leq c_\delta,
\]
then its support contains a simple \(d\)-regular spanning graph \(H\)
satisfying
\[
 \left\|\frac{A_H}{d}-P\right\|
 \leq C_\delta\left(
      d\alpha L+\sqrt{\frac{L}{d}}+\frac{L}{d}\right).
 \tag{1}
\]
The constants are uniform in \(P,N,d\).
They are not numerically optimized.
\end{theorem}

The simplicity conclusion requires controlling the forbidden edges
after every matching; independent Birkhoff samples alone would
produce a regular multigraph.

We use the following established result, with its precise bipartite
scope. By Brada\v{c}--Janzer~\cite[Corollary 1.4]{BJ}, for sufficiently
large \(N\), a simple \(D\)-regular bipartite graph with
\(\lambda_2\leq(1-\eta)D\) is Hamiltonian whenever
\[
 D>(3\eta^{-1}\log N)^K,\qquad K=10^8.                       \tag{2}
\]
Here \(\lambda_2\) is the algebraically second eigenvalue; the
bipartite eigenvalue \(-D\) is permitted.
The cited expansion definition applies up to \(2N/3\);
a spectral gap \(\eta\) supplies its edge-expansion constant
\(\eta/3\).

\begin{corollary}[Diffuse expanding kernels have Hamiltonian support]
\label{cor:weightedhamilton}
Fix \(a,\delta\in(0,1)\).
For all sufficiently large even \(N\), every stochastic bipartite
kernel with nontrivial gap at least \(\delta\) and
\(\alpha\leq N^{-a}\) has Hamiltonian support.
More generally there is \(b_\delta>0\) such that, for sufficiently
large \(N\), the same conclusion holds when
\[
 \alpha\leq b_\delta L^{-(K+3)}.                              \tag{3}
\]
Consequently a non-Hamiltonian support with this fixed gap must
have an entry larger than the right side of (3).
\end{corollary}

Let \(G_N\) have integer vertices \(1,\ldots,N\) and edges
\(ij\) for distinct \(i,j\) when \(i+j\) is prime.

\begin{theorem}[Eventual Filz and a regular spectral subgraph]
\label{thm:filz}
For every sufficiently large even \(N\), \(G_N\) contains a simple
\(d_N\)-regular spanning bipartite graph \(H_N\), where
\[
 d_N=\lfloor N^{1/4}\rfloor,\qquad
 \lambda_2(H_N)\leq(1/2+o(1))d_N.                            \tag{4}
\]
In particular \(G_N\) is Hamiltonian for every sufficiently large
even \(N\).
\end{theorem}

Filz's conjecture asks for Hamiltonicity at every even order at
least four. Chen, Fu and Guo~\cite{CFG} proved infinitely many cases.
Theorem~\ref{thm:filz} asserts an eventual statement and leaves a finite,
unbounded-in-this-draft remainder. We do not claim the complete
all-orders conjecture.
Shang--Li--Zhang~\cite{SLZ} construct an infinite prime circle,
a two-way ordering of all positive integers. That result concerns
an infinite graph, rather than every sufficiently large finite order. The arithmetic proof is unconditional and
uses no RH or QRH hypothesis.

The chain proving (4) and the Hamiltonicity conclusion is:
a signed prime-convolution operator limit; uniform prime-weighted
degrees; an exact stochastic scaling with controlled error;
Theorem~\ref{thm:rounding}; and finally the established theorem (2).
All row defects, forbidden-edge losses and martingale drift are
estimated below.

The uniform unweighted degrees allow a stronger integral-factor
consumer in Section~\ref{sec:regularfactors}: the degree improves
from \(N^{1/4}\) to \((1-o(1))N/\log N\).
The same section proves a positive-density packing of Hamilton
cycles and explicit local resilience.
Section~\ref{sec:filters} classifies all fixed congruence filters
by their exact finite Fourier obstruction.
Section~\ref{sec:paths} verifies all hypotheses for prescribed
spanning paths, including arbitrary parity-compatible endpoints.
Section~\ref{sec:enumeration} proves the unweighted counting
theorem with any fixed nontrivial spectral gap.
Section~\ref{sec:weightedentropy} gives weighted universality
for diffuse stochastic bipartite kernels.
Sections~\ref{sec:pathentropy}--\ref{sec:generalkernels} extend
it to prescribed spanning paths and general symmetric
aperiodic kernels at all orders, with quantitative error
\[
 \log\sum_{C\ {\rm Hamilton}}\prod_{ij\in E(C)}P_{ij}
       =-N+O_{\delta,C}(N\alpha^{1/8}\log(2N)).
\]
The resulting Sinkhorn entropy formula applies to sparse
nonuniform graphs, beyond a bipartite restriction.
Section~\ref{sec:coprimecircles} supplies an arithmetic
counting application with the modulus itself growing.
Section~\ref{sec:residuepressure} determines the partition
function under every fixed positive prime-residue tilt and
derives its full empirical relative-entropy law.
In particular
\[
 \log h(G_N)=N\log N-N\log\log N-N+o(N).
\]
The same formula includes \(N\log\sigma_R\) for every
unobstructed fixed congruence filter.
The prime-sum counting error is exponential in \(o(N)\).
Sections~\ref{sec:relativepermanents}--\ref{sec:relativecoprime}
prove relative formulas on the stated denser-kernel domains,
and in particular a full relative count for coprime circles.
Those initial estimates alone do not cover the prime-sum
density boundary. Sections~\ref{sec:secondpermanent}--\ref{sec:relativeprimecircles}
close that range by retaining the full first-order minor correction
and bounding its remainder only after the signed sum is evaluated.
This yields a power-saving relative count around the exact graph
entropy, without replacing that entropy by its coarse asymptotic.
Sections~\ref{sec:periodiccomparison}--\ref{sec:layeredcoprime}
extend the signed comparison to every fixed number of cyclic
layers, including noncommuting cross blocks.
The arithmetic consumer classifies all parity obstructions and
computes the higher-layer prime interactions explicitly.

\section{Exact scaling under a controlled row defect}
\label{sec:scaling}
For a vector \(x\), write \(\|x\|_\infty=\max_i|x_i|\);
matrix operator norms are Euclidean unless specified.
Let \(s=\one_A-\one_B\).

\begin{lemma}[A quantitative bipartite scaling repair]
\label{lem:scaling}
Fix \(\delta\in(0,1)\).
Let \(Q\) be a symmetric nonnegative matrix supported between
two parts of equal size, and put
\[
 d=Q\one,\qquad D=\diag(d),\qquad
 \|d-\one\|_\infty\leq\varepsilon<1/2,\qquad
 S=D^{-1/2}QD^{-1/2}.
\]
Assume all eigenvalues of \(S\) except its simple eigenvalues
\(1,-1\) have modulus at most \(1-\delta\).
There are constants \(c_\delta,C_\delta>0\) such that if
\(\varepsilon L^2\leq c_\delta\), there is a positive vector
\(u\) with
\[
 P=\diag(u)Q\diag(u),\quad P\one=\one,\qquad
 \|u-\one\|_\infty\leq C_\delta\varepsilon L,\qquad
 \|P-Q\|\leq C_\delta\varepsilon L.                           \tag{5}
\]
The support of \(P\) is exactly the support of \(Q\).
\end{lemma}
\begin{proof}
The bipartite symmetry gives
\(\sum_A d_i=\sum_Bd_i\).
Set \(T=D^{-1}Q\), a reversible stochastic matrix with stationary
measure \(\pi_i=d_i/\sum_jd_j\).
It is similar to \(S\), and
\(\pi(A)=\pi(B)=1/2\).
Its eigenvectors at \(1,-1\) are \(\one,s\).
The corresponding projections are
\[
 \Pi_+f=\one\sum_i\pi_i f_i,\qquad
 \Pi_-f=s\sum_i\pi_i s_i f_i,\qquad
 \Pi_0=I-\Pi_+-\Pi_-.
\]
For \(j\geq0\), the assumptions and \(d_i\in[1/2,3/2]\) give
\[
 \|T^j\Pi_0\|_{\infty\to\infty}
 \leq C\min\{1,\sqrt N(1-\delta)^j\}.                         \tag{6}
\]
For the first bound, use \(T^j\Pi_0=T^j-\Pi_+-(-1)^j\Pi_-\)
and the three bounded stochastic/projection norms.
For the second, use the \(L^2(\pi)\) spectral gap and
\(\min_i\pi_i\geq1/(3N)\).

Let \({\cal L}_Q=D+Q=D(I+T)\).
Its kernel is \(\R s\), by the hypothesis at \(-1\).
For \(r\perp s\), the vector \(f=D^{-1}r\) has
\(\Pi_-f=0\).
The convergent series
\[
 z_0=\tfrac12\Pi_+f+
                 \sum_{j\geq0}(-1)^jT^j\Pi_0f
\]
therefore solves \({\cal L}_Qz_0=r\).
The sum of the bounds in (6) is \(O_\delta(L)\).
Subtract the \(s\)-component of \(z_0\) to impose \(z\perp s\).
Thus the inverse on \(s^\perp\) satisfies
\[
 \|{\cal L}_Q^\dagger r\|_\infty
                    \leq C_\delta L\|r\|_\infty.             \tag{7}
\]
Here the dagger means the inverse with this fixed gauge;
no inverse is asserted on the one-dimensional kernel.

Put \(u=\one+z\).
The exact row equation, including its nonlinear cross term, is
\[
 u\odot(Qu)-\one
 =(d-\one)+{\cal L}_Qz+z\odot(Qz).                           \tag{8}
\]
Both \(d-\one\) and \(z\odot(Qz)\) are orthogonal to \(s\).
For the latter,
\[
 \sum_{i\in A}z_i(Qz)_i
 =\sum_{i\in A,j\in B}Q_{ij}z_i z_j
 =\sum_{j\in B}z_j(Qz)_j.
\]
For the former use equal part sizes and equal total degrees.
Consequently the map
\[
 z\longmapsto-{\cal L}_Q^\dagger
                    \bigl((d-\one)+z\odot(Qz)\bigr)
\]
maps \(s^\perp\) to itself.
By (7), it is a contraction on the ball
\(\|z\|_\infty\leq2C_\delta\varepsilon L\) when
\(c_\delta\) is sufficiently small.
Indeed \(\|Q\|_{\infty\to\infty}\leq3/2\), and the nonlinear
Lipschitz constant on a ball of radius \(r\) is at most \(3r\);
the contraction cost is \(O_\delta(\varepsilon L^2)\).
The map also preserves this ball by the same inequality.
The fixed point gives (8) equal to zero.
Choose \(c_\delta\) smaller so that \(\|z\|_\infty<1/2\);
then every \(u_i\) is positive.

Finally, with \(Z=\diag(z)\), retain all terms
\[
 P-Q=ZQ+QZ+ZQZ.
\]
Since \(\|Q\|\leq3/2\), their operator norm is
\(O_\delta(\varepsilon L)\), proving (5).
\end{proof}

\begin{corollary}[Repair after forbidden edges]
\label{cor:forbiddenrepair}
Let \(P\) satisfy the hypotheses of Theorem~\ref{thm:rounding},
and let \(F\) be an arbitrary set of its support edges.
Define the actual removed row mass
\[
 \eta=\max_i\sum_{j:ij\in F}P_{ij}.
\]
If \(\eta L^2\leq c_\delta\), then there is a symmetric
stochastic \(P_F\) supported exactly on the unused edges, with
\[
 \|P_F-P\|\leq C_\delta\eta L,\qquad
 \max_{ij}(P_F)_{ij}\leq2\alpha.                             \tag{9}
\]
It has nontrivial gap at least \(\delta/2\).
\end{corollary}
\begin{proof}
Set \(Q_{ij}=P_{ij}\one_{ij\notin F}\).
Then \(Q\one\) lies between \(1-\eta\) and 1.
Symmetry and the row bound give
\(\|Q-P\|\leq\eta\).
For \(D=\diag(Q\one)\), it follows that
\[
 \|D^{-1/2}QD^{-1/2}-P\|\leq C\eta.
\]
The normalized matrix has eigenvalues \(1,-1\), and Weyl's
inequality bounds every other eigenvalue in modulus by
\(1-\delta+C\eta\).
Thus Lemma~\ref{lem:scaling} applies with gap \(\delta/2\).
Its exact positive scaling preserves all unused entries and
has multipliers \(1+O_\delta(\eta L)\).
This gives (9).
After scaling, the constant and parity vectors are again exact
eigenvectors. The operator perturbation in (9), and a further
decrease of \(c_\delta\), give the claimed gap.
\end{proof}

\section{Adaptive spectral rounding and the Hamiltonicity consumer}
\label{sec:rounding}
We first recall the elementary decomposition used at each step.
If \(T\) is a doubly stochastic matrix, its positive support has
a perfect matching: a set of rows \(U\) sends mass \(|U|\)
into columns of total capacity \(|N(U)|\), so Hall's condition holds.
Choose such a permutation matrix \(M\) and let \(w>0\) be the
smallest matched entry.
Then
\[
 T=wM+(1-w)\frac{T-wM}{1-w}.
\]
If \(w=1\) the decomposition is complete; otherwise the second
matrix is doubly stochastic and has fewer positive entries.
Iteration gives a finite convex combination of permutation
matrices. In symmetric bipartite block form, these are adjacency
matrices of perfect matchings. In particular, for a stochastic
bipartite kernel \(P_i\) there is a random perfect matching \(M_i\)
with \(\E M_i=P_i\). This decomposition is an established input,
not a new claim.

\paragraph{Formal finite matching selector.}
The classical finite selection input has now been proved in Lean in
a form that can also drive a deterministic moment argument.
For every nonnegative doubly stochastic cross matrix \(P\) and
every real cost matrix \(C\), there is a permutation \(f\) with
\(P_{i,f(i)}>0\) for every \(i\) such that its permutation matrix
\(M_f\) satisfies the signed inequality
\[
 \sum_{i,j} C_{ij}\bigl((M_f)_{ij}-P_{ij}\bigr)\leq0.
\]
For completeness, choose a supported matching using the exact Hall
mass count above.  Let \(w>0\) be its smallest matched entry.
If \(w=1\), stochasticity forces \(P=M_f\).
Otherwise
\[
 R=\frac{P-wM_f}{1-w},\qquad P=wM_f+(1-w)R
\]
is nonnegative and doubly stochastic, its positive support is contained
in that of \(P\), and at least one matched entry has disappeared.
Induct on the number of positive entries.
If \(M_f\) already has cost at most that of \(P\), use it.
Otherwise the displayed convex identity forces \(R\) to have
strictly smaller cost than \(P\), and the induction hypothesis on
\(R\) supplies a supported matching with no larger cost.
This proves the signed conclusion without a probability assumption.
The maintained owners are \texttt{HamiltonianStochasticMatching},
\texttt{HamiltonianStochasticPeeling} and
\texttt{HamiltonianMatchingCostSelection}.
This is the established finite assignment principle; no originality
is claimed for it.  The full one-step matrix-moment estimate, finite adaptive
construction and same-graph normalized spectral and observable
readouts are now proved in Lean, as detailed below.

\paragraph{Formal trace-power estimate for deterministic selection.}
The next analytic input is now also proved in Lean, uniformly in the
matrix order and the moment order.
For real symmetric \(A,X\), put \(A_t=A+tX\) and
\(g_p(t)=\operatorname{tr}(A_t^p)\), where \(p\geq2\) is an integer.
Scalar differentiation of every matrix entry gives the exact formulas
\[
 g_p'(t)=p\operatorname{tr}(A_t^{p-1}X),\qquad
 g_p''(t)=p\sum_{j=0}^{p-2}
          \operatorname{tr}(A_t^{p-2-j}XA_t^jX).
\]
The matrix power direction retains every noncommuting word.
Diagonalize \(A_t\) only after differentiating:
write \(A_t=U\diag(\lambda_r)U^\top\) and \(B=U^\top XU\).
No derivative of the eigenbasis is used.
Every mixed trace is exactly the full ordered-pair sum
\[
 \operatorname{tr}(A_t^aXA_t^bX)
   =\sum_{r,s}B_{rs}^{\,2}\lambda_r^a\lambda_s^b.
\]
In particular all signed eigenvalues and return pairs remain present.
The actual operator norm gives both budgets
\[
 \sum_s B_{rs}^{\,2}\leq\|X\|^2,\qquad
 \sum_r B_{rs}^{\,2}\leq\|X\|^2.
\]
For \(x,y\geq0\) and \(0\leq j\leq n\),
\(x^jy^{n-j}\leq x^n+y^n\).
Apply this with \(n=p-2\) to each complete mixed trace,
then sum using the two square-mass budgets.  This proves
\[
 \boxed{\quad
 g_p''(t)\leq
 2p(p-1)\|X\|^2\sum_r|\lambda_r(A_t)|^{p-2}.
 \quad}
\]
The bound uses the eigenvalue mass instead of replacing it by the
ambient order times a maximum entry.
The maintained proof consists of
\texttt{HamiltonianTracePowerScalars},
\texttt{HamiltonianTraceSpectralBasis},
\texttt{HamiltonianTracePowerDerivative} and
\texttt{HamiltonianTraceHessianBudget}.
It is a supporting classical matrix estimate; no originality is
claimed for the estimate in isolation.
The eigenvalue-mass interpolation and the complete regularized
squared-moment step are now also proved in Lean.
For every finite nonnegative family \(a_r\) and \(0\leq q\leq1\),
\[
 \sum_r a_r^q\leq n^{1-q}\left(\sum_r a_r\right)^q.
\]
The proof raises the power-mean inequality with exponent \(1/q\)
back to the power \(q\), treating \(q=0\) separately.
It includes empty finite sets.
Apply it to \(a_r=|\lambda_r|^p\) and \(q=1-2/p\).
For even \(p\geq2\), the exact Hermitian trace identity gives
\[
 \sum_r|\lambda_r|^{p-2}
 \leq n^{2/p}\bigl(\operatorname{tr}A^p\bigr)^{1-2/p},
 \qquad \operatorname{tr}A^p\geq0.
\]
Now write \(g(t)=\operatorname{tr}(A+tX)^p\),
\(q=2/p\), and \(h(t)=g(t)+\varepsilon>0\), where
\(\varepsilon>0\).
The full differentiated expression is
\[
 \frac{d^2}{dt^2}h(t)^q
 =g''(t)q h(t)^{q-1}
  +(g'(t))^2q(q-1)h(t)^{q-2}.
\]
The second term is nonpositive, since \(0<q\leq1\).
In the first term the Hessian bound and the interpolation above give
\(g''(t)\leq K g(t)^{1-q}\), with
\(K=2p(p-1)n^{2/p}\|X\|^2\).
Using \(0\leq g(t)\leq h(t)\), the exact power cancellation
\(h^{1-q}h^{q-1}=1\) yields
\[
 \frac{d^2}{dt^2}h(t)^q
 \leq4(p-1)n^{2/p}\|X\|^2.
\]
This is uniform in \(A\) and \(\varepsilon>0\).
The one-sided Taylor theorem is proved by applying derivative
monotonicity first to \(f'(t)-Bt\), then to
\(f(t)-f'(0)t-Bt^2/2\); it needs an upper bound on \(f''\),
rather than an absolute second-derivative bound.
Consequently, if
\(\operatorname{tr}(A^{p-1}X)\leq0\), then
\[
 \boxed{\quad
 \bigl(\operatorname{tr}((A+X)^p)+\varepsilon\bigr)^{2/p}
 \leq
 \bigl(\operatorname{tr}A^p+\varepsilon\bigr)^{2/p}
 +2(p-1)n^{2/p}\|X\|^2 .
 \quad}
\]
These proofs are maintained in
\texttt{HamiltonianPowerMassInterpolation},
\texttt{HamiltonianTraceMomentMass},
\texttt{HamiltonianRegularizedPower},
\texttt{HamiltonianSquaredTraceCurvature},
\texttt{HamiltonianOneSidedTaylor} and
\texttt{HamiltonianSquaredTraceStep}.
They are classical analytic supporting estimates.
The complete analytic finite-path iteration and its operator-norm
readout are now also proved.
Suppose \(Z_0=0\), \(Z_{i+1}=Z_i+X_i\), the matrices are
Hermitian, \(\|X_i\|\leq2\), and
\(\operatorname{tr}(Z_i^{p-1}X_i)\leq0\) for every \(i<d\).
The previous step, summed with its full constants, gives
\[
 \bigl(\operatorname{tr}(Z_d^p)+\varepsilon\bigr)^{2/p}
 \leq \varepsilon^{2/p}+8(p-1)d\,n^{2/p}.
\]
To read the actual operator norm, the orthogonal spectral
decomposition gives
\(\|A\|\leq\|(\lambda_r)_r\|_\infty\), with the norm of the
empty vector equal to zero.
Every nonnegative even eigenvalue power is bounded by the full
trace sum, so
\[
 \|A\|^2\leq\bigl(\operatorname{tr}(A^p)\bigr)^{2/p}.
\]
Both conclusions include the empty finite index type.
Taking \(\varepsilon=1\), rather than discarding the initial
regularization, yields the compiled finite-path estimate
\[
 \boxed{\quad \|Z_d\|^2\leq1+8(p-1)d\,n^{2/p}.\quad}
\]
The two additional maintained owners are
\texttt{HamiltonianTraceNormReadout} and
\texttt{HamiltonianFiniteTracePath}.
This is the actual operator norm, with the dimension occurring only
as \(n^{2/p}\).
The path theorem explicitly requires the signed direction condition
at every step; no conditional concentration hypothesis is used.
The actual local matching step for the joint packet is now also
proved in Lean.
For a cross matrix \(C\), let
\[
 B(C)=\begin{pmatrix}0&C\\C^\top&0\end{pmatrix},\qquad
 {\cal L}(C)=\operatorname{diag}
 \bigl(B(C),\diag(B(C)\one_{S_1}),\diag(B(C)\one_{S_2})\bigr).
\]
Here \(S_1,S_2\) are arbitrary subsets of the full bipartite vertex set.
The packet has order \(n=3N=6m\) for an \(m\)-by-\(m\) cross matrix.
Its first block is exactly \(B(C)\), and its other two block
diagonals are exactly the requested row observables.
The map \({\cal L}\) is real linear and has Hermitian values.
If \(P\) is nonnegative and doubly stochastic, every row of \(B(P)\)
has absolute mass one.
Each observable diagonal entry is bounded in modulus by that same
row mass.
The established symmetric row-budget theorem therefore gives
\[
 \|{\cal L}(P)\|\leq1.
\]
The same proof controls the lift by any uniform absolute row-mass
budget; it uses the full repaired-kernel subset estimate.

For any current Hermitian packet \(Z\), the complete trace
functional
\(\phi(C)=\operatorname{tr}(Z^{p-1}{\cal L}(C))\)
is linear.
With \(E_{ij}\) the single-entry matrix units, its exact finite
coordinate representation is
\[
 \phi(C)=\sum_{i,j}\phi(E_{ij})C_{ij}.
\]
Thus the already proved arbitrary signed-cost matching selector
produces a supported permutation matrix \(M\) with
\[
 X={\cal L}(M-P),\qquad
 X=X^\top,\qquad \|X\|\leq2,\qquad
 \operatorname{tr}(Z^{p-1}X)\leq0.
\]
The norm bound uses linearity and the two actual stochastic-kernel
bounds \(\|{\cal L}(M)\|,\|{\cal L}(P)\|\leq1\).
Every hypothesis needed by the finite moment path theorem has
therefore been discharged for this actual local supported matching
step, for the same unrestricted \(S_1,S_2\).
The maintained owners are
\texttt{HamiltonianLinearMatchingDirection},
\texttt{HamiltonianJointObservableLift} and
\texttt{HamiltonianBipartiteRoundingStep}.
These are classical supporting constructions.
The exact repair-and-freshness step is now also proved in Lean.
Suppose the current cross adjacency matrix \(H\) has entries zero
or one, every row and column sums to \(i\), and it is supported
on the original \(P\).
Define the actual deleted mass entrywise by
\[
 Q_{ab}=P_{ab}H_{ab}.
\]
If \(0\leq P_{ab}\leq\alpha\), then
\[
 0\leq Q\leq P,\qquad
 \sum_b Q_{ab}\leq i\alpha,\qquad
 \sum_a Q_{ab}\leq i\alpha.
\]
These are exact mask bounds, with no independence or disjoint
pairing assumption.
Under \(i\alpha L^2\leq c_\delta\), apply the full repair theorem
to the original \(P\), yielding a stochastic \(W\) with
\[
 W_{ab}=0\ \Longleftrightarrow\
 (P-Q)_{ab}=P_{ab}(1-H_{ab})=0,\qquad
 \|{\cal L}(W-P)\|\leq C_\delta^{\rm obs}i\alpha L.
\]
Select the actual joint signed matching using \(W\).
Its positive support forces \(H_{a,f(a)}=0\) and
\(P_{a,f(a)}>0\).
Thus \(H+M_f\) is still binary, is supported on \(P\), and
has every row and column sum exactly \(i+1\).
The same witness supplies all stochastic, support, drift and
joint-moment direction properties.
Every repair uses the original kernel and its original gap;
the argument does not repeatedly halve the gap of successive repairs.
The maintained owners are \texttt{HamiltonianUsedEdgeMass} and
\texttt{HamiltonianUsedEdgeJointStep}.
The complete finite-prefix induction and its quantitative error are
now also proved in Lean.
The maintained prefix records the actual binary supported \(H_i\),
the whole centered matrix path, every Hermitian increment and signed
direction, and a drift matrix satisfying
\[
 {\cal L}(H_i-iP)=Z_i+R_i,\qquad
 \|R_i\|\leq C_\delta^{\rm obs}\alpha L\,i^2 .
\]
At extension, append \(X_i={\cal L}(M_{i+1}-W_i)\) and set
\[
 Z_{i+1}=Z_i+X_i,\qquad
 R_{i+1}=R_i+{\cal L}(W_i-P).
\]
Every earlier path entry is retained.
The exact linear matrix identity proves the new error identity.
The full repair bound gives
\[
 \|R_{i+1}\|
 \leq C_\delta^{\rm obs}\alpha L(i^2+i)
 \leq C_\delta^{\rm obs}\alpha L(i+1)^2.
\]
The global smallness hypothesis implies each prefix repair's
smallness hypothesis.
Fix an even \(p\geq2\) before this construction.
Induction starts with zero adjacency, centered path and drift.
It constructs an actual final binary supported cross matrix
\(H=H_p\), with every row and column sum \(d\).
The witness may depend on the chosen moment order.
Applying the compiled path moment estimate to this same \(H_p\)
gives
\[
 \boxed{\quad
 \|{\cal L}(H-dP)\|
 \leq\sqrt{1+8(p-1)d(6m)^{2/p}}
       +C_\delta^{\rm obs}\alpha Ld^2.
 \quad}
\]
The owners are \texttt{HamiltonianRoundingPrefix},
\texttt{HamiltonianRoundingPrefixExtension} and
\texttt{HamiltonianRoundingConstruction}.
The full actual finite construction is proved, rather than assumed
as a favorable matrix path.
The logarithmic moment choice and the final same-graph readouts are
now also proved in Lean.
Put \(n=6m\) and choose the even integer
\[
 p=2\lceil2\log n\rceil.
\]
For \(m\geq1\), the scalar proof gives
\[
 p\geq2,\qquad p-1\leq5\log n,\qquad
 n^{2/p}\leq2,\qquad
 1\leq L=\log(4m),\quad \log n\leq2L.
\]
Consequently, for every integer \(d\geq1\),
\[
 1+8(p-1)d\,n^{2/p}\leq256dL.
\]
The initial unit is included: \(dL\geq1\) absorbs it in this
explicit bound.
For the actual matrix constructed at this chosen \(p\), division
by the proved positive \(d\) yields
\[
 \boxed{\quad
 \|{\cal L}(H/d-P)\|
 \leq16\sqrt{L/d}+C_\delta^{\rm obs}d\alpha L.
 \quad}
\]
The normalization uses the exact identity
\(H/d-P=(H-dP)/d\), with its nonzero denominator checked.
Compression onto the first block controls the full spectral norm.
The other two diagonal blocks control the signed row sums for the
same unrestricted subsets \(S_1,S_2\), on this same matrix.

Finally set \(E=\{(a,b):H_{ab}=1\}\).
The binary-entry proof identifies its genuine adjacency matrix
with \(B(H)\); the row and column sum identities give exactly
\(d\) distinct incident edges at each vertex.
Every edge lies in the positive support of the original \(P\).
For every finite vertex subset \(S\), the adjacency sum is exactly
\(\deg_E(v,S)\), proved by a finite indicator count.
Thus both observables are actual neighbor counts of one simple
regular bipartite graph, not separate matrix witnesses.
The additional maintained owners are
\texttt{HamiltonianLogMomentChoice},
\texttt{HamiltonianLogarithmicRounding},
\texttt{HamiltonianNormalizedRounding},
\texttt{HamiltonianRoundingReadout} and
\texttt{HamiltonianRoundedEdgeSet}.
The last owner also proves positivity of the dimension-independent
constants
\[
 c_\delta=c_\delta^{\rm repair},\qquad
 C_\delta=\max\{16,C_\delta^{\rm obs}\},
\]
and the complete advertised bound (1), simultaneously with both
neighbor-count conclusions.
In fact the displayed stronger estimate does not require the
additional \(L/d\) term in (1).

\begin{proof}[Proof of Theorem~\ref{thm:rounding}]
Apply the complete deterministic joint construction above, with any
two specified subsets (or with both empty for the spectral assertion).
The full signed prefix identity and every repair are taken relative
to the original kernel and its original gap.
At the even moment \(p=2\lceil2\log(6m)\rceil\), the explicit
scalar estimate and positive-degree normalization give
\[
 \left\|\frac{A_H}{d}-P\right\|
 \leq16\sqrt{L/d}+C_\delta^{\rm obs}d\alpha L.
\]
The binary support and exact degree identities give a simple
\(d\)-regular spanning graph.
Take \(c_\delta=c_\delta^{\rm repair}>0\) and
\(C_\delta=\max\{16,C_\delta^{\rm obs}\}>0\).
Both positive terms are bounded by their corresponding terms
with coefficient \(C_\delta\); adding the nonnegative \(C_\delta L/d\)
gives (1).
All constants depend only on \(\delta\).
This proof is the maintained Lean proof; it requires no probabilistic
concentration input.
\end{proof}

\begin{lemma}[Rounding preserves two specified row observables]
\label{lem:observablerounding}
Under the hypotheses of Theorem~\ref{thm:rounding}, fix any
two vertex subsets \(S_1,S_2\).
The graph \(H\) can be chosen to satisfy its operator bound
and, simultaneously,
\[
 \max_{\substack{v\\S\in\{S_1,S_2\}}}
 \left|\frac{\deg_H(v,S)}d-(P\one_S)(v)\right|
 \leq C_\delta\left(d\alpha L+\sqrt{L/d}+L/d\right).
\]
There is no size restriction on these subsets.
\end{lemma}
\begin{proof}
Use the same joint construction with these \(S_1,S_2\).
The corresponding observable blocks have diagonal entries
\[
 \frac{\deg_H(v,S_j)}d-(P\one_{S_j})(v),\qquad j=1,2.
\]
The modulus of each entry is at most the norm of the complete
joint packet.
The explicit normalized bound and the same
\(C_\delta=\max\{16,C_\delta^{\rm obs}\}\) prove the assertion.
The finite edge-set dictionary and both counts use the very same
supported graph as the spectral estimate.
\end{proof}

\paragraph{Formalization status.}
The repaired stochastic witness now has a Lean proof of the complete
row and column absolute-mass estimate, uniformly over every subset
observable.  With \(L=\log(2N)\), its explicit constant is
\[
 C_\delta^{\mathrm{obs}}
 =3C_\delta^{\mathrm{scale}}+\frac1{\log4}.
\]
Indeed, write \(T=P-Q\), \(r=C_\delta^{\mathrm{scale}}\eta L\)
and \(W=\diag(1+z_L)T\diag(1+z_R)\), using the same scaling
witness as in Corollary~\ref{cor:forbiddenrepair}.
The smallness condition gives \(0\leq r\leq1\), and the exact
entrywise difference is
\[
 W_{ij}-P_{ij}
 =(z_L(i)+z_R(j)+z_L(i)z_R(j))T_{ij}-Q_{ij}.
\]
Consequently every row and every column has absolute mass at most
\[
 \eta+2r+r^2\leq\eta+3r
 \leq C_\delta^{\mathrm{obs}}\eta L .
\]
This retains the deleted mass and the quadratic cross term.
Summing over any subset gives the stated predictable observable drift
with no restriction on its size.  The maintained owners are
\texttt{HamiltonianScaledCrossMass} and
\texttt{HamiltonianForbiddenObservableDrift}.
The logarithmic moment choice, full normalization, simple edge-set
construction, spectral readout and both actual neighbor-count readouts
are now all formalized in the five owners listed above.
Theorem~\ref{thm:rounding} and Lemma~\ref{lem:observablerounding}
have complete maintained Lean proofs.
Other headline claims of the paper remain under formalization.


\begin{proof}[Proof of Corollary~\ref{cor:weightedhamilton}]
First suppose \(\alpha\leq N^{-a}\) and take
\(d=\lfloor N^{a/2}\rfloor\).
Then \(d\alpha L^2\to0\), and every term in (1) tends to zero.
Since both \(A_H/d\) and \(P\) have the exact eigenvectors
\(\one,s\) at eigenvalues \(1,-1\), Weyl's inequality on
their common complement gives
\[
 \lambda_2(H)\leq(1-\delta/2)d
\]
for sufficiently large \(N\).
The polynomial degree exceeds
\((6\delta^{-1}\log N)^K\).
Equation (2) makes \(H\), and hence the support of \(P\),
Hamiltonian.

For (3), choose a fixed \(A_\delta>0\) sufficiently large
and set \(d=\lceil A_\delta L^{K+1}\rceil\).
For all sufficiently large \(N\), this is at most \(N/2\)
and exceeds \((6\delta^{-1}\log N)^K\).
Under (3),
\[
 d\alpha L^2\leq 2A_\delta b_\delta,\qquad
 d\alpha L\leq 2A_\delta b_\delta/L.
\]
Choose \(b_\delta\) small enough for the first quantity to
be at most the constant \(c_\delta\) in Theorem~\ref{thm:rounding}.
The second tends to zero, as do the two sampling errors.
The same spectral-gap and Hamiltonicity consumer applies.
The concentration obstruction is its contrapositive.
\end{proof}

\begin{remark}[What is independent and what is quantitative]
The finite rounding bound (1) is proved without a Hamiltonicity
input. It preserves the spectrum, exact degrees, and simplicity
simultaneously. The Hamiltonicity conclusion then uses the
independent recent theorem (2).
The constants \(c_\delta,C_\delta\) are controlled by the
scaling contraction and the deterministic matrix-moment budget; their numerical
values have not been optimized.
The starting order in the Hamiltonicity corollary is
asymptotic and has not been computed.
\end{remark}

\section{The integer prime-sum operator and its application}
\label{sec:primesum}
Let \(N\) be even. For \(1\leq i,j\leq N\), first put
\[
 A_N(i,j)=\frac{\Lambda(i+j)}{N}.
\]
The normalized vector inner product is
\(\langle f,g\rangle_N=N^{-1}\sum_i\overline{f_i}g_i\);
its induced operator norm agrees with the Euclidean one.
The adjacency support used at the end will exclude both
proper prime powers and diagonal entries.

\begin{lemma}[Signed integer-input transfer]
\label{lem:integertransfer}
For every fixed \(B>0\), there is a fixed \(C>0\) such that,
with \(Q=(\log N)^C\),
\[
 \left\|A_N-
    \left(\frac1N\sum_{q\leq Q}
                \frac{\mu(q)}{\phiG(q)}c_q(i+j)\right)_{ij}
       \right\|\ll_B(\log N)^{-B}.                          \tag{16}
\]
Here \(c_q(t)=\sum_{a\bmod q,(a,q)=1}\e(at/q)\)
is the Ramanujan sum. The complete signed main term is retained.
\end{lemma}
\begin{proof}
We record the argument orientation explicitly.
For distinct symbolic vectors \(f,g\), set
\(F_f(\theta)=\sum_{i=1}^N f_i\e(i\theta)\) and
\(S_M(\theta)=\sum_{r\leq M}\Lambda(r)\e(r\theta)\),
where \(M=4N+4\). Exact Fourier inversion gives
\[
 \langle f,A_Ng\rangle_N
 =N^{-2}\int_0^1 S_M(\theta)
             \overline{F_f(\theta)}F_g(-\theta)\,d\theta .
 \tag{17}
\]
Thus the left matrix factor is the conjugate of the first
argument. Parseval and Cauchy--Schwarz give
\[
 \int_0^1|F_f(\theta)F_g(-\theta)|\,d\theta
                  \leq N\|f\|_N\|g\|_N.                    \tag{18}
\]

Use major arcs \(|\theta-a/q|\leq T/N\), \(q\leq Q\),
where \(T=(\log N)^D\) and \(D\) is a sufficiently large
fixed exponent. These arcs are disjoint for large \(N\).
Take the integer \(H=\lceil N/T\rceil\).
Dirichlet approximation supplies coprime \(a,q\) with
\(1\leq q\leq H\) and
\[
 |\theta-a/q|\leq1/(qH)\leq T/(qN).
\]
If \(q\leq Q\), this lies in the stated major arc.
On their complement \(Q<q\leq H\leq2N/T\), and
Vaughan's standard estimate~\cite{Vaughan} implies
\[
 |S_M(\theta)|
 \ll N(\log N)^4
       \left(Q^{-1/2}+T^{-1/2}+N^{-1/5}\right).
\]
For completeness, the estimate being used is
\[
 |S_M(\theta)|\ll
 (Mq^{-1/2}+M^{4/5}+M^{1/2}q^{1/2})(\log M)^4
\]
when \(|\theta-a/q|\leq q^{-2}\).
The approximation also satisfies \(|\theta-a/q|\leq q^{-2}\),
because \(q\leq H\). On the minor arcs its denominator
satisfies \(Q<q\leq2N/T\).
Equation (18) therefore turns this into an operator error
smaller than \((\log N)^{-B}\) if \(C,D\) are large enough.

On each major arc, Siegel--Walfisz and partial summation give,
with arbitrary prescribed logarithmic accuracy,
\[
 S_M(a/q+\beta)
  =\frac{\mu(q)}{\phiG(q)}V_M(\beta)
          +O(N(\log N)^{-A}),\qquad
 V_M(\beta)=\int_0^M\e(t\beta)\,dt .
\]
The ranges \(q\leq Q\) and \(|\beta|\leq T/N\) are fixed
logarithmic ranges. Choose the underlying prime estimate
with accuracy sufficient to absorb the partial-summation
factor \(1+M|\beta|\).
The uniform error, inserted in (17)--(18), is again
\(O((\log N)^{-B})\).

It remains to justify the entire retained main kernel,
including the arc tails.
For \(s=i+j\in[2,2N]\), Fourier inversion of the interval
\([0,M]\), followed by one integration by parts in each
oscillatory tail, gives
\[
 \left|\int_{-T/N}^{T/N}
            V_M(\beta)\e(-s\beta)\,d\beta-1\right|
 \leq C\left(\frac{N}{T(i+j)}+\frac1T\right).
\]
The normalized matrix \(N^{-1}\) times this error has norm
\(O(1/T)\): the Hilbert matrix \((1/(i+j))\) has bounded
norm, for example by the Schur test with weights \(i^{-1/2}\).
The phase at \(a/q\) is a diagonal unitary on each side,
so does not change this bound.
Summing the explicit remainder over reduced \(a\bmod q\)
and \(q\leq Q\) costs at most \(O(Q/T)\).
Choose \(D>C+B+10\).
The retained phases sum to
\(\mu(q)c_q(-(i+j))/\phiG(q)\), which equals the displayed
main term since \(c_q(-s)=c_q(s)\).
This proves (16) without applying an unsigned estimate to
that main term.
\end{proof}

\paragraph{Formalized whole-prefix minor-grid input.}
The maintained owners \texttt{HamiltonianPrimeMinorArc} and
\texttt{HamiltonianPrimePrefixMinor} consume the existing actual
coefficient-mean Vaughan theorem and its proved physical cutoffs.
Precisely, let \(M\geq2^{18}\), \(1\leq P\leq M^{1/40}\), and
let the positive auxiliary period \(R\) be divisible by \(M\)
and by every integer \(1\leq q\leq P\). Put
\[
 W=(R/M)P,\qquad
 \mathfrak M_R(P,W)=
 \left\{\frac{Ra}{q}+j\pmod R:
  1\leq q\leq P,\ 0\leq a<q,\ (a,q)=1,\ |j|\leq W\right\}.
\]
For one absolute \(C>0\), independent of every displayed
parameter, the complete actual prefix satisfies
\[
 \left|\sum_{r=1}^M\Lambda(r)\e(kr/R)\right|^2
 \leq C\left(\frac{M^2}{P}+M^{8/5}\right)
                     (1+\log(2M))^{10},
 \qquad k\in\mathbb Z,\quad k\bmod R\notin\mathfrak M_R(P,W).
\]
The initial Vaughan cutoff is retained and bounded explicitly;
every prime power remains in the prefix. There is no divisor
power loss in this estimate. All five new declarations pass
the focused foundational-axiom check.
The full sampled-operator transfer, retained major-arc main term
and spectral comparison remain formalization obligations.

\begin{lemma}[The nontrivial ratio is one half]
\label{lem:integerspectrum}
Apart from at most one eigenvalue near \(1\) and at most one
near \(-1\), every eigenvalue of \(A_N\) has modulus at most
\(1/2+o(1)\). More precisely, in decreasing algebraic order,
\[
 \lambda_2(A_N)\leq1/2+o(1),\qquad
 \lambda_{N-1}(A_N)\geq-1/2-o(1).
 \tag{19}
\]
\end{lemma}
\begin{proof}
Use the orthonormal model basis indexed by reduced fractions
\(a/q\bmod1\), with squarefree \(q\leq Q\), including
\(0/1\) and \(1/2\).
On this basis define the self-adjoint operator
\[
 C_Q e_{a/q}=\frac{\mu(q)}{\phiG(q)}e_{-a/q}.
\]
Its eigenvalues at \(0/1,1/2\) are \(1,-1\).
All other blocks have norm \(1/\phiG(q)\leq1/2\).
It is the finite cutoff of a compact direct sum of these
orthogonal blocks.

Let \(E_Qe_{a/q}(j)=\e(aj/q)\), using the normalized sample
inner product. Its Gram matrix has diagonal 1.
For two different fractions, geometric summation bounds
the off-diagonal entry by \(C Q^2/N\).
There are \(O(Q^2)\) basis vectors, so
\[
 \|E_Q^*E_Q-I\|\ll Q^4/N=o(1).
\]
Direct multiplication, with the first argument on the left
as in (17), gives
\[
 (E_QC_QE_Q^*)(i,j)
       =\frac1N\sum_{q\leq Q}
                    \frac{\mu(q)}{\phiG(q)}c_q(i+j).
\]
The polar map \(U_Q=E_Q(E_Q^*E_Q)^{-1/2}\) is an isometry,
and the evaluated operator differs from
\(U_QC_QU_Q^*\) by \(O(Q^4/N)\) in operator norm.
Append zero eigenvalues to this finite model.
Equation (16) and Weyl's inequality give (19).
There is one exceptional eigenvalue at each end of the
spectrum, not two positive exceptions.
\end{proof}

\begin{lemma}[Uniform degrees and a true prime support]
\label{lem:primeweights}
Define the true prime-edge weight matrix
\[
 Q_N(i,j)=\frac{\log(i+j)}{N}
       \one_{\,i\ne j,\ i+j\ {\rm prime}}.
\]
It is symmetric and supported between the odd and even
integers, which form equal parts.
For every fixed \(B>0\),
\[
 \|Q_N\one-\one\|_\infty\ll_B(\log N)^{-B},\qquad
 \|Q_N-A_N\|\ll N^{-1/2}(\log N)^2,\qquad
 \max_{ij}Q_N(i,j)\leq\frac{\log(2N)}N.                      \tag{20}
\]
If \(D_N=\diag(Q_N\one)\), its normalized operator
\(D_N^{-1/2}Q_ND_N^{-1/2}\) has simple eigenvalues \(1,-1\)
and all others of modulus at most \(1/2+o(1)\).
\end{lemma}
\begin{proof}
The exact raw row sum is
\[
 (A_N\one)_i
       =\frac{\psi(i+N)-\psi(i)}{N}.
\]
The classical strong prime number theorem,
\(\psi(x)=x+O_B(x/(\log(x+2))^B)\), gives \(1+
O_B((\log N)^{-B})\) uniformly for \(1\leq i\leq N\).
Constants for bounded \(x\) can be absorbed; equivalently
split the range before applying the estimate.

For a fixed input, every proper-prime-power output
determines at most one other vertex.
There are \(O(N^{1/2}\log N)\) such outputs up to \(2N\),
each of weight at most \(\log(2N)\).
The normalized removed row sum is therefore
\(O(N^{-1/2}(\log N)^2)\).
The diagonal costs at most \(\log(2N)/N\) per row.
Symmetry and the row-sum norm bound prove the second
estimate in (20); the same errors preserve the first.
The entry bound is immediate.
Every off-diagonal prime output is odd, because the sole
representation of prime 2 by positive equal-bounded inputs
is the removed loop at vertex 1.
Thus the remaining support is exactly bipartite.

Its degree matrix is \(I+o(1)\), so normalization changes
the operator by \(o(1)\).
The normalized matrix has eigenvalues \(1,-1\), with vectors
\(D_N^{1/2}\one,D_N^{1/2}s\).
Equation (19) permits no further eigenvalue of modulus
larger than \(1/2+o(1)\).
It also forces these two eigenvalues to be simple;
multiple components would supply extra eigenvalues at 1.
This proves the final assertion.
\end{proof}

\begin{proof}[Proof of Theorem~\ref{thm:filz}]
Apply Lemma~\ref{lem:scaling} to \(Q_N\), with any fixed gap
smaller than \(1/2\), using (20).
For instance take the row accuracy \(B=10\);
then \(\varepsilon L^2\to0\).
The result is a symmetric stochastic bipartite \(P_N\)
with exactly the true prime-edge support and
\[
 \|P_N-Q_N\|=o(1),\qquad
 \max_{ij}(P_N)_{ij}\ll\frac{\log N}{N}.
\]
Its nontrivial spectral ratio is \(1/2+o(1)\).
Fix a positive gap such as \(1/3\) for all sufficiently
large \(N\), and take \(d_N=\lfloor N^{1/4}\rfloor\).
The finite rounding theorem applies, since
\[
 d_N\alpha_N L^2\ll N^{-3/4}(\log N)^3\longrightarrow0.
\]
Its complete operator error is
\[
 O\left(N^{-3/4}(\log N)^2+
             N^{-1/8}\sqrt{\log N}+N^{-1/4}\log N\right)=o(1).
\]
It gives a simple \(d_N\)-regular spanning subgraph of
\(G_N\), whose second eigenvalue is at most
\((1/2+o(1))d_N\).
This is (4).
For sufficiently large \(N\) its gap is at least \(1/3\).
The polynomial degree exceeds \((9\log N)^K\).
The established bipartite theorem (2) therefore supplies
a Hamilton cycle in this subgraph, and hence in \(G_N\).
All proper-prime-power and loop contributions were removed
before the support was used.
\end{proof}

\begin{remark}[Quantifiers and the finite remainder]
Every sufficiently large even \(N\) is covered, rather than
a selected subsequence of prime gaps.
The starting threshold depends on classical prime estimates
and on the large-order threshold in (2), and has not been
made numerical. This draft therefore leaves the finite
remainder of the all-orders Filz conjecture unresolved.
No finite prime experiment is used to fill that remainder.
\end{remark}

\section{Near-optimal factors, resilience and many prime Hamilton cycles}
\label{sec:regularfactors}
The weighted rounding theorem permits very nonuniform edge weights.
In the integer prime-sum application there is additional uniformity:
ordinary, unweighted degrees are asymptotically equal.
This allows an integral regular-factor argument, proved here, which
strengthens the arithmetic conclusion substantially.
It is an alternative application mechanism; adaptive rounding is
not claimed to be necessary for this particular uniform graph.

\begin{lemma}[Regular factors from near regularity and cuts]
\label{lem:regularfactor}
Let \(G\) be a simple balanced bipartite graph on \(N=2m\)
vertices. Suppose, for real \(D>0\), \(0<g\leq1\) and
\(0\leq\varepsilon<g/2\), that
\[
 (1-\varepsilon)D\leq\deg_G(v)\leq(1+\varepsilon)D
 \quad\hbox{for every }v,
\]
and
\[
 e_G(W,W^c)\geq gD\min\{|W|,N-|W|\}
 \quad\hbox{for every }W\subseteq V(G).
 \tag{21}
\]
Then \(G\) contains an \(r\)-regular spanning subgraph whenever
\(r\) is a nonnegative integer with
\[
 r\leq(1-2\varepsilon/g)D.                                  \tag{22}
\]
In particular one may take the floor of the right side.
Every such factor \(H\) satisfies the finite operator bound
\[
 \|A_G-A_H\|\leq(1+\varepsilon)D-r.                          \tag{23}
\]
\end{lemma}
\begin{proof}
Write the parts as \(A,B\).
The finite matching criterion for an \(r\)-factor is
\[
 e_G(S,B\setminus T)\geq r(|S|-|T|)
 \quad(S\subseteq A,\ T\subseteq B).                         \tag{24}
\]
Here is a capacity-one construction proving sufficiency, including
\(r=0\). Form a bipartite matching gadget with left vertices
\(A\times\{1,\ldots,r\}\) and one \(L_e\) for each edge, and right
vertices \(B\times\{1,\ldots,r\}\) and one \(R_e\) for each edge.
Join an \(A\)-slot to \(R_e\) when \(e\) starts at that vertex,
join \(L_e\) to every slot of its \(B\)-endpoint, and join
\(L_e\) to \(R_e\).

For an arbitrary left set, let \(K\) be its slots, \(F\) its edge
vertices, \(S\) the vertices touched by \(K\), and \(T\) the
\(B\)-endpoints of \(F\). Its neighborhood contains all \(r|T|\)
slots over \(T\), all \(R_e\) with \(e\in F\), and all \(R_e\)
with \(e\in E(S,B\setminus T)\). These last two edge sets are
disjoint. Thus (24) gives
\[
 |N(K\cup F)|\geq r|T|+|F|+e_G(S,B\setminus T)
               \geq r|S|+|F|\geq |K|+|F|.
\]
Finite Hall matching applies. Both gadget parts have size
\(rm+|E(G)|\), so the matching is perfect. An edge \(e\) is
either matched internally \(L_eR_e\), or its two vertices are
matched to an \(A\)-slot and a \(B\)-slot. Select the latter
edges. Each slot is matched once and each edge has only one
private pair, so the selected graph is simple and has degree
exactly \(r\) at every vertex.

Only \(x=|S|-|T|>0\) matters.
By exchanging the two sides and replacing \((S,T)\) with
\((B\setminus T,A\setminus S)\), if necessary, we may assume
\[
 u=|S|+|T|\leq m.
\]
Both \(x\) and the edge count in (24) are preserved by this
replacement. Put
\[
 a=e_G(S,B\setminus T),\qquad b=e_G(A\setminus S,T).
\]
The full boundary of \(S\cup T\) is \(a+b\), so (21) and
the degree bounds give
\[
 a+b\geq gDu,\qquad a-b\geq Dx-\varepsilon Du.
\]
Since \(b\geq0\), the two simultaneous lower bounds are
\[
 a\geq D(x-\varepsilon u),\qquad
 a\geq \frac D2\bigl(x+(g-\varepsilon)u\bigr).                \tag{25}
\]
If \(x\leq(g-\varepsilon)u\), the second gives \(a\geq Dx\).
Otherwise the first gives
\[
 a\geq Dx\left(1-\frac{\varepsilon}{g-\varepsilon}\right)
       \geq Dx(1-2\varepsilon/g).
\]
Thus (24) follows from (22).
Finally \(A_G-A_H\) is a symmetric nonnegative matrix with
maximum row sum at most the right side of (23).
Its operator norm obeys the same bound.
\end{proof}

The complete finite lemma is formalized, with the actual edge set
and adjacency matrix. \texttt{HamiltonianFiniteHall} provides
the finite matching proof without rebuilding a dependency cache.
\texttt{HamiltonianRegularFactorHall} retains the complete slot
and private-edge gadget; \texttt{HamiltonianRegularFactorMatching}
extracts the simple factor and proves both exact degree equalities.
\texttt{HamiltonianFactorIncidence} proves the signed crossing
balance and its complementary degree bound.
\texttt{HamiltonianRegularFactorCriterion} combines these with
the same full cut and proves (22).
\texttt{HamiltonianRegularFactorOperator} proves the exact
adjacency row counts and (23), and assembles both conclusions
for the same factor. All thirty-four declarations in these six
owners pass the focused foundational-axiom audit. The actual
near-optimal prime-factor consumer, including the single family
and its native adjacency eigenvalues, is formalized below.

Let \(B_N\) denote the ordinary adjacency matrix of \(G_N\)
and put \(D_N^*=N/\log N\).
This scalar is distinct from the degree matrix \(D_N\)
in Lemma~\ref{lem:primeweights}.

\begin{lemma}[Unweighted spectral transfer]
\label{lem:unweightedtransfer}
As even \(N\) tends to infinity,
\[
 \max_i\left|\frac{\deg_{G_N}(i)}{D_N^*}-1\right|
       =O\left(\frac{\log\log N}{\log N}\right),\qquad
 \|B_N/D_N^*-Q_N\|
       =O\left(\frac{\log\log N}{\log N}\right).              \tag{26}
\]
There is a stochastic kernel \(P_N\) on the same support with
\[
 \|B_N/D_N^*-P_N\|=o(1),\qquad
 \|P_N|_{\{\one,s\}^{\perp}}\|\leq1/2+o(1).
 \tag{27}
\]
Consequently, for all sufficiently large \(N\),
\[
 e_{G_N}(W,W^c)\geq\frac15D_N^*
                  \min\{|W|,N-|W|\}                        \tag{28}
\]
for every vertex set \(W\).
\end{lemma}
\begin{proof}
The difference in (26) has entries
\[
 \frac{\log N-\log(i+j)}N
                    \one_{\,i\ne j,\ i+j\ {\rm prime}}.
\]
When \(i+j\geq N/(\log N)^4\), the absolute numerator is at
most \(4\log\log N+\log2\), and each row has at most
\(\pi(2N)=O(N/\log N)\) nonzero entries.
The remaining outputs number at most \(N/(\log N)^4\)
per row, each with absolute numerator at most \(\log(2N)\).
Thus the maximum absolute row sum of the difference is
\(O(\log\log N/\log N)\).
Symmetry proves the operator bound, while the row-sum version
and (20) prove the degree assertion.
The exact scaling already used in the proof of
Theorem~\ref{thm:filz} gives (27).

For the boundary estimate put \(w=|W|\) and
\(x=\one_W-(w/N)\one\).
The identity for the combinatorial Laplacian gives
\[
 e_{G_N}(W,W^c)
 =x^\top(\diag(B_N\one)-B_N)x.
\]
On \(\one^\perp\), the algebraically largest eigenvalue of
\(P_N\) is at most \(1/2+o(1)\); its parity eigenvalue is \(-1\).
Equations (26)--(27), with the degree lower bound, therefore give
\[
 e_{G_N}(W,W^c)
 \geq(1/2-o(1))D_N^*\|x\|_2^2
 =(1/2-o(1))D_N^*\frac{w(N-w)}N.
\]
Since \(w(N-w)/N\geq\frac12\min(w,N-w)\), this implies (28).
\end{proof}

\paragraph{Formalized complete unweighted arithmetic transfer.}
The actual constant-edge matrix
\[
 U_n(i,j)=\frac{\log n}{n}
      \one_{i\ne j,\ i+j\ {\rm prime}}
\]
is now defined in \texttt{HamiltonianUnweightedPrimeTransfer}.
The original loop and prime masks are preserved exactly.
For every positive integer \(n\) with \(\log n\geq64\), its
comparison with the true logarithmic prime matrix satisfies
\[
 \max_i\sum_j|U_n(i,j)-Q_n(i,j)|\leq E_n,\qquad
 \max_j\sum_i|U_n(i,j)-Q_n(i,j)|\leq E_n,\qquad
 \|U_n-Q_n\|\leq E_n,
\]
where the explicit error is
\[
 E_n=\frac{1+4\log4\,(2\log\log n+\log2)}{\log n}.
\]
\texttt{HamiltonianLogWeightComparison} proves the two
pointwise regimes and the complete finite packing estimate.
\texttt{HamiltonianPrimeLogTransfer} uses
\(K=\lfloor n/(\log n)^2\rfloor\), charges every small
output to its integer count and every large output to its
actual prime logarithmic weight, and consumes
\(\theta(2n)\leq2n\log4\).
No prime-count asymptotic or independent prime-pair estimate
is assumed in this transfer.
The injective physical row intervals, all endpoints and both
matrix directions are included before the separate operator
error is bounded.

\texttt{HamiltonianPrimeLogTransferDecay} proves the uniform
decay on an entire terminal interval, including
\[
 E_n\leq\frac{1+16\log4+4\log4\log2}{\sqrt{\log n}}.
\]
It also proves, for every natural \(b\), one \(D_b>0\)
and one starting order such that every later row satisfies
\[
 \left|\sum_jU_n(i,j)-1\right|
                   \leq E_n+D_b/(\log n)^b.
\]
Thus the logarithmic rate in (26) follows from an explicit
uniform inequality. The actual constant-edge spectral consumer
is formalized without an auxiliary grid:
\[
 \forall\varepsilon>0\ \exists n_0 \forall n\geq n_0:\qquad
 \|U_n-P_{0,n}\|\leq\tfrac12+\varepsilon .
\]
All twenty-four new declarations in these four maintained owners
pass the focused foundational-axiom audit.
The balanced odd/even cross-block identification and the actual
support-preserving stochastic repair are formalized, as detailed
in Section~\ref{sec:relativeprimecircles}. The same witness gives
the exact entropy maximum and the actual unoriented prime-circle
count, including the factorial normalization and the quantitative
Stirling conversion. Its two relative errors are respectively
at most \(m^{-3/5}\) and \(2m^{-3/5}\), \(m=N/2\).
\paragraph{Formalized sharp repair and the full prime-edge cut bound.}
\texttt{HamiltonianPrimeSharpScaling} retains the same positive
counting witness while making its perturbation arbitrarily small.
\texttt{HamiltonianPrimePhysicalScaling} transfers it back through
the exact odd/even bijection. For every \(\eta>0\), for every
sufficiently large \(m\), the actual physical matrix \(P_{2m}\)
is symmetric, nonnegative and exactly stochastic, with
\[
 P_{2m}(i,j)=0\ \Longleftrightarrow\ U_{2m}(i,j)=0,\qquad
 \|P_{2m}-U_{2m}\|\leq\eta,\qquad
 \|P_{2m}-P_{0,2m}\|\leq\tfrac12+\eta .
\]
\texttt{HamiltonianPrimeOrthogonalGap} proves the exact
constant/parity multiplication formula and the operator bound
on the subspace with both corresponding sums zero.
This completes the literal sharp-gap consumer in (27).

For (28), \texttt{HamiltonianSpectralCuts} and
\texttt{HamiltonianCenteredCuts} prove the full signed
Dirichlet identity and exact centered indicator norm.
For every centered cut vector \(x\),
\texttt{HamiltonianPrimeCuts} retains
\[
 x^\top P_{0,N}x=-\frac1N
       \left(\sum_i(-1)^i x_i\right)^2\leq0.
\]
The degree lower bound \(19/20\) and operator error at most
\(11/20\) give the complete finite bound
\[
 e_{U_N}(W,W^c)\geq\frac25\,\frac{|W|(N-|W|)}N
                       \geq\frac15\min\{|W|,N-|W|\}.
\]
The same owner defines the actual finite set of distinct
prime-sum edges crossing the cut, proves its exact weighted
cardinality identity, and concludes
\[
 e_{G_N}(W,W^c)\geq\frac{N}{5\log N}
                           \min\{|W|,N-|W|\}
\]
on an entire terminal interval, uniformly over every \(W\).
No cut or overlap is removed. All thirty-nine new declarations
in these six owners pass focused foundational-axiom checks;
the two audit batches are disjoint. Equations (26)--(28)
now have complete maintained arithmetic consumers.

\begin{theorem}[Near-optimal regular prime-sum factors]
\label{thm:nearlyoptimalfactor}
For every sufficiently large even \(N\), \(G_N\) has a simple
\(r_N\)-regular spanning subgraph \(J_N\) such that
\[
 r_N=(1-o(1))\frac N{\log N},\qquad
 \|A_{J_N}/r_N-P_N\|=o(1).
 \tag{29}
\]
In particular every eigenvalue other than its exact
\(r_N,-r_N\) has modulus at most \((1/2+o(1))r_N\).
The factor degree is asymptotically best possible.
\end{theorem}
\begin{proof}
Use Lemma~\ref{lem:regularfactor} with \(g=1/5\),
\(D=D_N^*\), and the vanishing degree defect from (26).
It supplies
\[
 r_N=\left\lfloor(1-10\varepsilon_N)D_N^*\right\rfloor
                      =(1-o(1))D_N^* .
\]
The discarded adjacency has norm \(o(D_N^*)\) by (23).
Together with (27) and \(r_N/D_N^*\to1\), this proves (29).
The vectors \(\one,s\) remain exact, so the asserted
nontrivial spectral bound follows on their common complement.
Any regular factor has degree at most the minimum degree of
\(G_N\), which is \((1+o(1))D_N^*\).
\end{proof}

The finite prime-factor consumer is now formalized.
\texttt{HamiltonianFactorCutGeometry} identifies both crossing
counts with the full cut boundary and transports it losslessly
through the odd/even vertex bijection.
\texttt{HamiltonianPrimeFactorInputs} defines the actual prime
edge set, proves both ordinary degree estimates with
\(D=2m/\log(2m)\), and transfers the cut coefficient \(1/5\)
to that same edge set.
\texttt{HamiltonianPrimeRegularFactor} proves that, for every
fixed \(0<\varepsilon<1/10\), there is \(M_0\) such that
every \(m\geq M_0\) and every integer
\(0\leq r\leq(1-10\varepsilon)D\) have a simple factor \(H\)
on the original integer vertices with
\[
 \deg_H(v)=r,\qquad E(H)\subseteq E(G_{2m}),\qquad
 \|A_{G_{2m}}-A_H\|\leq(1+\varepsilon)D-r.
\]
Its declarations check the original distinct-integer prime
support and exact physical degree and norm transfers.
All twenty-three declarations in these three owners pass
focused foundational-axiom checks.

The single shrinking-defect factor family and its complete
operator normalization are now formalized as well.
The simultaneous degree estimate supplies one fixed \(B>0\).
Write \(L=\log(2m)\) and \(D=2m/L\). For every sufficiently
large \(m\), the chosen factor has the fixed degree
\[
 r_m=\left\lfloor
       \left(1-\frac{10B}{\sqrt L}\right)D
     \right\rfloor,\qquad D/2\leq r_m\leq D.
\]
\texttt{HamiltonianFactorFloor} retains the full one-unit
floor error and proves \(D\geq L/4\).
\texttt{HamiltonianPrimeFactorFamily} constructs one family,
with both original degree estimates retained, and proves
\[
 \left|\frac{r_m}{D}-1\right|
       \leq\frac{10B}{\sqrt L}+\frac4L.
\]
Its terminal threshold is fixed; no moving defect is
substituted into a fixed-defect eventual theorem.
\texttt{HamiltonianFactorNormalization} uses the exact identity
\[
 \frac{A_{J_m}}{r_m}-U_{2m}
 =\frac{A_{J_m}-D U_{2m}}{r_m}
       +\frac{D-r_m}{r_m}U_{2m}.
\]
\texttt{HamiltonianPrimeFactorNormalization} consumes both
terms for that same chosen family and proves
\[
 \left\|\frac{A_{J_m}}{r_m}-U_{2m}\right\|
                 \leq\frac{62B}{\sqrt L}+\frac{24}L.
\]
The shared \texttt{HamiltonianSupportMixture} theorem proves
that a positive mixture of a stochastic factor with a
full-support stochastic reference retains every original
edge, with perturbation at most twice the mixing weight.
Twenty-one core declarations in these five owners pass
focused foundational-axiom checks.

Theorem~\ref{thm:nearlyoptimalfactor} now has its complete
maintained consumer. \texttt{HamiltonianPrimeFactorReference}
proves the normalized factor's actual stochastic law and
constructs one full-support reference family
\[
 P_m=(1-\delta_m)\frac{A_{J_m}}{r_m}+\delta_m R_m,
 \qquad \delta_m=\frac1{m+1},
\]
using one fixed-terminal full-support stochastic family \(R_m\).
It retains exact symmetry, nonnegativity, every row and column
sum, and every original prime edge, with
\[
 \left\|\frac{A_{J_m}}{r_m}-P_m\right\|\leq\frac2{m+1},
 \qquad
 \|P_m-U_{2m}\|\leq
       \frac{62B}{\sqrt L}+\frac{24}L+\frac2{m+1}.
\]
\texttt{HamiltonianPrimeFactorLimits} proves the complete
terminal epsilon bounds; the same chosen family works for
every precision. \texttt{HamiltonianBipartiteEigenvalues}
derives both vanishing block sums from the actual eigenvector
equations whenever the normalized eigenvalue differs from
\(1,-1\). Thus the constant/parity reference acts as zero,
without an added orthogonality hypothesis.
\texttt{HamiltonianPrimeFactorEigenvalues} transfers this
argument through the exact vertex bijection and applies it
to the native real Hermitian eigenvalues of \(A_{J_m}\).

Finally, \texttt{HamiltonianPrimeFactorHeadline} proves,
for every \(\eta>0\), one terminal threshold for which the
same \(J_m,P_m\) satisfy
\[
 \left|\frac{r_m}{D}-1\right|\leq\eta,\qquad
 \left\|\frac{A_{J_m}}{r_m}-P_m\right\|\leq\eta,\qquad
 \left\|\frac{A_{J_m}}{r_m}-P_{0,2m}\right\|
       \leq\frac12+\eta .
\]
Every native adjacency eigenvalue unequal to \(r_m,-r_m\)
then has modulus at most \((1/2+\eta)r_m\).
The exact constant and parity eigenvector identities are
also proved. Uniformly over every other regular factor
\(K\subseteq G_{2m}\) of degree \(s\), the same terminal
interval gives \(s/D\leq1+\eta\), proving asymptotic degree
optimality. All twenty-nine core declarations in these five
new owners pass focused foundational-axiom checks.
The resilience theorem below now has its complete signed-count
consumer. The cycle-packing theorem retains its separate
formalization obligation.

\begin{theorem}[A positive proportion of the prime edges in Hamilton cycles]
\label{thm:manycycles}
For every fixed \(\kappa<1/8\) with \(\kappa>0\), and every
sufficiently large even \(N\), \(G_N\) contains at least
\[
 \left\lfloor\kappa\,\frac N{\log N}\right\rfloor
 \tag{30}
\]
pairwise edge-disjoint Hamilton cycles.
Consequently the supremal number of such cycles, divided by
\(N/\log N\), has lower limit at least \(1/8\).
\end{theorem}
\begin{proof}
Start with \(J_N\) from (29), of degree \(r=r_N\).
After deleting \(j\) edge-disjoint Hamilton cycles, the remaining
graph \(J^{(j)}\) is simple, bipartite and \((r-2j)\)-regular.
The deleted adjacency has norm at most \(2j\). Weyl's inequality,
including this full signed subtraction, gives
\[
 \lambda_2(J^{(j)})\leq(1/2+o(1))r+2j,
\]
and hence
\[
 (r-2j)-\lambda_2(J^{(j)})
                   \geq(1/2-o(1))r-4j.                    \tag{31}
\]
Fix any \(\kappa'\) strictly between \(\kappa\) and \(1/8\).
For all \(j\leq\kappa'r\), (31) leaves a fixed positive
spectral gap as a proportion of the remaining degree.
That degree is of order \(N/\log N\), so exceeds the
polylogarithmic requirement in (2), uniformly in these \(j\),
for all sufficiently large \(N\).
Equation (2) supplies another Hamilton cycle at each step.
Since \(r/(N/\log N)\to1\), this produces (30).
No independence of the successive cycles is used.
\end{proof}

\paragraph{Formalized retained-mode input for cycle packing.}
The full uniform packing conclusion above remains a formalization
obligation. Its arithmetic finite-mode input is now proved.
Let \(Q_n\) be the proved moving prime-operator cutoff, fix a
positive integer \(K\), and put \(H=27K^4\). Index the retained
modes by
\[
 {\cal U}_H=\{(q,a):1\leq q\leq H,\quad
                          a\in(\mathbb Z/q\mathbb Z)^\times\}.
\]
The exact fixed-index model \(F_{n,R,K}\) retains the coefficient
\(\mu(q)/\varphi(q)\) when \(q\leq Q_n\), and zero otherwise,
together with the original sampled character phase. The
auxiliary grid has \(R>2n\), \(2n\mid R\), and \(q\mid R\)
for every \(1\leq q\leq Q_n\).

\texttt{HamiltonianRationalTailNorm} consumes the already proved
uniform inequality \(2q^3\leq27\varphi(q)^4\), deriving
\[
 q>27K^4\quad\Longrightarrow\quad
 |\mu(q)/\varphi(q)|\leq1/K.
\]
It proves the full two-frame remainder bound
\[
 \|\mathsf R_{n,R,Q}-\mathsf R_{n,R,Q}^{\leq H}\|
 \leq K^{-1}(1+4Q^4/n).
\]
\texttt{HamiltonianRationalCutoffTransport} exchanges the moving
and fixed denominator indices exactly, retaining every signed
coefficient, phase and inactive denominator.
\texttt{HamiltonianPrimeRationalApproximation} exposes the
stronger part of the complete major/minor calculation before
its two-mode relaxation: uniformly on a terminal interval,
\[
 \|\mathsf M_n-\mathsf R_{n,R,Q_n}\|+2Q_n^4/n
                                                    \leq\varepsilon,
\]
where \(\mathsf M_n\) is the actual Mangoldt sum matrix.
Thus \texttt{HamiltonianPrimeFiniteModeApproximation} proves
\[
 \|\mathsf M_n-F_{n,R,K}\|
       \leq\varepsilon+K^{-1}(1+2\varepsilon).
\]
For every \(\eta>0\), one fixed \(K\) and one initial order
therefore give error at most \(\eta\) for every subsequent
order and every admissible grid.
The same owner proves \(F_{n,R,K}=A_nB_n\) through the
fixed vector space \(\mathbb C^{{\cal U}_H}\), with
\(|{\cal U}_H|\leq H^2\).
All seventeen core declarations in these four owners pass
focused foundational-axiom checks.
Transfer to the normalized simple prime factors, the retained
finite-mode determinant comparison and the successive
edge-disjoint cycle construction remain to be formalized.
This completed arithmetic input is not a completed packing
certificate.

\begin{theorem}[Explicit positive local resilience]
\label{thm:resilience}
For every sufficiently large even \(N\), if an arbitrary set
of edges \(F\) of \(G_N\) has
\[
 \max_v\deg_F(v)\leq \frac1{1000}\frac N{\log N},             \tag{32}
\]
then \(G_N-F\) is Hamiltonian.
\end{theorem}
\begin{proof}
Write \(D=2m/\log(2m)\), \(N=2m\), and choose the order
large enough that the ordinary degrees lie in
\([0.999D,1.001D]\), the full cut bound (28) holds, and
\[
 \|U_{2m}-P_{0,2m}\|\leq0.501,\qquad
 \|U_{2m}\|\leq1.001.
\]
Let \(Q\) be the deleted edges in the canonical odd/even
cross block. Every row and column of \(Q\) has at most
\(D/1000\) edges.

Here is the complete deletion accounting. For arbitrary
\(S,T\) in the two parts, write \(a_E,b_E\) for the two
complementary crossing counts in edge set \(E\).
Their exact finite identities are
\[
 a_G=a_{G\setminus Q}+a_Q,\qquad
 b_G=b_{G\setminus Q}+b_Q.
\]
Counting the deleted boundary from either side gives
\[
 a_Q+b_Q\leq\frac D{1000}
 \min\{|S|+|T|,\ |S^c|+|T^c|\}.
\]
Thus the remaining degrees lie in
\([0.998D,1.001D]\subseteq[0.998D,1.002D]\), and the
remaining full cut coefficient is at least \(0.199\).
Lemma~\ref{lem:regularfactor} applies with
\(\varepsilon=0.002\), \(g=0.199\).
For \(D\geq1000\), it permits
\[
 r=\lfloor(39/40)D\rfloor,\qquad 0.97D\leq r\leq D,
 \qquad 39/40\leq195/199=1-\frac{2(0.002)}{0.199}.
\]
Obtain a simple \(r\)-regular factor \(H\subseteq G_N-Q\).
Measure the full discarded adjacency from the original graph:
\[
 \|A_{G_N}-A_H\|\leq1.001D-r.
\]
This includes \(Q\) once. No independent additional loss
for \(Q\) is charged.

Both normalization terms are retained:
\[
 \frac{A_H}{r}-U_{2m}
 =\frac{A_H-DU_{2m}}r+\frac{D-r}rU_{2m}.
\]
Consequently
\[
 \left\|\frac{A_H}r-U_{2m}\right\|
 \leq \frac{1.001D-r+1.001(D-r)}r
 =2.002\,\frac Dr-2.001<0.064,
\]
and hence
\[
 \left\|\frac{A_H}r-P_{0,2m}\right\|
       \leq0.565<3/5.
\]
The cross block \(T\) of \(A_H/r\) is doubly stochastic.
Put \(a=(3/5)\log(2m)\). Its entries satisfy
\[
 0\leq T_{ij}\leq1/r\leq a/m,\qquad
 a-1\leq(3/5)\log m.
\]
The exact centered cross norm equals the preceding physical
norm. The \(3/5\) specialization of
Theorem~\ref{thm:criticalrelative} therefore gives a relative
Hamilton-cycle error at most \(m^{-3/5}<1\).
The resulting cycle weight is strictly positive.
A positive term in the finite permutation sum supplies a
one-orbit permutation, and every edge of that permutation
belongs to \(H\). Through the exact odd/even vertex
bijection this is a Hamilton cycle on the original integers,
using no deleted edge. All thresholds precede the choice of
\(Q\), so the conclusion is uniform over (32).
\end{proof}

The complete resilience consumer is formalized in
\texttt{HamiltonianPrimeResilience}.
\texttt{HamiltonianFactorDeletion} proves the exact finite
partition of every crossing count, the two-sided deleted
boundary bound, and extraction from the remaining graph,
while measuring the total adjacency loss from the original
graph. \texttt{HamiltonianResilientNormalization} proves
the floor constraints and the full \(0.064\) budget.
\texttt{HamiltonianRegularFactorKernels} constructs the
actual regular cross kernel, proves its exact stochastic
law and physical centered-norm identity, and extracts an
actual one-orbit cycle from a positive signed count.
The final theorem keeps every original distinct-integer
prime-sum edge condition and proves that no deleted edge is
used. All nineteen core declarations in these four owners
pass focused foundational-axiom checks. This proof consumes
the signed counting theorem directly; the cycle-packing
theorem still has its own formalization obligation.

\section{An exact classification for congruence-filtered prime sums}
\label{sec:filters}
Fix an even integer \(q\geq2\) and a nonempty subset
\(R\subseteq(\mathbb Z/q\mathbb Z)^\times\).
For each even \(N\), let \(G_N(q,R)\) join distinct integers
\(i,j\in[1,N]\) precisely when \(i+j\) is prime and
\(i+j\bmod q\in R\).
The modulus and the set \(R\) are fixed before \(N\) tends to infinity.
Put
\[
 \sigma_R=\frac{|R|}{\phiG(q)},\qquad
 g_R=\gcd\bigl(q,\{r-r':r,r'\in R\}\bigr).
 \tag{33}
\]
For a singleton set the difference list consists of zero,
so \(g_R=q\). Representatives do not affect the gcd.
Every element of \(R\) is odd, so \(g_R\) is even.

Write
\[
 \rho_R(b)=\frac1{|R|}\sum_{r\in R}\e(br/q),\qquad
 p_q=\min\{p:p\ {\rm prime},\ p\nmid q\},
\]
and, when \(g_R=2\), define the positive explicit spectral gap
\[
 \delta_R=1-\max\left\{\frac1{p_q-1},
       \max_{\substack{b\bmod q\\b\ne0,\ q/2}}|\rho_R(b)|
                  \right\}>0.                              \tag{34}
\]
An empty inner maximum is zero.

\begin{theorem}[The finite obstruction is complete]
\label{thm:filterclassification}
For fixed \(q,R\), the following are equivalent:
\begin{enumerate}
\item \(g_R=2\).
\item \(G_N(q,R)\) is connected for every sufficiently large even \(N\).
\item \(G_N(q,R)\) is Hamiltonian for every sufficiently large even \(N\).
\end{enumerate}
When these hold, \(G_N(q,R)\) has a spanning regular subgraph
of degree
\[
 r_N=(1-o(1))\sigma_R\frac N{\log N}
\]
with nontrivial spectral ratio at most \(1-\delta_R+o(1)\).
For every fixed \(0<\kappa<\delta_R/4\) it contains at least
\[
 \left\lfloor\kappa\sigma_R\frac N{\log N}\right\rfloor
 \tag{35}
\]
edge-disjoint Hamilton cycles.
It remains Hamiltonian after arbitrary edge deletions of
maximum degree at most
\[
 \frac{\delta_R^2}{1000}\sigma_R\frac N{\log N}.              \tag{36}
\]
If \(g_R>2\), the graph is disconnected for every \(N\geq g_R\).
\end{theorem}

\begin{lemma}[The complete signed Fourier model]
\label{lem:filtermodel}
Let \(\Gamma=\mathbb Q/\mathbb Z\).
The compact self-adjoint model for the normalized weight
\[
 A_N^R(i,j)=\frac{\phiG(q)}{|R|N}
                     \Lambda(i+j)\one_{\,i+j\bmod q\in R}
\]
acts on the orthonormal basis \((e_\gamma)_{\gamma\in\Gamma}\) by
\[
 C_Re_\gamma=w_R(\gamma)e_{-\gamma}.
 \tag{37}
\]
The coefficient is zero unless
\[
 \gamma=b/q+a/c,\qquad
 (c,q)=1,\quad c\ {\rm squarefree},\quad(a,c)=1,
\]
where \(c=1,a=0\) is included.
In the nonzero cases its exact value is
\[
 w_R(\gamma)=\rho_R(b)\frac{\mu(c)}{\phiG(c)}.                \tag{38}
\]
If \(g_R=2\), the only eigenvalues of modulus 1 are the simple
constant and parity eigenvalues \(1,-1\).
Every other eigenvalue has modulus at most \(1-\delta_R\).
The sampled operator has the same spectral bounds with an \(o(1)\)
error.
\end{lemma}
\begin{proof}
One can view (38) as a finite character computation, with no
hypothesis about unproved prime correlations.
The limiting prime residue law is the product of the normalized
uniform unit laws at each prime power, conditioned on its residue
modulo \(q\) lying in \(R\).
At primes dividing \(q\), additive characters of conductor
larger than the corresponding prime-power factor of \(q\)
average to zero over the uniform lifts.
At any prime not dividing \(q\), a primitive additive character
averages to \(-1/(p-1)\) at conductor \(p\), and to zero at
conductor \(p^k\), \(k\geq2\).
The Chinese remainder theorem makes these latter coordinates
independent of the residue modulo \(q\).
The remaining finite average at \(q\) is exactly \(\rho_R(b)\).
This proves every zero coefficient and the entire signed formula
(38). In particular
\(w_R(-\gamma)=\overline{w_R(\gamma)}\), so (37) is self-adjoint;
its two-dimensional reflection blocks have eigenvalues
\(\pm|w_R(\gamma)|\).
Conductor tails tend to zero, since \(q\) is fixed and
\(\phiG(c)\geq\sqrt{c/2}\), so the model is compact.

Here are details linking this exact model to the sampled matrix.
Let \(C\) be the unfiltered reflection model with coefficients
\(\mu(k)/\phiG(k)\) at reduced conductor \(k\).
For \(t=b/q\), let \(U_t e_\gamma=e_{\gamma+t}\).
The finite Fourier expansion of the residue mask gives the
operator identity
\[
 C_R=\frac{\phiG(q)}{|R|q}
       \sum_{b\bmod q}
       \left(\sum_{r\in R}\e(-br/q)\right)U_{b/q}CU_{b/q}.
 \tag{39}
\]
Indeed the coefficient at \(\gamma\) in its reflection is the
finite masked average of \(\e(\gamma t)\) under the prime
residue law just computed.
Notice the same shift on both sides of \(C\), rather than an
adjoint conjugation: this is a sum kernel.

Truncate \(C\) to reduced conductors \(k\leq U\).
The operator tail is at most \(\sqrt{2/U}\).
The finite sum (39) therefore has tail \(O_{q,R}(U^{-1/2})\).
All frequencies in its truncated range and domain have
denominator at most \(qU\).
Their sampled Gram matrix differs from the identity by
\(O_q(U^4/N)\), by geometric summation as in
Lemma~\ref{lem:integerspectrum}.
On the integer samples the same finite Fourier expansion
takes \(A_N\) to \(A_N^R\): multiplication by
\(\e(b(i+j)/q)\) is a diagonal unitary on each side.
Lemma~\ref{lem:integertransfer}, with \(U\) a sufficiently
large fixed power of \(\log N\), bounds the unfiltered
operator error by any prescribed negative power of \(\log N\).
The finite mask expansion multiplies this bound only by a
constant depending on \(q,R\).
The Gram bound and the truncated model tail tend to zero.
Polar orthogonalization and Weyl's inequality therefore
transfer the spectrum of \(C_R\), with appended zeros, to
\(A_N^R\). This also justifies the model without appealing to
a pointwise infinite Ramanujan expansion.

It remains to locate its unit-modulus coefficients.
A finite average of unit complex numbers has modulus 1
exactly when all its terms agree. Thus
\[
 |\rho_R(b)|=1
 \quad\Longleftrightarrow\quad
 q\mid b(r-r')\ \hbox{for all }r,r'\in R
 \quad\Longleftrightarrow\quad
 b\equiv kq/g_R\pmod q
\]
for some integer \(k\).
When \(g_R=2\), these are just \(b=0,q/2\), with averages
\(1,-1\).
For \(c>1\), (38) has modulus at most \(1/(p_q-1)\).
Every remaining coefficient is bounded by (34).
This proves all the assertions.
\end{proof}

\begin{proof}[Proof of Theorem~\ref{thm:filterclassification}]
First suppose \(g_R>2\), and fix \(r_0\in R\).
Every edge satisfies
\[
 i+j\equiv r_0\pmod {g_R}.
\]
Since \(g_R\) is even and \(r_0\) odd, the involution
\(a\mapsto r_0-a\) on residues modulo \(g_R\) has no fixed point.
Its \(g_R/2\) distinct pairs give disjoint vertex sets with
no edge between them. All are nonempty once \(N\geq g_R\).
There are at least two such sets, so the graph is disconnected.

Now suppose \(g_R=2\).
Remove proper prime powers and diagonal entries from \(A_N^R\),
and denote the resulting true prime weight matrix by \(Q_N^R\).
The removed maximum row sum is
\(O_{q,R}(N^{-1/2}(\log N)^2)\).
The prime number theorem in the fixed reduced classes modulo
\(q\), with arbitrary logarithmic accuracy, gives uniformly
\[
 \|Q_N^R\one-\one\|_\infty\ll_{q,R,B}(\log N)^{-B}
 \quad\hbox{for every fixed }B.
\]
Indeed the exact raw row is the normalized difference of the
restricted Chebyshev function at \(i+N\) and \(i\).
The normalized spectrum has the two simple exceptions and
gap \(\delta_R-o(1)\), by Lemma~\ref{lem:filtermodel}.
All entries lie between opposite parities, since \(q\) is even
and \(R\) contains only units.
Lemma~\ref{lem:scaling} therefore gives a true-support
stochastic \(P_N^R\), with
\[
 \|P_N^R-Q_N^R\|=o(1),\qquad
 \|P_N^R|_{\{\one,s\}^\perp}\|\leq1-\delta_R+o(1).
\]
Put \(D=\sigma_R N/\log N\), and let \(B_N^R\) be the
ordinary adjacency.
The row estimate in Lemma~\ref{lem:unweightedtransfer}, with
the fixed normalization factor \(1/\sigma_R\), gives
\[
 \max_i|\deg(i)/D-1|=o(1),\qquad
 \|B_N^R/D-P_N^R\|=o(1).                                    \tag{40}
\]
The Laplacian calculation proving (28) now gives the cut
constant \(g=\delta_R/3\) for all sufficiently large \(N\).
Lemma~\ref{lem:regularfactor} gives an
\(r_N=(1-o(1))D\) factor, discarding only \(o(D)\)
edges at each vertex.
Equation (40) transfers the nontrivial spectral bound to this
factor. Its polynomial-order degree dominates the fixed
polylogarithmic threshold in (2), giving a Hamilton cycle.
The iterative estimate (31), with initial gap \(\delta_R\)
in place of \(1/2\), proves (35).
Hamiltonicity implies connectedness, completing the equivalence.

For completeness we verify the uniform deletion constant (36).
Write \(\delta=\delta_R\) and \(\rho=\delta^2/1000\).
Choose \(N\) so large that the initial normalized degree defect
is at most \(\rho\), the initial second eigenvalue is at most
\((1-\delta+\rho)D\), and the cut constant is \(\delta/3\).
After deletion of maximum degree \(\rho D\), the degree defect
is at most \(2\rho\), and the cut constant is at least
\(\delta/3-\rho\geq\delta/4\).
The regular-factor lemma gives
\[
 r=\left\lfloor(1-16\rho/\delta)D\right\rfloor
                     \geq(1-17\rho/\delta)D
\]
for sufficiently large \(D\).
The total discarded degree from the original graph is at most
\(18\rho D/\delta\). Hence its factor satisfies
\[
 r-\lambda_2(H)
 \geq\left(\delta-\rho-\frac{35\rho}{\delta}\right)D
 \geq \frac{\delta}{2}D.
\]
The factor is regular, bipartite, and of degree comparable to
\(D\), so (2) applies again. This proves (36).
\end{proof}

\begin{corollary}[Arbitrarily small fixed prime density suffices]
\label{cor:sparseprimeclasses}
For every fixed even \(q\geq4\), permit only prime sums
congruent to \(1\) or \(-1\) modulo \(q\).
The resulting graphs are Hamiltonian at every sufficiently
large even order.
The permitted primes have relative density \(2/\phiG(q)\)
among all primes. This can be made arbitrarily small by
choosing \(q\) first and then allowing \(N\) to tend to infinity.
\end{corollary}
\begin{proof}
The difference of the two residues is 2, so \(g_R=2\).
The classification applies, and the fixed-progression prime
number theorem gives the stated density.
There is no claim uniform in a growing modulus.
\end{proof}

\begin{remark}[What the criterion detects]
A single reduced residue class modulo even \(q>2\) fails the
criterion and forces disconnection. Two opposite classes pass.
Thus density alone is not the deciding invariant.
The decisive datum is the subgroup generated by differences
of allowed prime residues, and the retained Fourier phases
compute its obstruction exactly.
\end{remark}

\begin{remark}[Equal density, different outcomes, and a nonreal phase]
For \(q=8\), all three rows below retain exactly half of
the primes asymptotically.
\[
 \begin{array}{c|c|c|c}
 R&\sigma_R&g_R&\hbox{eventual outcome}\\ \hline
 \{1,3\}&1/2&2&\hbox{Hamiltonian}\\
 \{1,5\}&1/2&4&\hbox{disconnected}\\
 \{1,7\}&1/2&2&\hbox{Hamiltonian}
 \end{array}
\]
For the first row, the exact coefficient at \(b=1\) is
\[
 \rho_{\{1,3\}}(1)
 =\tfrac12\bigl(\e(1/8)+\e(3/8)\bigr)
 =\mathrm{i}/\sqrt2 .
\]
Its reflection pair has eigenvalues \(\pm1/\sqrt2\);
the imaginary phase is retained before the eigenvalues are
taken. Here \(\delta_R=1-1/\sqrt2\) and the counting law is
\[
 h(G_N(8,\{1,3\}))
   =\left(\frac{N}{2e\log N}\right)^N\exp(o(N)).
\]
The middle row instead has no Hamilton cycle for any
\(N\geq4\), because its two residue-pair components cannot
communicate.
\end{remark}

\section{Prescribed spanning paths and arbitrary endpoints}
\label{sec:paths}
A Hamilton cycle does not itself give a Hamilton path with
arbitrarily specified endpoints.
We verify the stronger spanning-path input separately.

We use Theorem 6.4 and Definition 3.9 of~\cite{BJ}.
For a regular balanced bipartite graph of degree \(d\), their
bi-uniformity condition for \(S\) in one part is
\[
 \left||N(v)\cap S|-\frac{2|S|d}{N}\right|\leq\alpha d
 \quad\hbox{at every vertex of the other part}.              \tag{41}
\]
The theorem supplies internally disjoint paths joining prescribed
opposite-part pairs, with all remaining vertices in their interiors,
provided the endpoint sets and multisets obey (41), have equal
cardinality, and the following conditions hold:
\[
 0<c<1/4,\quad 0<\gamma<(\log N)^{-2},\quad
 d'\leq d^{1-c},\quad\alpha\leq d^{-c},\quad
 m\leq d^{-c/900}N,\quad d\geq\gamma^{-10^6/c}.               \tag{42}
\]
Here the graph is \((d\pm d')\)-nearly regular and a
\(\gamma\)-expander in the sense used in (2), and \(N\)
is sufficiently large depending on \(c\).
We apply it to exact regular factors, so \(d'=0\).

\begin{lemma}[Arbitrary small endpoint sets meet the source condition]
\label{lem:smallendpoints}
Let \(H\) be a simple \(d\)-regular balanced bipartite graph
on \(N\) vertices.
Any subset \(S\) of one part of cardinality \(m\) obeys
(41) with \(\alpha=m/d\).
The same holds when its distinct elements are viewed as a multiset.
\end{lemma}
\begin{proof}
Both \(|N(v)\cap S|\) and \(2md/N\) lie in \([0,m]\),
since \(d\leq N/2\). Their difference has modulus at most \(m\).
This is (41), including the set and multiset versions.
\end{proof}

\begin{theorem}[Almost-linear prescribed spanning linkages]
\label{thm:linkages}
Fix \(q,R\) as in Theorem~\ref{thm:filterclassification},
with \(g_R=2\), and fix \(0<\eta<1\).
For all sufficiently large even \(N\), choose any integer
\(1\leq m\leq N^{1-\eta}\), any distinct odd vertices
\(a_1,\ldots,a_m\), and any distinct even vertices
\(b_1,\ldots,b_m\).
Then \(G_N(q,R)\) contains pairwise vertex-disjoint paths
\(P_1,\ldots,P_m\), where \(P_i\) joins \(a_i\) to \(b_i\)
and their vertex sets partition all of \([1,N]\).
In particular it has a Hamilton path between every prescribed
odd-even pair.
\end{theorem}
\begin{proof}
Use the exact regular factor from
Theorem~\ref{thm:filterclassification}.
Its degree is \(d=(1-o(1))\sigma_R N/\log N\),
and its fixed positive spectral gap supplies a fixed positive
edge-expansion constant \(\gamma_0\), under the definition (2).
Choose
\[
 c=\min\{1/5,\eta/2\},\qquad
 \gamma=(\log N)^{-3},\qquad d'=0,\qquad\alpha=m/d .
\]
For large \(N\), \(\gamma<\gamma_0\), so expansion with this
smaller constant holds throughout the required cut range.
The arbitrary endpoint sets meet (41), by
Lemma~\ref{lem:smallendpoints}, and each has cardinality \(m\).
Their vertices are distinct, so the multiset requirements are
the same.

We now check every quantitative condition in (42).
The choice of \(c\) is in \((0,1/4)\), and the choice of
\(\gamma\) is smaller than \((\log N)^{-2}\).
Also
\[
 m/d^{1-c}
 =O_{q,R}\left(N^{c-\eta}(\log N)^{1-c}\right)\longrightarrow0,
\]
since \(c<\eta\).
Thus \(\alpha=m/d\leq d^{-c}\).
Similarly
\[
 \frac{m}{d^{-c/900}N}
 =O_{q,R}\left(N^{-\eta+c/900}
                    (\log N)^{-c/900}\right)\longrightarrow0.
\]
Finally
\[
 d\asymp_{q,R}N/\log N
       \gg(\log N)^{3\cdot10^6/c}
       =\gamma^{-10^6/c}.
\]
The degree error \(d'=0\) obeys its condition exactly.
The cited theorem gives the desired paths.
Their interiors avoid every prescribed endpoint; since all
endpoints are distinct, internal disjointness becomes full
vertex disjointness, and its covering conclusion partitions
the entire vertex set.
\end{proof}

\begin{corollary}[All parity-compatible prime-sum endpoints]
\label{cor:oddendpoints}
For all sufficiently large \(N\), in the ordinary prime-sum
graph on \([1,N]\), every distinct pair \(a,b\) satisfying
the necessary parity condition is joined by a Hamilton path:
opposite parities for even \(N\), and both odd for odd \(N\).
\end{corollary}
\begin{proof}
The even case is Theorem~\ref{thm:linkages} with \(q=2\)
and \(R=\{1\}\).
Suppose \(N\) is odd and \(a,b\) are distinct odd vertices.
Choose an even neighbour \(v\) of \(a\) in \(G_N\).
Such a neighbour exists for every large \(N\), uniformly in
\(a\): the degree estimate (26) for \(G_{N+1}\) loses at
most one edge on deleting vertex \(N+1\).

Consider \(G_N-a\), a balanced bipartite graph on \(N-1\)
vertices. It is obtained from \(G_{N+1}\) by deleting the
odd vertex \(a\) and the even vertex \(N+1\).
Its degrees remain \((1+o(1))(N+1)/\log(N+1)\),
uniformly in the two deleted vertices.
If \(W\) is a subset of the remaining vertices of size at
most \((N-1)/2\), the cut bound (28) in \(G_{N+1}\) loses
at most \(2|W|\) edges. Thus its boundary is at least
\[
 \left(\frac15\frac{N+1}{\log(N+1)}-2\right)|W|.
\]
The complementary smaller-side bound is identical.
The regular-factor lemma supplies a regular spanning factor
of degree \(d=(1-o(1))(N+1)/\log(N+1)\), discarding only
\(o(d)\) edges per vertex.
It retains a fixed positive edge-expansion constant.
The argument proving Theorem~\ref{thm:linkages}, or (42)
with \(m=1\), supplies a Hamilton path from \(v\) to \(b\)
in this factor. Prepend the edge \(av\).
The resulting path begins at \(a\), ends at \(b\), and
covers exactly all \(N\) vertices of \(G_N\).
\end{proof}

\begin{corollary}[Diffuse kernel supports admit prescribed spanning linkages]
\label{cor:weightedlinkages}
Fix \(0<a\leq1\), \(C>0\), \(0<\delta<1\), and
\(0<\zeta<a\).
For all sufficiently large \(N\), every stochastic bipartite
kernel \(P\) with nontrivial gap at least \(\delta\) and
\(\alpha\leq C N^{-a}\) has the following property:
any \(m\leq N^{a-\zeta}\) prescribed distinct opposite-part
pairs can be joined by vertex-disjoint paths partitioning
all \(N\) vertices of its support.
\end{corollary}
\begin{proof}
In Theorem~\ref{thm:rounding}, choose
\[
 d=\lfloor N^{a-\zeta/2}\rfloor.
\]
For all large \(N\), \(d\leq N/2\) and
\(d\alpha L^2=O(N^{-\zeta/2}L^2)\to0\).
The entire error in (1) tends to zero, so its simple regular
factor has nontrivial gap at least \(\delta/2\).
Choose
\[
 c=\min\{1/5,\zeta/(4a)\},\qquad\gamma=(\log N)^{-3}.
\]
As before take \(d'=0\), \(\alpha'=m/d\).
The arbitrary endpoint sets satisfy (41) with \(\alpha'\).
Since
\[
 (a-\zeta/2)(1-c)>(a-\zeta),
\]
we have \(m\leq d^{1-c}\) for all large \(N\).
Moreover
\[
 a-\zeta+(a-\zeta/2)c/900
 \leq 1-\zeta+\zeta/3600<1,
\]
so \(m\leq d^{-c/900}N\).
The polynomial degree dominates
\(\gamma^{-10^6/c}\).
Every hypothesis of (42) is satisfied.
The same spanning-path theorem therefore applies.
\end{proof}

\section{Hamilton-cycle entropy with a fixed spectral gap}
\label{sec:enumeration}
This section proves a counting result under a fixed spectral gap;
no assumption that the nontrivial spectral ratio tends to zero
is made. For comparison, the bipartite enumeration theorem of
Byrne--Tait~\cite[Theorem 7]{BT}, extending Krivelevich~\cite{Krivelevich},
requires \(d/\lambda\geq(\log N)^{1+\epsilon}\).
Our explicit polylogarithmic degree threshold permits a
nontrivial ratio arbitrarily close to 1.

For a simple graph \(G\), let \(h(G)\) count unoriented Hamilton
cycles, with cyclic rotations identified.
Let \(J_m\) denote the \(m\)-by-\(m\) all-ones matrix.

\begin{theorem}[Full exponential Hamilton entropy at fixed gap]
\label{thm:hamiltonentropy}
Fix \(0<\delta<1\).
Let \(G=G_N\) be a sequence of simple balanced bipartite
graphs on \(N=2m\) vertices, with biadjacency matrix \(B\).
Suppose there are positive numbers \(D=D_N\) such that
\[
 \begin{gathered}
 D\geq(\log N)^{K_0},\qquad K_0=2\cdot10^7,\\
 \max_v|\deg_G(v)/D-1|=o(1),\qquad
 \left\|B/D-J_m/m\right\|\leq1-\delta .
 \end{gathered}
 \tag{43}
\]
Then
\[
 \log h(G)=N(\log D-1)+o(N),\qquad
 h(G)=(D/e)^N\exp(o(N)).                                    \tag{44}
\]
The logarithmic assertion includes eventual positivity.
The error is uniform along any sequence meeting the displayed
degree defect and fixed parameters.
\end{theorem}

We first supply the random-restriction estimate used in its proof.
This is a standard matrix-concentration mechanism, with
related formulations in~\cite[Lemmas 2.2--2.4]{Muyesser}
and the earlier work credited there.
We give the particular estimate needed below, without
claiming a new sampling principle or independence of
arithmetic entries.

\begin{lemma}[Two independent coordinate restrictions]
\label{lem:randomrestriction}
Let \(E\) be an \(m\)-by-\(m\) real or complex matrix, and put
\[
 v=\max\left\{\max_i\|E_{i,*}\|_2,\max_j\|E_{*,j}\|_2\right\},
 \qquad k=\max_{ij}|E_{ij}|,\qquad \ell=\log(4m).
\]
Choose independent row and column sets by Bernoulli sampling
with probability \(p\), with diagonal selectors \(R_X,R_Y\).
For every fixed \(A>0\), with probability at least \(1-m^{-A}\),
\[
 \|R_XER_Y\|
 \leq p\|E\|+C_A\bigl(v\sqrt{p\ell}+k\ell\bigr).             \tag{45}
\]
For independent uniform sets of common prescribed size \(b\),
take \(p=b/m\); the same bound fails with probability at most
\((m+1)^2m^{-A}\).
\end{lemma}
\begin{proof}
Write \(c_j\) for the columns of \(E\).
For the first column restriction the centered rank-one sum has
increments of norm at most \(v^2\) and matrix variance at most
\[
 p\sum_j\|c_j\|_2^2c_jc_j^*
             \preceq pv^2EE^*.
\]
Matrix Freedman applied to these independent increments gives,
with arbitrarily prescribed polynomially high probability,
\[
 \|ER_YE^*-pEE^*\|
 \leq C_A\bigl(v\|E\|\sqrt{p\ell}+v^2\ell\bigr).
\]
Taking square roots, and increasing \(C_A\), yields
\[
 \|ER_Y\|\leq\sqrt p\|E\|+C_Av\sqrt\ell.                     \tag{46}
\]
For each row, the sum
\(\sum_j\one_{j\in Y}|E_{ij}|^2\) has mean at most \(pv^2\),
bounded summands \(k^2\), and scalar variance at most
\(pk^2v^2\).
The scalar version of the same inequality and a union bound
over rows give
\[
 \max_i\|(ER_Y)_{i,*}\|_2\leq\sqrt p\,v+C_Ak\sqrt\ell.        \tag{47}
\]
Condition on \(Y\), and apply the rank-one matrix argument to
the independent row restriction of \(Z=ER_Y\).
It gives
\[
 \|R_XZ\|\leq\sqrt p\|Z\|
                   +C_A\sqrt\ell\max_i\|Z_{i,*}\|_2.
\]
Substitute (46)--(47), retaining both error terms, to obtain
(45). Choosing the failure exponents a little larger before
the three union bounds gives the stated exponent \(A\).

If \(p=b/m\), the event that a Bernoulli set has size exactly
\(b\) has probability at least \(1/(m+1)\): \(b\) is a mode
of its binomial distribution.
Conditional on that event the set is uniform.
The two independent conditioning events therefore cost at
most \((m+1)^2\) in failure probability.
The cases \(b=0,m\) follow directly, or can be omitted in the
application where \(0<b<m\).
\end{proof}

\begin{lemma}[Regular expanding factors in random blocks]
\label{lem:randomblockfactor}
Assume the uniform-degree and centered-norm hypotheses
of (43), without imposing its global lower bound on \(D\).
Let \(X,Y\) be independent uniform subsets of the two parts,
each of size \(s\), and put \(p=s/m\).
If \(pD/\log N\to\infty\), then, with probability
\(1-N^{-A}\) for every fixed prescribed \(A\), the block
\(G[X,Y]\) has a regular spanning factor of degree
\[
 d_s=(1-o(1))pD
\]
and a fixed positive nontrivial gap depending only on \(\delta\).
All vertex degrees before factoring are \((1+o(1))pD\).
If the host is exactly \(D\)-regular, the factor can be chosen
so that its degree differs from \(pD\), and its discarded
maximum degree is bounded, by
\[
 C_{\delta,A}\bigl(\sqrt{pD\log N}+\log N\bigr).
 \tag{48}
\]
The assertions hold simultaneously for polynomially many
specified block pairs of a random ordered partition.
A common factor degree may be used for pairs with the same size.
\end{lemma}
\begin{proof}
Use \(E=B/D-J_m/m\).
The degree condition in (43) and \(D\leq(1+o(1))m\)
give
\[
 v=O(D^{-1/2}),\qquad k=O(D^{-1}).
\]
Lemma~\ref{lem:randomrestriction}, conditioned on the two
prescribed sizes, gives
\[
 \left\|\frac{B[X,Y]}{pD}-\frac{J_s}{s}\right\|
 \leq1-\delta+
        O_A\left(\sqrt{\frac{\log N}{pD}}
                        +\frac{\log N}{pD}\right).
 \tag{49}
\]
Sampling neighbours without replacement gives, uniformly for
all vertices in the original graph,
\[
 \deg(v,Y)\ \hbox{or}\ \deg(v,X)
 =pD\left(1+O\left(\varepsilon_N+
       \sqrt{\frac{\log N}{pD}}+\frac{\log N}{pD}\right)\right),
 \tag{50}
\]
where \(\varepsilon_N=\max_v|\deg(v)/D-1|\).
The usual scalar concentration inequality follows also by
Bernoulli sampling and the same conditioning argument.
The constants can be chosen for any polynomial failure exponent.

For large \(N\), (49) is at most \(1-\delta/2\) and the
degree defect in (50) tends to zero.
The Laplacian argument proving (28), with the bipartite
constant block \(J_s/s\), now gives a fixed positive cut
constant, for example \(g=\delta/8\).
Lemma~\ref{lem:regularfactor} extracts a factor of degree
\((1-o(1))pD\), deleting \(o(pD)\) edges per vertex.
The discarded adjacency has the same operator bound;
(49) transfers a fixed positive nontrivial gap to the factor.
All errors depend on the displayed deterministic defect and
on \(pD\), so one may take their common upper bound before
choosing the integer factor degree.
For an exactly regular host, (50) has no \(\varepsilon_N\)
term. The explicit loss in (22)--(23), including its floor,
is then at most a constant times
\(\sqrt{pD\log N}+\log N\), which proves (48).
The marginals of any two opposite-part blocks in a random
ordered partition are independent uniform subsets.
A union bound with a larger prescribed exponent proves the
simultaneous assertion; independence between different block
pairs is not required.
\end{proof}

\begin{proof}[Proof of Theorem~\ref{thm:hamiltonentropy}]
We first regularize the host before sampling.
The centered norm in (43) and its vanishing degree defect
give a fixed positive cut constant by the Laplacian argument
used in (28).
Lemma~\ref{lem:regularfactor} gives an exactly \(d\)-regular
spanning subgraph \(H\), where
\[
 d=(1-o(1))D,\qquad
 \|A_G-A_H\|=o(D).
\]
Consequently its biadjacency satisfies
\(\|B_H/d-J_m/m\|\leq1-\delta/2\) for all large \(N\).
Also \(d\geq\tfrac12(\log N)^{K_0}\) eventually.
For the lower bound work in this exact regular host;
every resulting cycle belongs to the original graph.
The difference \(\log d-\log D=o(1)\) will cost only \(o(N)\).

Take
\[
 t=\lfloor d^{1/4}\rfloor,\qquad
 b=\left\lfloor\frac{m-m/(\log N)^3}{t}\right\rfloor,\qquad
 r=m-tb.
\]
These parameters satisfy
\[
 t\asymp d^{1/4},\quad b\asymp N d^{-1/4},\quad
 r=\frac{m}{(\log N)^3}+O(d^{1/4}),\quad
 (r+b)\log N=o(N),\quad t\log N=o(N).                        \tag{51}
\]
Partition the first part into ordered sets
\(A_1,\ldots,A_t,R_A\), of sizes \(b,\ldots,b,r\),
and independently the other part into
\(B_1,\ldots,B_t,R_B\) with the same sizes.
The number of ordered partition pairs is
\[
 {\cal P}=\left(\frac{m!}{r!(b!)^t}\right)^2.
\]
Stirling's formula and (51) give, with \(p=b/m\),
\[
 \log{\cal P}=2tb\log(1/p)+o(N).                            \tag{52}
\]

Require the regular-factor conclusions in every adjacent
pair of the alternating list
\[
 A_1,B_1,A_2,B_2,\ldots,A_t,B_t.
\]
Their sampled degree is \(pd\asymp d^{3/4}\),
so it exceeds \((\log N)^{3K_0/4}\) up to a fixed factor.
Require the same conclusions on
\[
 K=H[R_A\cup A_1,\ R_B\cup B_t],\qquad
 n_K=2(r+b),\qquad p_K=(r+b)/m .
\]
The augmented reservoir has independent uniform part sets
of size \(r+b\), and \(p_Kd\asymp d/(\log N)^3\).
In addition require the scalar endpoint concentration
\[
 |\deg_H(v,A_1)-pd|,\ |\deg_H(v,B_t)-pd|
       \leq C\bigl(\sqrt{pd\log N}+\log N\bigr)              \tag{53}
\]
for every vertex \(v\) of the appropriate opposite part.
This holds by sampling without replacement in the exact
regular host.
There are only \(2t\) block requirements and \(2N\) scalar
requirements. Lemma~\ref{lem:randomblockfactor} and scalar
concentration, with suitably larger failure exponents,
therefore hold for a proportion \(1-o(1)\) of all partitions.
Their events depend on the partition, not on any matching.

Fix a good partition.
For every consecutive block choose a factor with the same
degree \(d_0=(1-o(1))pd\).
The van der Waerden permanent bound gives at least
\[
 d_0^b\,\frac{b!}{b^b}\geq(d_0/e)^b
\]
perfect matchings in each.
This established permanent inequality, together with
Bregman's upper bound used below, is recorded in
\cite[Section 2.4]{BT}.
Choose these \(2t-1\) matchings as independent combinatorial
choices. Their union covers the \(2tb\) nonreservoir vertices
by \(b\) paths, each with \(2t\) vertices.
The first endpoints are \(A_1\), the last endpoints \(B_t\).
The layer structure prevents cycles, and each resulting path
cover determines all its chosen matchings.

We verify the source hypotheses for closing every cover.
Choose in \(K\) an exact regular factor \(K'\) of degree \(d_K\)
using (48), so that
\[
 |d_K-p_Kd|,\quad
 \max_v\deg_{K-K'}(v)
       \leq C_\delta\bigl(\sqrt{p_Kd\log N}+\log N\bigr).
 \tag{54}
\]
It has fixed positive nontrivial gap.
The endpoint sets can be too large for the arbitrary-set
bound in Lemma~\ref{lem:smallendpoints}; their distribution
must be retained instead.
Combining (53)--(54), for every \(v\) of the opposite part,
\[
 \left|\deg_{K'}(v,A_1)-\frac b{r+b}d_K\right|,
 \quad
 \left|\deg_{K'}(v,B_t)-\frac b{r+b}d_K\right|
 \leq C_\delta\bigl(\sqrt{p_Kd\log N}+\log N\bigr).           \tag{55}
\]
Indeed the target before factoring is \(pd=(b/(r+b))p_Kd\).
The three errors are the original endpoint fluctuation,
the actual number of edges lost at the vertex, and
\((b/(r+b))|d_K-p_Kd|\).
Each is bounded by the right side, since \(p\leq p_K\).
Thus no endpoint compensation is erased by regularization.

The two distinct endpoint sets, both of size \(b\), satisfy
the set and multiset bi-uniformity condition (41), now on
\(K'\), with
\[
 \alpha=C_\delta\left(
       \sqrt{\frac{\log N}{p_Kd}}+\frac{\log N}{p_Kd}\right).
\]
Take
\[
 c=1/5,\qquad \gamma=(\log n_K)^{-3},\qquad d'=0 .
\]
The degree satisfies
\(d_K\asymp d/(\log N)^3\gg(\log N)^{K_0-3}\).
In particular
\[
 \alpha d_K^c
 \ll_\delta(\log N)^{1/2}d_K^{c-1/2}
                       +(\log N)d_K^{c-1}\longrightarrow0.
\]
Thus \(\alpha\leq d_K^{-c}\) for all large \(N\).
For the number \(b\) of prescribed pairs, charge the full range:
\[
 \frac{b}{d_K^{-c/900}n_K}
 \ll (\log N)^3 d^{-1/4+c/900}\longrightarrow0.
 \tag{56}
\]
The limit follows from \(K_0(1/4-1/4500)>3\).
The fixed positive spectral gap gives a fixed positive
edge-expansion constant, eventually larger than \(\gamma\).
Also
\[
 \gamma^{-10^6/c}=(\log n_K)^{15\cdot10^6}
                   =o(d_K),
\]
since \(K_0-3>15\cdot10^6\).
The remaining source conditions are
\(0<c<1/4\), \(\gamma<(\log n_K)^{-2}\), equal endpoint
cardinalities, \(d'=0\leq d_K^{1-c}\), and large \(n_K\)
relative to fixed \(c\); all hold.

Label the long paths \(Q_1,\ldots,Q_b\), with endpoints
\(a_i\in A_1,z_i\in B_t\), and apply (42) to
\((a_{i+1},z_i)\), taking indices cyclically.
It supplies disjoint connecting paths with interiors
exactly \(R_A\cup R_B\).
Together with the long paths, in the prescribed cyclic order,
these form one Hamilton cycle.
For all large \(N\), \(t\geq2\), so no long-path edge has both
ends in \(K\); the two edge collections do not overlap.

Choose one such Hamilton cycle for each configuration by a
fixed deterministic rule.
For a fixed output cycle \(C\), its long-path cover in any
preimage is a subset of its \(N\) edges with \(2tb-b\) edges.
The number of possibilities is at most
\[
 \binom N{\,2r+b\,}=\exp(o(N)).                              \tag{57}
\]
Each subset determines its entire configuration:
its isolated vertices are the reservoir, and each nontrivial
component is one long path.
Orient each path from its endpoint in the first part;
successive positions fix all ordered layers and matchings.
No permutation of paths or extra orientation factor is
left uncharged.

The number of good configurations is at least
\[
 (1-o(1)){\cal P}\,(d_0/e)^{b(2t-1)}.
\]
Its logarithm, by (51)--(52), is
\[
 2tb\log(1/p)+(2tb-b)(\log(pd)-1)+o(N)
                 =N(\log d-1)+o(N).
\]
The missing \(2r+b\) matching edges cost
\(O((r+b)\log N)=o(N)\).
The ordered-partition entropy cancels the entire layer
degree loss. Division by (57), followed by
\(\log d=\log D+o(1)\), proves the required lower bound
for the original graph.

For its upper bound, the full adjacency permanent counts
all directed cycle covers and includes both orientations
of every Hamilton cycle.
Bregman's inequality gives
\[
 2h(G)\leq\operatorname{per}(A_G)
            \leq\prod_v(\deg_G(v)!)^{1/\deg_G(v)}.
\]
Stirling and (43) imply
\[
 \log\operatorname{per}(A_G)
 \leq\sum_v(\log\deg_G(v)-1)
       +O\left(\frac{N\log D}{D}\right)
 =N(\log D-1)+o(N).
\]
This completes the proof with its explicit polylogarithmic
degree threshold.
\end{proof}

\begin{corollary}[The exponential asymptotic for prime circles]
\label{cor:primecircleentropy}
Fix an even \(q\) and a nonempty allowed unit set \(R\).
If \(g_R=2\), then along the even integers
\[
 h(G_N(q,R))
 =\left(\frac{\sigma_R N}{e\log N}\right)^N\exp(o(N)).
 \tag{58}
\]
If \(g_R>2\), the count is zero for all \(N\geq g_R\).
In particular the ordinary prime-sum graph satisfies
\[
 \log h(G_N)
   =N\log N-N\log\log N-N+o(N).
 \tag{59}
\]
\end{corollary}
\begin{proof}
The norm estimate (40) and the stochastic cross-block
spectral gap imply
\[
 \left\|B_N^R/D-J_m/m\right\|
                    \leq1-\delta_R+o(1),
 \qquad D=\sigma_R N/\log N .
\]
The uniform degrees are (40), and
\(D\geq(\log N)^{K_0}\) once \(N\) is sufficiently large.
Apply Theorem~\ref{thm:hamiltonentropy} with any smaller
fixed positive gap.
The disconnected case is the exact obstruction in
Theorem~\ref{thm:filterclassification}.
For \(q=2,R=\{1\}\), the relative prime density is one,
giving (59).
\end{proof}

\begin{remark}[Precision of the enumeration]
The factor \(\exp(o(N))\) is part of the assertion.
This determines the exponential base, including its constant
\(e^{-1}\); it is not a relative-error \(1+o(1)\) formula for
the total count.
It leaves, for example, any bounded singular-series factor
and finer subexponential terms undetermined.
\end{remark}

\begin{corollary}[An entropy and concentration inverse for kernel supports]
\label{cor:entropyinverse}
Fix \(0<\delta<1\), and let \(K_0=2\cdot10^7\).
There is \(a_\delta>0\) such that, for sufficiently large even
\(N\), a stochastic bipartite kernel \(P\) with nontrivial
gap at least \(\delta\) and
\[
 \alpha=\max_{ij}P_{ij}\leq a_\delta L^{-(K_0+2)}             \tag{60}
\]
has Hamiltonian support.
More quantitatively, with \(c_\delta\) from the finite
rounding theorem and
\[
 d=\left\lfloor\frac{c_\delta}{2\alpha L^2}\right\rfloor,
 \qquad
 \log h(\supp P)\geq N(\log d-1)-o(N),                       \tag{61}
\]
the displayed lower bound holds along every sequence
satisfying (60), with fixed \(\delta\).
Consequently a non-Hamiltonian support with this fixed gap
forces \(\alpha>a_\delta L^{-(K_0+2)}\).
\end{corollary}
\begin{proof}
Stochasticity on equal parts gives \(\alpha\geq1/m=2/N\).
Decrease \(c_\delta\) if necessary so that \(c_\delta\leq1\),
and choose \(a_\delta\leq c_\delta/8\).
Then the degree in (61) is at most \(m\) and, for all
sufficiently large \(N\), is at least
\(2L^{K_0}\), hence at least \((\log N)^{K_0}\).
It obeys \(d\alpha L^2\leq c_\delta/2\).
The full rounding error in (1) is bounded by
\[
 C_\delta\left(\frac{c_\delta}{2L}
       +\sqrt{\frac Ld}+\frac Ld\right)=o(1).
\]
The resulting simple regular graph \(H\) therefore has
nontrivial gap at least \(\delta/2\), and its cross block
satisfies (43) with degree \(d\).
Theorem~\ref{thm:hamiltonentropy} gives
\(\log h(H)=N(\log d-1)+o(N)\).
Since \(H\) is contained in the support graph, (61) follows.
The inverse is the contrapositive of its positivity conclusion.
The \(o(N)\) term denotes a uniform vanishing exponential
loss under the displayed hypotheses, not a numerical
starting threshold.
\end{proof}

\section{Weighted Hamiltonian entropy and nonuniform graph counts}
\label{sec:weightedentropy}
The near-regular counting theorem is not the limit of the method.
Here the row observables in Lemma~\ref{lem:observablerounding}
let us retain the endpoint distribution for nonuniform kernels.
For a symmetric kernel \(P\), define
\[
 Z_{\rm H}(P)=
  \sum_{\substack{C\ {\rm unoriented}\\
                   {\rm Hamilton\ cycle\ on}\ [1,N]}}
                        \prod_{ij\in E(C)}P_{ij}.
\]
Zero entries simply give zero summands.

\begin{theorem}[Hamiltonian-weight universality]
\label{thm:weightedentropy}
Fix \(0<\delta<1\) and \(C>0\), and put \(K_{\rm w}=3\cdot10^7\).
Let \(P=P_N\) be symmetric stochastic bipartite kernels on
equal parts, with nontrivial gap at least \(\delta\).
If
\[
 \alpha=\max_{ij}P_{ij}\leq C(\log(2N))^{-K_{\rm w}},
\]
then, along even \(N\),
\[
 \log Z_{\rm H}(P)
       =-N+O_{\delta,C}\bigl(N\alpha^{1/8}\log(2N)\bigr).
\]
In particular the error is \(o(N)\).
There is no positive lower bound assumption on the nonzero entries.
The bound holds for all sufficiently large even \(N\);
its starting threshold may depend on the fixed \(C,\delta\).
\end{theorem}

\begin{lemma}[A diffuse stochastic permanent upper bound]
\label{lem:diffusepermanent}
If an \(N\)-by-\(N\) nonnegative matrix \(P\) has row sums 1
and maximum entry \(\alpha\to0\), then
\[
 \operatorname{per}(P)\leq\exp(-N+o(N)).
\]
\end{lemma}
\begin{proof}
Let \(k=\lceil1/\alpha\rceil\).
Each row of \(P/\alpha\) lies in \([0,1]^N\) and has sum
\(1/\alpha\leq k\leq N\).
Extend it coordinatewise to a vector of sum \(k\).
The polytope \(\{x\in[0,1]^N:\sum x_i=k\}\) has the
zero-one \(k\)-subsets as its vertices.
Since the permanent is linear in each row and is nonnegative,
one can replace the extended rows successively by such
vertices without decreasing the permanent.
Bregman's bound therefore gives
\[
 \operatorname{per}(P)
       \leq\alpha^N(k!)^{N/k}.
\]
Now \(\alpha k=1+O(\alpha)\), and Stirling yields
\[
 \log\operatorname{per}(P)
 \leq-N+O\bigl(N\alpha\log(1/\alpha)\bigr)
 =-N+o(N).
\]
This argument does not require column sums.
\end{proof}

\begin{proof}[Proof of Theorem~\ref{thm:weightedentropy}]
The upper bound follows from
\(2Z_{\rm H}(P)\leq\operatorname{per}(P)\) and the quantitative
proof of Lemma~\ref{lem:diffusepermanent}:
\(N\alpha\log(1/\alpha)\ll N\alpha^{1/8}\log(2N)\).

For the lower bound first remove a potential tiny-edge
obstruction without assuming it away.
Delete every positive entry smaller than \(N^{-10}\),
symmetrically. The removed mass at any row is at most \(N^{-9}\).
Corollary~\ref{cor:forbiddenrepair} gives
\[
 P'=\diag(u)P_{\rm trimmed}\diag(u),\qquad
 \|u-\one\|_\infty=O_\delta(N^{-9}\log N).
\]
It has fixed nontrivial gap at least \(\delta/2\), maximum
entry at most \(2\alpha\), and every nonzero entry at least
\(N^{-11}\), for all large \(N\).
Every Hamilton cycle of its support satisfies
\[
 \prod_{ij\in E(C)}P'_{ij}
       =\left(\prod_vu_v^2\right)
                       \prod_{ij\in E(C)}P_{ij}.
\]
The common multiplier has logarithm
\(O_\delta(N^{-8}\log N)=o(1)\).
Thus a lower bound for \(Z_{\rm H}(P')\) transfers to the
original \(P\), including its actual weights.
Rename the repaired kernel \(P\), its maximum entry \(\alpha\),
and its fixed gap for the rest of the lower-bound proof.
We may now use the proved lower weight \(N^{-11}\);
the original theorem did not assume it.

Let \(N=2m\), \(L=\log(2N)\), and choose
\[
 t=\lfloor\alpha^{-1/4}\rfloor,\qquad
 r_0=\lceil m\alpha^{1/8}\rceil,\qquad
 b=\left\lfloor\frac{m-r_0}{t}\right\rfloor,\qquad r=m-tb.
\]
Stochasticity gives \(\alpha\geq1/m\), and our upper bound
on \(\alpha\) ensures
\[
 \begin{gathered}
 t\leq m^{1/4},\qquad b\asymp N\alpha^{1/4},\qquad
 r=m\alpha^{1/8}+O(t+1),\\
 (r+b)L+tL=O(N\alpha^{1/8}L)=o(N).
 \end{gathered}
\]
This smaller reservoir supplies the quantitative error.
Partition the parts into
\(A_1,\ldots,A_t,R_A\) and \(B_1,\ldots,B_t,R_B\),
with the same layer sizes as in the unweighted proof.
Put \(p=b/m\asymp\alpha^{1/4}\).
Stirling's formula, including the reservoir and all factorial
remainders, gives
\[
 \log{\cal P}
   =2tb\log(1/p)+O(N\alpha^{1/8}L).
\]

We record precisely the weighted random-block estimates.
The cross block \(T\) of \(P\) is doubly stochastic, and
\(\|T-J_m/m\|\leq1-\delta\), with the renamed fixed gap.
Every row and column has squared norm at most \(\alpha\);
the centered matrix therefore has row and column norms
\(O(\sqrt\alpha)\) and entries \(O(\alpha)\).
Lemma~\ref{lem:randomrestriction}, with weighted scalar
sampling applied to the rows, gives for a block of common
sampling proportion \(z\)
\[
 \left\|P[X,Y]/z-J_s/s\right\|
       \leq1-\delta+O(\tau_z),\qquad
 \|z^{-1}P[X,Y]\one-\one\|_\infty=O(\tau_z),
\]
where
\[
 \tau_z=\sqrt{\alpha L/z}+\alpha L/z,\qquad s=zm.
\]
The degree estimate holds on both parts.
Its scalar variance uses \(\sum_jP_{ij}^2\leq\alpha\),
rather than an independence assumption on arithmetic entries.
The estimates hold with any fixed polynomially high probability.

For all consecutive layer pairs take \(z=p\).
Since \(\tau_p L^2\to0\), the normalized bipartite operator
has a fixed gap and Lemma~\ref{lem:scaling} applies.
Its multipliers satisfy \(u_v=1+O_\delta(\tau_pL)\).
If \(W=P[X,Y]\) is a weighted cross block, its scaled
doubly stochastic matrix is
\(\widetilde W_{ij}=u_iW_{ij}u_j/p\).
Consequently, with every multiplier charged,
\[
 \operatorname{per}(W)
  =\frac{p^b}{\prod_{v\in X\cup Y}u_v}
                      \operatorname{per}(\widetilde W)
  \geq (p/e)^b
                  \exp\bigl(-O_\delta(b\tau_pL)\bigr).
\]
The error is uniform across the \(2t-1\) layer pairs and
its total is \(O_\delta(N\tau_pL)
=O_\delta(N\alpha^{1/8}L)\).
The product of these weighted matching sums is exactly
the total weight of their long-path covers.

The augmented reservoir uses
\[
 K=R_A\cup A_1\cup R_B\cup B_t,\quad
 n_K=2(r+b),\quad p_K=(r+b)/m\asymp\alpha^{1/8}.
\]
Its raw kernel is \(P[K,K]/p_K\).
The same estimates, now at \(z=p_K\), and exact scaling
produce a stochastic kernel \(\widehat P\) on \(K\), with
fixed positive gap, maximum entry
\(\beta\leq C_\delta\alpha^{7/8}\), and multipliers
\(1+O_\delta(\tau_{p_K}L)\).

We must retain the endpoint row masses, not just the spectrum.
For the two endpoint sets \(S=A_1,B_t\), on each opposite
part, weighted scalar concentration gives
\[
 \left|\sum_{w\in S}P_{vw}-p\right|
       \leq C\bigl(\sqrt{p\alpha L}+\alpha L\bigr).
\]
Divide by \(p_K\) and include the scaling multipliers.
Writing \(\rho=b/(r+b)=p/p_K\), the result is
\[
 |(\widehat P\one_S)(v)-\rho|
 \leq C_\delta\left(
       \frac{\sqrt{p\alpha L}+\alpha L}{p_K}
                               +\tau_{p_K}L\right)
 \leq C_\delta\alpha^{7/16}L^2.
\]
This common majorant includes the sampling fluctuation
and the actual scaling change.
A union bound makes all layer and reservoir estimates
hold for a proportion \(1-o(1)\) of partitions.

Round \(\widehat P\), preserving these two row observables,
with
\[
 d=\lfloor\alpha^{-2/3}L^{-3}\rfloor.
\]
Here \(d\leq n_K/2\) for all large \(N\), since
\(\alpha^{-1}\leq m\) and \(n_K\asymp N\alpha^{1/8}\).
Also \(n_K\gg N^{7/8}\), so \(\log n_K\asymp L\).
The finite rounding condition holds:
\[
 d\beta(\log(2n_K))^2
           \ll_\delta\alpha^{5/24}/L\longrightarrow0.
\]
The resulting graph \(H_K\) is simple and exactly regular,
has fixed positive nontrivial gap, and has simultaneous
operator and row-observable error at most
\[
 C_\delta\left(d\beta L+\sqrt{L/d}+L/d\right)
                   \leq C_\delta\alpha^{5/24}L^2.
\]
Indeed the drift term is \(O_\delta(\alpha^{5/24}/L^2)\)
and the centered fluctuation is \(O(\alpha^{1/3}L^2)\);
both are included.
Thus its endpoint sets are bi-uniform, as sets and multisets,
with source parameter
\[
 \alpha_{\rm src}=C_\delta\alpha^{5/24}L^2.
\]
This is where the observable-preserving rounding, including
its predictable drift, is used essentially.

Verify every remaining condition of (42).
Take \(c=1/5\), \(d'=0\), and
\(\gamma=(\log n_K)^{-3}\).
Then
\[
 \alpha_{\rm src}d^c
 \ll_\delta\alpha^{3/40}L^{7/5}\to0,
\]
so \(\alpha_{\rm src}\leq d^{-c}\).
For the full number \(b\) of endpoint pairs,
\[
 \frac{b}{d^{-c/900}n_K}
 \ll\alpha^{1/8-1/6750}\longrightarrow0.
\]
Finally
\[
 d\gg L^{2K_{\rm w}/3-3}
        =L^{20,000,000-3}
        \gg(\log n_K)^{15,000,000}
        =\gamma^{-10^6/c}.
\]
The fixed gap supplies expansion larger than \(\gamma\)
eventually, the two endpoint sets have equal size and
distinct elements, and \(0<c<1/4\).
The source size and degree-error conditions hold as well.
Therefore every long-path cover closes to a Hamilton cycle
using only \(R_A\cup R_B\) for its new interior vertices,
in precisely the cyclic pairing used in the unweighted proof.

It remains to charge weighted completion and multiplicity.
Each completion uses exactly \(2r+b\) new edges.
Every such edge is a positive entry of the trimmed original
kernel and has weight at least \(N^{-11}\).
Thus an output cycle has weight at least its long-path
cover weight times \(N^{-11(2r+b)}\).
Its logarithmic cost is
\(O((r+b)L)=O(N\alpha^{1/8}L)\).
A fixed output cycle has at most
\[
 \binom N{2r+b}
       \leq\exp\bigl(O(N\alpha^{1/8}L)\bigr)
\]
preimage configurations, by the same edge-subset
reconstruction as (57).
All good-partition events hold with probability tending to
one, hence at least \(1/2\) for sufficiently large \(N\).
Summing the weighted matching bounds over those partitions
therefore gives, retaining every loss,
\[
 \begin{split}
 \log Z_{\rm H}(P)
 &\geq
 2tb\log(1/p)+(2tb-b)(\log p-1)
                          -O_\delta(N\alpha^{1/8}L)\\
 &=-N-O_\delta(N\alpha^{1/8}L).
 \end{split}
\]
For the last equality the residual terms are \(2r+b\) and
\(b\log(1/p)\), each of the stated order.
The lost matching edges, all closing weights, all scaling
multipliers, and the complete preimage multiplicity have
each been included in this inequality.
The initial trimming repair changes \(\log Z_{\rm H}\) by
only \(o(1)\), and its new maximum entry is at most twice the
original one. The same error bound therefore holds for the
original kernel. Together with the upper bound this proves
the quantitative theorem.
\end{proof}

\begin{corollary}[The entropy cost of the one-cycle constraint]
\label{cor:onecyclecost}
Under Theorem~\ref{thm:weightedentropy}, give a permutation
\(\sigma\) probability proportional to
\(\prod_{i=1}^N P_{i,\sigma(i)}\).
Then
\[
 0\leq-\log\mathbb P\{\sigma\hbox{ is one cycle}\}
       \leq C_{\delta,C}N\alpha^{1/8}\log(2N).
\]
Equivalently, with the Hamiltonian permanent counting both
orientations,
\[
 \log\frac{\operatorname{per}(P)}{2Z_{\rm H}(P)}
          =O_{\delta,C}(N\alpha^{1/8}\log(2N)).
\]
Thus enforcing a single global cycle costs only sublinear
entropy, although the support can have zeros.
The same comparison holds between the full adjacency
permanent and \(2h(G)\) in the graph-scaling corollary below.
\end{corollary}
\begin{proof}
The exact probability is \(2Z_{\rm H}(P)/\operatorname{per}(P)\).
It is at most one. The permanent upper bound and the weighted
Hamiltonian lower bound give the other inequality;
all smaller errors and the factor two are absorbed.
If \(P=\diag(u)A\diag(u)\), every permutation product and
every directed Hamilton-cycle product is multiplied by
the same factor \(\prod_vu_v^2\), so their ratio is unchanged.
\end{proof}

\begin{corollary}[Sinkhorn entropy determines nonuniform graph counts]
\label{cor:sinkhornentropy}
Let \(G_N\) be simple balanced bipartite graphs with
adjacency matrices \(A_N\).
Suppose positive vectors \(u\) give symmetric stochastic
scalings
\[
 P_N=\diag(u)A_N\diag(u)
\]
with a fixed positive nontrivial gap and maximum entry
\(O((\log N)^{-K_{\rm w}})\).
Then, with \(\alpha=\max_{ij}(P_N)_{ij}\),
\[
 \log h(G_N)=-N-2\sum_{v=1}^N\log u_v
                  +O_{\delta,C}(N\alpha^{1/8}\log(2N)).
\]
Equivalently, if
\[
 {\cal H}(P)=-\sum_{ij\in E(G_N)}P_{ij}\log P_{ij},
\]
with each undirected edge counted once, then
\[
 \log h(G_N)=2{\cal H}(P_N)-N+o(N).
\]
The displayed edge entropy is maximal among all nonnegative
proper edge weightings of \(G_N\) with row sums 1.
\end{corollary}
\begin{proof}
Every Hamilton cycle has degree two at each vertex, so
\[
 \prod_{ij\in E(C)}P_{ij}=\prod_vu_v^2
\]
independently of \(C\).
The weighted theorem gives the first formula.
Since the edge weights have row sums 1,
\[
 \sum_{ij\in E(G_N)}P_{ij}\log P_{ij}
 =\sum_{ij\in E(G_N)}P_{ij}(\log u_i+\log u_j)
 =\sum_v\log u_v.
\]
This proves the entropy formula and fixes its factor of two.

For any other proper edge weighting \(x\), its total mass
and that of \(P\) are both \(N/2\), and
\(\sum_e x_e\log P_e=\sum_v\log u_v\).
The log-sum inequality gives
\(\sum_e x_e\log(x_e/P_e)\geq0\).
Thus \({\cal H}(x)\leq{\cal H}(P)\), including zero \(x_e\)
by continuity.
\end{proof}

\begin{corollary}[A Markov-chain formulation]
\label{cor:markovhamiltonentropy}
Under Theorem~\ref{thm:weightedentropy}, start a Markov chain
with transition matrix \(P_N\) in its uniform stationary law.
Then
\[
 {\mathbb P}\{X_0,\ldots,X_{N-1}\ {\rm all\ distinct},
                        \ X_N=X_0\}=\exp(-N+o(N)).
\]
\end{corollary}
\begin{proof}
Each unoriented Hamilton cycle has exactly \(2N\)
choices of starting point and direction.
The stationary starting probability is \(1/N\).
The displayed probability is therefore exactly
\(2Z_{\rm H}(P_N)\); the quantitative logarithmic error is
\(O_{\delta,C}(N\alpha^{1/8}\log(2N))\) as well.
\end{proof}

\begin{remark}[Comparison with prior entropy methods]
Cuckler--Kahn~\cite{CK} establish graph-entropy counting in
the Dirac setting, and also describe dense balanced bipartite
analogues with minimum degree exceeding half a part.
The present hypothesis permits sparse zero support and
does not impose a linear minimum degree.
Barvinok's permanent/Hamiltonian-permanent comparison
\cite[Theorem 1.2]{Barvinok} assumes all entries are positive
and bounded below relative to their maximum.
Here zero entries are allowed, and the initial trimming
argument removes any separate positive-entry lower bound.
The fixed spectral gap, quantified diffuseness, and the
recent spanning-path input replace those dense-positivity
hypotheses. No novelty is claimed for entropy maximization
or Sinkhorn scaling themselves.
\end{remark}

\section{Weighted spanning paths with prescribed endpoints}
\label{sec:pathentropy}
For distinct vertices \(a,b\), let
\[
 Z_{\rm path}(P;a,b)
  =\sum_{\substack{Q\ {\rm Hamilton\ path}\\
                       Q\ {\rm oriented\ from}\ a\ {\rm to}\ b}}
               \prod_{\substack{(i,j)\ {\rm successive}\\
                                 {\rm in\ oriented}\ Q}}P_{ij}.
\]
A path is counted once with this prescribed orientation.

\begin{lemma}[A row-cap upper bound for path minors]
\label{lem:pathminor}
Let \(P\) be an \(N\)-by-\(N\) nonnegative row-stochastic
matrix with largest entry \(\alpha\to0\).
For every two distinct vertices,
\[
 Z_{\rm path}(P;a,b)
 \leq\operatorname{per}P^{(b,a)}
 \leq\exp\bigl(-N+
       O(N\alpha\log(1/\alpha)+\log N)\bigr),
\]
where \(P^{(b,a)}\) removes row \(b\) and column \(a\).
The constant is absolute, uniformly over the endpoints.
\end{lemma}
\begin{proof}
A directed Hamilton path maps every vertex other than \(b\)
to its unique successor, a bijection onto all vertices
other than \(a\). It is therefore one term of the minor
permanent, with exactly its original weight.

Put \(n=N-1\) and \(k=\min(n,\lceil1/\alpha\rceil)\).
Each row of the minor divided by \(\alpha\) belongs to
\([0,1]^n\), with sum at most both \(n\) and \(1/\alpha\),
hence at most \(k\).
Extend the row to sum \(k\).
The row-polytope and Bregman argument in
Lemma~\ref{lem:diffusepermanent} gives
\[
 \operatorname{per}P^{(b,a)}
               \leq\alpha^n(k!)^{n/k}.
\]
Stochasticity of the original matrix gives \(\alpha\geq1/N\).
Thus \(k\geq1/\alpha-1\), while \(\alpha k\leq1+\alpha\).
In particular \(k\asymp1/\alpha\).
Stirling gives the claimed bound, with the change from
\(n\) to \(N\) absorbed into \(O(\log N)\).
This proof also covers a minor with zero permanent.
It does not assume that deleting a column leaves row sum one.
\end{proof}

\paragraph{Verified path-minor consumer.}
The maintained module
\texttt{HamiltonianPrescribedPathMinor} now defines the actual prescribed
path sum by replacing column \(a\) by the unit vector at \(b\) and taking
the one-orbit cycle sum. The artificial closing edge \(b\to a\) has weight
one. Its permutation expansion retains every original path weight.
An exact Laplace identity gives the minor obtained by deleting row \(b\)
and column \(a\), with both endpoint labels explicitly transported.
The zero-permanent case is included.

For \(N\geq2\), nonnegative row-substochastic \(P\) with entries at most
\(\alpha\) and \(N\alpha\geq1\), the compiled endpoint-uniform consequence is
\[
 Z_{\rm path}(P;a,b)
       \leq \exp\bigl(-(N-1)+16N\alpha\log(2N)\bigr).
\]
The proof applies the existing row-cap estimate to the actual minor;
it never assumes that its rows retain sum one.
This upper bound suffices for the upper half of the weighted path-entropy
theorem, since \(\alpha\leq\alpha^{1/8}\) for \(0<\alpha\leq1\).
All nine new declarations passed a focused foundational-axiom audit.
The sharper Bregman error in Lemma~\ref{lem:pathminor} and the spanning-path
lower bound remain separate formal obligations.

The maintained module \texttt{HamiltonianPathContraction} additionally
proves a lower-bound consumer. Label the initial vertex \(0\), let
\(b\ne0\), and index the remaining \(N-1\) labels by \(i+1\). Put
\[
 Q_{ij}=P_{\,(0\ b)(i+1),\,j+1},
\]
where \((0\ b)\) is the indicated transposition.
Then
\[
 Z_{\rm cycle}(Q)\leq Z_{\rm path}(P;0,b)
                         \leq\operatorname{per}Q.
\]
The lower inequality is proved by lifting every one-orbit minor
permutation through the exact permutation decomposition. The proof
retains fixed-point orbits, all vertex labels and every edge weight;
the hypothesis \(b\ne0\) is explicit. Five new declarations passed
their focused foundational-axiom audit. A quantitative lower estimate
on this actual minor is still required before claiming path entropy.

\paragraph{Verified marked-edge suppression.}
The maintained consumer \texttt{HamiltonianPermanentRowMass}
retains the full row-normalization product. If a nonnegative
\(m\)-by-\(m\) matrix \(M\) has row sums \(s_i\geq r>0\),
\(r\leq1\), entries at most \(a/X\), \(a\geq1\), \(X\geq\max(1,m)\), then
\[
 \operatorname{per}M
 \leq\exp\!\left(-m-\sum_i(1-s_i)
                   +16(a/r)\log(2X)\right).
\]
Indeed \(M_{ij}/s_i\) is row-stochastic with entry cap
\(a/(rX)\). The permanent scales by \(\prod_i s_i\), and
\(\prod_i s_i\leq\exp(-\sum_i(1-s_i))\).
Every row deficit is retained.

Let \(E\) be any set of directed marked edges and suppose the
original row sums lie in \([1/2,1]\). Define
\[
 \mu_E=\sum_i\sum_j {\bf1}_{E}(i,j)M_{ij}.
\]
Halve precisely the marked entries. Every resulting row has sum
at least \(1/4\), and its deficit is at least half its original
marked mass. Every cycle avoiding all marked edges retains
exactly its original weight. Consequently the compiled bound is
\[
 Z_{\rm avoid}
 \leq\exp\!\left(-m-\mu_E/2+64a\log(2X)\right).
\]
The proof gives the exact partition and signed lower bound
\[
 Z_{\rm hit}=Z_{\rm cycle}-Z_{\rm avoid}
 \geq Z_{\rm cycle}
       -\exp\!\left(-m-\mu_E/2+64a\log(2X)\right).
\]
Each marked cycle is counted once, regardless of how many marked
edges it contains. No independence of edge events is assumed.
Positive \(Z_{\rm hit}\) yields an actual one-orbit permutation
of positive weight containing a marked edge.
All nine declarations passed their focused foundational-axiom audit.
The arithmetic marked-mass bound for coprime kernels is now compiled
in \texttt{HamiltonianCoprimeMarkedMass}, including both endpoint
losses and their positive intersection correction.

\paragraph{Verified injective weighted cutting.}
The maintained module \texttt{HamiltonianMarkedPathTransfer} gives
a general cycle-to-path count transfer. Label the endpoints \(0,1\),
let \(Q\) be the actual principal matrix on the other \(m\) vertices,
and let \(E\) be any set of marked entries of \(Q\).
For every real \(\kappa\), if \(M\) is nonnegative and
\[
 \kappa M_{ij}\leq M_{i1}M_{0j}
       \qquad\text{for each marked interior entry }(i,j),
\]
then the compiled comparison is
\[
 \kappa Z_{\rm marked}(Q,E)\leq Z_{\rm path}(M;0,1).
\]
Choose one usable edge of each marked source cycle.
Insert the terminal endpoint at that cut and then insert the
initial endpoint through the artificial closing edge.
Both nonzero insertions preserve the full single orbit.
Applying the exact permutation decomposition recovers the source
cycle, proving injectivity even when the chosen edge depends on it.
All uncut edge weights remain unchanged; the displayed local
inequality accounts for the replaced edge.
All nine new declarations passed their focused foundational-axiom audit.
In the uniform coprime case, the local condition holds with
\(\kappa=1/\varphi(n)\): both new endpoint entries equal
\(1/\varphi(n)\), and every removed entry is at most that value.
The maintained module \texttt{HamiltonianEndpointPathTransfer}
now transports this comparison to arbitrary distinct endpoints on
an arbitrary finite support. The full equivalence sends \(0,1\)
to the endpoints, and its tail is an equivalence onto their exact
principal complement. Both the marked cycle mass and prescribed
path mass are transported with their original weights.

\paragraph{Verified native signed path consumer and uniform budgets.}
The module \texttt{HamiltonianCoprimePathSigned} instantiates the
closing condition with \(1/d\), \(d=\varphi(n)\), on the actual
coprime endpoint-deleted support. Writing \(m=n-2\), \(A=n/d\),
and \(Z_{\rm base}=Z_{\rm cycle}(P_n|_{V'\times V'})\), the mathematical
form of its compiled native inequality, after simplifying the
exponent, is
\[
 Z_{\rm path}(P_n;a,b)\geq
 \frac1d\left[
 Z_{\rm base}-
 \exp\!\left(-m-\frac{d^2}{4n}+1+64A\log(2n)\right)\right].
\]
It holds for distinct compatible endpoints and \(d\geq4\).
The positive base count and the entire avoidance allowance remain
coupled. The overlap correction and every original edge weight
have already been retained in the preceding inputs.

The module \texttt{HamiltonianCoprimePathBudget} proves that for
every \(\beta,\gamma>0\), one \(N_0\geq512\) covers all \(n\geq N_0\)
with
\[
 \begin{gathered}
 d\geq4,\qquad
 4A^2\leq\beta\log(n-2),\qquad
 2A\leq\gamma\log\!\left\lfloor\frac{n-2}{2}\right\rfloor,\\
 \frac{16A^3}{n}\leq10^{-4},\\
 1+170A+64A\log(2n)+\log(8(n-2))
          \leq\frac{n}{4A^2}.
 \end{gathered}
\]
The proof uses the unconditional squared-sublogarithmic totient
bound with
\(\eta=\min(1,\beta/8,\gamma/4)>0\).
For \(L=\log(2n)\), the already proved log-power domination gives
\(160000L^3\leq n\) and \(1000L^4\leq n\) eventually;
\(\log n\leq2\log\lfloor(n-2)/2\rfloor\) preserves the rank shift.
These estimates discharge both repair-density ranges and the
complete half-mass allowance, uniformly over all growing moduli.
All nine new endpoint, native signed-consumer and scalar-budget
declarations passed their combined focused foundational-axiom audit.
The native odd/even lower-count specializations and their complete
positive signed path consumer are now compiled as follows.
The final actual integer path-count normalization and entropy assembly
remain formal obligations. The full general wide-cap spanning-path
entropy theorem remains a separate obligation.

\paragraph{Verified native positive path count.}
The modules \texttt{HamiltonianCoprimeOddBaseLower} and
\texttt{HamiltonianCoprimeEvenBaseLower} instantiate the independent
lower counts on the actual native principal and balanced supports.
The module \texttt{HamiltonianCoprimeInteriorLower} transports the
full retained cycle count and treats arbitrary compatible endpoint
labels, including both orderings of the even parity classes.
Uniformly over every sufficiently large modulus and such endpoints,
\[
 Z_{\rm base}\geq B_n:=\frac{\exp(-m-170A)}{4m}.
\]
The module \texttt{HamiltonianCoprimePathLower} consumes the entire
proved scalar budget without dropping any signed term. It proves
\[
 Z_{\rm avoid}\leq\frac{B_n}{2},\qquad
 Z_{\rm path}(P_n;a,b)\geq\frac{B_n}{2d}
             =\frac{\exp(-m-170A)}{8md}>0.
\]
All ten new native support, parity, uniform lower-count and final
signed-consumer declarations pass their combined focused
foundational-axiom audit. This certifies positive actual weighted
path mass with an explicit lower bound; the integer-count
normalization and logarithmic entropy comparison are the remaining
native endpoint-count consumers.

\paragraph{Verified independent retained-cycle supply.}
The modules \texttt{HamiltonianCycleExponentialLower},
\texttt{HamiltonianPrincipalDeletedCycleLower} and
\texttt{HamiltonianBalancedDeletedCycleLower} now give independent
positive lower counts on the actual retained supports.
Let \(T\) be a nonnegative stochastic \(n\)-by-\(n\) matrix,
with entries at most \(a/n\), \(a\geq1\), and
\(\|T-J/n\|\leq\rho\), where \(0\leq\rho<1\).
Delete the same \(r\)-element set from both coordinates,
put \(m=n-r\) and \(\Delta=(1-\rho)/2\), and assume
\[
 \frac r n\left(\frac a{1-\rho}\right)^3\leq10^{-4},
 \qquad
 (2a)^2\leq b(\Delta)\log m,
 \qquad
 b(\Delta)=\min\left(1,\frac1{18(16e^2+1+1/\Delta)}\right).
\]
For \(n\geq4\) and all sufficiently large \(m\), uniformly over
the matrix and deleted set, the compiled bound is
\[
 Z_{\rm cycle}(T|_{D^c\times D^c})
 \geq\frac{\exp\!\left(
       -m-((2a-1)+(2a-1)/\Delta)
          -40ra/(1-\rho)\right)}{4m}.
\]
Square repair initially identifies the two complements arbitrarily.
The proof undoes that row identification before counting cycles;
row permutation invariance of the permanent is not used as cycle
invariance. The resulting matrix on the actual principal complement
is stochastic, has cap \(2a/m\) and gap at least \(\Delta\).
The signed cycle/permanent comparison, positive permanent and
determinant lower bound supply its positive cycle mass.
The entire two-sided scaling product is then removed, with logarithmic
cost at most \(40ra/(1-\rho)\).

For balanced deletions \(S,U\) of size \(r\) from the two sides of
the actual bipartite lift, the corresponding compiled bound is
\[
 Z_{\rm cycle}((\operatorname{Lift}T)|_{(S\sqcup U)^c})
       \geq\frac{\exp(-2m-80ra/(1-\rho))}{2m}.
\]
Here the same repair-smallness condition holds, and its density
condition is \(2a\leq b_{\rm bip}(\rho')\log m\), where
\(\rho'=1-(1-\rho)/2\) and
\(b_{\rm bip}(\rho')=1/(1+|\kappa(\rho')|)>0\) is the proved
density from \texttt{HamiltonianSublogBipartiteCount};
\(\kappa\) is its existing Gram-amplification exponent.
This proof uses a full conjugation of the retained lift and retains
the square of the scaling product. It does not apply an aperiodic
gap theorem to a period-two matrix.
The nine new declarations passed their combined focused
foundational-axiom audit. The thresholds depend only on the fixed gap;
no numerical threshold is asserted. These explicit small-cap
consumers do not certify the general wide-cap path theorem below.




\begin{theorem}[Uniform weighted path entropy in the bipartite case]
\label{thm:bipartitepathentropy}
Under Theorem~\ref{thm:weightedentropy}, for every
\(a\in A,b\in B\),
\[
 \log Z_{\rm path}(P;a,b)
       =-N+O_{\delta,C}(N\alpha^{1/8}\log(2N)).
\]
The error and sufficiently-large-order threshold are uniform
over all opposite-part endpoints.
\end{theorem}
\begin{proof}
The upper bound follows from Lemma~\ref{lem:pathminor}.
We give the complete changes to the weighted lower-bound
construction.

First trim entries below \(N^{-10}\) and repair the kernel,
exactly as in Theorem~\ref{thm:weightedentropy}.
A spanning path from \(a\) to \(b\) changes its weight by
\[
 \frac{\prod_v u_v^2}{u_a u_b},
\]
because its endpoints have degree one and all other vertices
have degree two.
The logarithm of this factor is \(O_\delta(N^{-8}\log N)\).
After renaming the repaired kernel, all used positive
weights are at least \(N^{-11}\).

Use exactly the parameters
\[
 L=\log(2N),\quad
 t=\lfloor\alpha^{-1/4}\rfloor,\quad
 r_0=\lceil m\alpha^{1/8}\rceil,\quad
 b_0=\lfloor(m-r_0)/t\rfloor,\quad r=m-tb_0 .
\]
Here \(b_0\) is the layer size, distinct from the endpoint \(b\).
Require \(a\in R_A,b\in R_B\) in the ordered partitions.
The proportion of all partitions satisfying this requirement
is exactly \((r/m)^2\asymp\alpha^{1/4}\).
Its logarithmic cost is \(O(\log N)\), because
\(\alpha\geq1/m\).
All layer, reservoir, and endpoint-concentration failures
in the weighted proof have arbitrarily high polynomial
decay. Conditioning on this event costs at most a fixed
multiple of \(N^{1/4}\), so a proportion \(1-o(1)\) of the
conditioned partitions are still good, uniformly in \(a,b\).
This is a conditional counting argument, not an assumption
that fixed endpoints are randomly distributed.

Every good partition and choice of the \(2t-1\) matchings
gives \(b_0\) long paths, with endpoints
\(a_i\in A_1,z_i\in B_t\).
The weighted sum of all these covers is bounded below
by the same product of layer permanents as before.
Round the same augmented reservoir
\[
 K=R_A\cup A_1\cup R_B\cup B_t
\]
to the simple \(d\)-regular graph, where
\(d=\lfloor\alpha^{-2/3}L^{-3}\rfloor\).
The sets \(A_1,B_t\) have bi-uniformity error
\(O_\delta(\alpha^{5/24}L^2)\), with all scaling and
adaptive drift terms retained.

Use now the endpoint sets
\[
 A_1\cup\{a\},\qquad B_t\cup\{b\}.
\]
Adding one vertex changes normalized bi-uniformity error
by at most \(1/d+1/(r+b_0)\leq2/d\).
Consequently the source parameter is
\[
 \alpha_{\rm src}
    =O_\delta(\alpha^{5/24}L^2+d^{-1}).
\]
With \(c=1/5\),
\[
 \alpha_{\rm src}d^c
       =O_\delta(\alpha^{3/40}L^{7/5}+d^{-4/5})\to0 .
\]
The number of pairs is \(b_0+1\), instead of \(b_0\).
Since \(b_0\to\infty\), the already verified full pair-range
condition remains valid:
\[
 (b_0+1)d^{c/900}/|K|\to0.
\]
The degree bound, expansion, choice
\(\gamma=(\log|K|)^{-3}\), exact regularity, distinctness
and equal endpoint-set cardinalities are unchanged.
Thus every hypothesis of Brada\v{c}--Janzer
\cite[Theorem 6.4]{BJ} holds.

Apply it to the \(b_0+1\) opposite-part pairs
\[
 (a,z_1),\quad
 (a_1,z_2),\ldots,(a_{b_0-1},z_{b_0}),\quad(a_{b_0},b).
\]
Traverse the first connector, the first long path from
\(z_1\) to \(a_1\), then the next connector and so on.
The resulting object is a single spanning path from \(a\)
to \(b\). The source paths use every remaining reservoir
vertex as an interior vertex, while none of the prescribed
endpoints is an interior vertex.
All old long-path interiors lie outside \(K\);
for large \(N\), \(t\geq2\), so the two edge sets do not overlap.

A completion uses exactly \(2r+b_0-1\) new edges.
Their total weight is at least
\(N^{-11(2r+b_0-1)}\).
A fixed output path has at most
\[
 \binom{N-1}{2r+b_0-1}
       \leq\exp\bigl(O(N\alpha^{1/8}L)\bigr)
\]
preimages. Indeed its long-path subset determines all
isolated reservoir vertices, the unique orientation of
every long component from its first-part endpoint, and
hence every layer and matching.
Both the closing-weight cost and the full multiplicity
are \(O(N\alpha^{1/8}L)\) in logarithmic scale.
The original weighted partition and matching calculation,
with the additional \(O(\log N)\) endpoint-conditioning
cost, gives
\[
 \log Z_{\rm path}(P;a,b)
       \geq-N-O_{\delta,C}(N\alpha^{1/8}L).
\]
Transfer back through the initial endpoint-correct scaling.
Every constant used was independent of the chosen endpoints.
\end{proof}

\section{General symmetric kernels, complementary cuts, and odd orders}
\label{sec:generalkernels}
The bipartite hypothesis can be removed when the negative spectrum
is also separated from \(-1\).
The essential new transfer is a dependent, complementary-cut
estimate. Independent row and column selectors cannot simply be
conditioned to be complements: that event can have exponentially
small probability. The argument below keeps their dependence exactly.

\begin{lemma}[Complementary restriction with the full covariance retained]
\label{lem:complementaryrestriction}
Let \(N\geq2\) and let \(E=E^\top\) be an \(N\)-by-\(N\) matrix with
\(\|E\|\leq1\) and all squared column norms at most \(\alpha\geq0\).
Let \(R\) select each coordinate independently with probability
\(1/2\), and put \(R'=I-R\).
For every fixed \(A>0\), with probability \(1-N^{-A}\),
\[
 \|RER'\|^2\leq\tfrac14\|E\|^2+
       C_A\bigl(\sqrt{\alpha\log(2N)}+\alpha\log(2N)\bigr).
\]
The assertion remains true with polynomially high probability
after conditioning \(|R|=N/2\), when \(N\) is even, or after
also prescribing one coordinate in each part.
\end{lemma}
\begin{proof}
For columns \(c_j\) of \(E\), write
\[
 F=ERE-\tfrac12 E^2
       =\sum_j(\one_{j\in R}-1/2)c_jc_j^\top .
\]
The summands have norm at most \(\alpha\), and their matrix
variance is at most
\[
 \tfrac14\sum_j\|c_j\|_2^2c_jc_j^\top
              \preceq\tfrac{\alpha}{4}E^2 .
\]
Here is a finite signed-moment proof of the required tail.
Put \(A_j=c_jc_j^\top/2\), write
\(F_\epsilon=\sum_{j=1}^N\epsilon_jA_j\), and let
\(M_{2k}=2^{-N}\sum_{\epsilon\in\{-1,1\}^N}
                   \operatorname{tr}F_\epsilon^{2k}\).
For \(\epsilon^{(j)}\) obtained by flipping coordinate \(j\),
the exact identities are
\[
 \Delta_j=F_\epsilon-F_{\epsilon^{(j)}}=2\epsilon_jA_j,
 \qquad \Delta_j^2=4A_j^2,\qquad \sum_{j=1}^N\Delta_j=2F_\epsilon .
\]
The flip is an involution of the whole sign cube. Hence
\[
 4M_{2k}=2^{-N}\sum_{\epsilon\in\{-1,1\}^N}\sum_{j=1}^N
 \operatorname{tr}\!\left(
 \Delta_j(F_\epsilon^{2k-1}-F_{\epsilon^{(j)}}^{2k-1})\right).
\]
The trace mean-value inequality proved below retains both
powers and every mixed term. Apply it before the flip
bijection combines the two power contributions, giving
\[
 M_{2k}\leq(2k-1)\,2^{-N}\sum_{\epsilon\in\{-1,1\}^N}
                 \operatorname{tr}(V F_\epsilon^{2k-2}),
 \qquad V=\sum_{j=1}^NA_j^2 .
\]
In full, \(V=\tfrac14E\operatorname{diag}
(\|c_j\|_2^2)E^\top\). For every symmetric \(Z\) and
even \(q\), the diagonals of \(E^\top Z^qE\) are
nonnegative: for \(q=2r\) this is the Gram matrix
\((Z^rE)^\top(Z^rE)\).
Cyclic trace therefore gives
\[
 \operatorname{tr}(VZ^q)
 \leq\tfrac{\alpha}{4}\operatorname{tr}(E^2Z^q)
 \leq\tfrac{\alpha}{4}\operatorname{tr}(Z^q).
\]
For the last inequality, use the independent eigenbasis
overlap squares of \(E,Z\), the bounds
\(\lambda_i(E)^2\leq1\) and \(\lambda_j(Z)^q\geq0\).
Consequently
\[
 M_{2k}\leq N(2k-1)!!(\alpha/4)^k
          \leq N(\alpha k/2)^k .
\]
The even trace controls \(\|F_\epsilon\|^{2k}\), so, for
every \(k\geq1,t>0\),
\[
 2^{-N}\#\{\epsilon:\|F_\epsilon\|>t\}
          \leq\frac{N(\alpha k/2)^k}{t^{2k}}.
\]
For \(N\geq2,A>0,\alpha>0\), take
\(k=\lceil(A+4)\log(2N)\rceil\).
At \(t=\sqrt{2\alpha k}\) the bound is \(N4^{-k}\).
The inequalities \(\log4\geq1\) and
\(k\leq(A+5)\log(2N)\) give the explicit uniform form
\[
 \Pr\!\left\{\|F\|>
       \sqrt{2(A+5)\alpha\log(2N)}\right\}
             \leq N^{-A-3}.
\]
If \(\alpha=0\), every column vanishes and \(E=F=0\),
so the same conclusion holds. This proves the asserted
tail, with room for the conditioning costs below.
Crucially the other restriction has the exact relation
\[
 ER'E=\tfrac12 E^2-F .
\]
Consequently, retaining the same covariance \(F\),
\[
 \begin{split}
 \|RER'\|^2
 &=\|R'ERER'\|\\
 &\leq\tfrac12\|ER'\|^2+\|F\|\\
 &\leq\tfrac14\|E\|^2+\tfrac32\|F\|.
 \end{split}
\]
This proves the assertion without independent complementary
selectors or a discarded cross term.

The balanced event has probability at least \(1/(N+1)\).
With prescribed distinct coordinates \(a\in R,b\notin R\),
its probability is
\(2^{-N}\binom{N-2}{N/2-1}\geq1/(4(N-1))\).
Choosing the initial failure exponent larger absorbs these
polynomial conditioning costs.
\end{proof}

\paragraph{Maintained full complementary covariance and actual cut bound.}
The modules \texttt{HamiltonianComplementaryCovariance} and
\texttt{HamiltonianComplementaryCut} now prove the deterministic
part on the actual coordinate supports.
For any coordinate set \(D\), let \(R\) and \(S\) be its two
complementary diagonal projections, and put
\(F=ERE-\tfrac12E^2\).
The exact identities and bound are
\[
 ESE=\tfrac12E^2-F,\qquad F=F^\top,\qquad
 \|E[D,D^c]\|^2\leq\tfrac14\|E\|^2+\tfrac32\|F\|.
\]
The same \(F\) occurs with both signs before the operator estimates;
no independent-selector assumption is used.
The formal proof includes projection symmetry, idempotence,
complementarity and norm, the exact retained matrix, and its
dimension-preserving injective restriction.
Every column remains in the full signed covariance entry:
\[
 F_{ij}=\sum_{k=1}^N
       \bigl(\mathbf1_{k\in D}-\tfrac12\bigr)E_{ik}E_{kj}.
\]
It also proves the actual spectral consumer:
if \(0<\delta\leq1\), \(\|E\|\leq1-\delta\) and
\(\|F\|\leq\delta/128\), then
\[
 2\|E[D,D^c]\|\leq1-\tfrac34\delta .
\]
All eleven new deterministic covariance, coordinate-support and
gap declarations pass their combined focused foundational-axiom audit.
These statements are uniform over the finite matrix, its dimension,
and every cut; they do not restrict \(N\alpha\) to a sublogarithmic
range. The uniform probability tail for \(F\) is now formalized below.
Balanced and prescribed-coordinate conditioning in
Lemma~\ref{lem:complementaryrestriction} are now formalized below.
The full wide-cap entropy theorem is not yet claimed as formalized.

\paragraph{Maintained signed trace moment input.}
The three modules \texttt{HamiltonianTraceMeanValueScalar},
\texttt{HamiltonianMixedTraceReadout}, and
\texttt{HamiltonianTraceMeanValue} now prove the complete
trace mean-value inequality for all real symmetric matrices
\(X,Y\) and every integer \(k\geq1\):
\[
 \operatorname{tr}\!\left((X-Y)(X^{2k-1}-Y^{2k-1})\right)
 \leq \frac{2k-1}{2}
       \operatorname{tr}\!\left((X-Y)^2(X^{2k-2}+Y^{2k-2})\right).
\]
The scalar proof pairs every exponent of the exact difference
quotient with its reflection. Independent orthogonal
diagonalizations of \(X,Y\) then express both full trace
polynomials through the same nonnegative overlap squares
\((U^\top V)_{ij}^2\). Zero powers are included, so the pure
terms and every mixed compensation term retain exactly the
same finite support. No assumption on the signs of the
eigenvalues or on their multiplicities is imposed.
All seventeen new scalar, mixed-readout and trace declarations
compile and pass their focused foundational-axiom audit.

For the covariance above, put \(A_j=c_jc_j^\top/2\) and
\(F=\sum_{j=1}^N\epsilon_jA_j\), with independent uniform signs.
The compiled finite sign-flip identity yields the exact
variance calculation
\[
 V=\sum_j A_j^2
  =\frac14\sum_j\|c_j\|_2^2c_jc_j^\top
  \preceq \frac{\alpha}{4}E^2,
 \qquad
 \mathbb E\operatorname{tr} F^{2k}
 \leq (2k-1)\frac{\alpha}{4}\mathbb E\operatorname{tr} F^{2k-2}.
\]
These finite sign-flip moment and variance statements are now
Lean theorems, together with their uniform probability tail.
The eight maintained owners from \texttt{HamiltonianSignCube}
through \texttt{HamiltonianComplementaryPolynomialTail}
include the exact involutions, squared increments, full
variance factorization, all-order moments, actual norm
readout, the zero-variance case, and an exact bijection
between sign vectors and every coordinate cut.
All forty-eight new declarations compile and pass their
combined focused foundational-axiom audit.

In particular, their actual covariance tail is
\[
 \frac{\#\{D\subseteq[N]:
   \|ERE-\tfrac12E^2\|>
       \sqrt{2(A+5)\alpha\log(2N)}\}}{2^N}
          \leq N^{-A-3}.
\]
The same finite count bound holds for cuts violating the
actual quarter estimate
\[
 \|E[D,D^c]\|^2
 \leq\tfrac14\|E\|^2+
          \tfrac32\sqrt{2(A+5)\alpha\log(2N)}.
\]
The exact fixed-gap failure inclusion is also compiled.
These are uniform over every finite dimension \(N\geq2\)
and every \(\alpha\geq0\), with no small-\(N\alpha\)
restriction. Balanced and prescribed-coordinate conditioning
are also formalized below. The full wide-cap entropy
headline remains unfinished.

\paragraph{Maintained balanced and prescribed finite laws.}
The modules \texttt{HamiltonianBalancedCutCounts},
\texttt{HamiltonianPrescribedCutCounts},
\texttt{HamiltonianFiniteConditioning}, and
\texttt{HamiltonianConditionedComplementaryTail} complete
the conditioning step.
For \(N=2m\), \(m\geq1\), write
\[
 \mathcal B_m=\{\epsilon\in\{-1,1\}^N:|D_\epsilon|=m\},
 \quad
 \mathcal B_m(a,b)=\{\epsilon\in\mathcal B_m:
                           a\in D_\epsilon,\ b\notin D_\epsilon\}.
\]
For every distinct \(a,b\), their exact counts and budgets are
\[
 \begin{aligned}
 |\mathcal B_m|&=\binom Nm,&
 2^N&\leq(N+1)|\mathcal B_m|,\\
 |\mathcal B_m(a,b)|&=\binom{N-2}{m-1},&
 2^N&\leq4(N-1)|\mathcal B_m(a,b)|.
 \end{aligned}
\]
The two prescribed-coordinate count is a genuine
insertion/deletion bijection on the finite supports.
Both budgets follow from
\(2^{2r}\leq(2r+1)\binom{2r}{r}\), proved by summing
every binomial coefficient and comparing it to the middle one.
No endpoint case or coinciding-coordinate exception is suppressed;
the distinctness hypothesis is explicit, and \(N=2\) is included.

For any exceptional subset \(\mathcal E\) of the full sign cube
with \(|\mathcal E|/2^N\leq N^{-A-3}\), the exact
intersection counts give
\[
 \frac{|\mathcal E\cap\mathcal B_m|}{|\mathcal B_m|}
       \leq N^{-A},
 \qquad
 \frac{|\mathcal E\cap\mathcal B_m(a,b)|}{|\mathcal B_m(a,b)|}
       \leq N^{-A}.
\]
The denominators are proved positive from their count budgets.
The final factor absorption uses
\(N+1\leq N^3\) and \(4(N-1)\leq N^3\) for \(N\geq2\).
Applying these compiled statements to both the actual
covariance exceptional set and the actual coordinate
quarter-bound exceptional set proves the full finite
uniform-cut form of Lemma~\ref{lem:complementaryrestriction},
including both conditioning clauses.
Its stronger explicit error is
\[
 \tfrac32\sqrt{2(A+5)\alpha\log(2N)}.
\]
All eighteen new finite-count and conditional-tail declarations
compile and pass their focused foundational-axiom audit.
Weighted row concentration, the repaired balanced kernels and
the full general weighted entropy count remain separate
unfinished formal obligations.

\begin{lemma}[Scaling repair without a parity kernel]
\label{lem:aperiodicrepair}
Let \(Q\) be symmetric nonnegative of any order, and put
\[
 d=Q\one,\quad D=\diag(d),\quad
 \|d-\one\|_\infty\leq\varepsilon<1/2,\quad
 S=D^{-1/2}QD^{-1/2}.
\]
Assume \(S\) has simple eigenvalue \(1\) and all other
eigenvalues of modulus at most \(1-\delta\).
Then the scaling conclusion (5) holds. No bipartition is needed.
For sufficiently small \(\varepsilon L^2\), the repaired matrix
has nontrivial gap at least \(\delta/2\) and largest entry at most
twice that of \(Q\).
\end{lemma}
\begin{proof}
Use \(D,T,\pi\) as in Lemma~\ref{lem:scaling}, and let
\(\Pi_+f=\one\sum_i\pi_i f_i\), \(\Pi_0=I-\Pi_+\).
The same reversible spectral calculation gives
\[
 \|T^j\Pi_0\|_{\infty\to\infty}
       \leq C\min(1,\sqrt N(1-\delta)^j).
\]
There is no eigenvalue \(-1\), so \(D+Q=D(I+T)\) is invertible
on the whole space. Its inverse is specified by
\[
 (D+Q)^{-1}r
   =\tfrac12\Pi_+(D^{-1}r)
      +\sum_{j\geq0}(-1)^jT^j\Pi_0(D^{-1}r).
\]
Its infinity norm is \(O_\delta(L)\).
Equation (8) therefore has a contraction solution on the full
ball \(\|z\|_\infty\leq C_\delta\varepsilon L\);
the nonlinear term and its Lipschitz cost are exactly those
already calculated there. No orthogonality or gauge restriction
is imposed or required.
All entries of \(u=\one+z\) are positive.
The complete expansion \(ZQ+QZ+ZQZ\) gives the operator bound,
and the multiplier bound gives the entry bound.
Comparison with the normalized input and Weyl's inequality
give the stated gap after decreasing the admissible constant.
\end{proof}

\begin{lemma}[Most balanced cuts admit controlled stochastic scaling]
\label{lem:balancedcutscaling}
Let \(N\) be even and \(P=P^\top\) be nonnegative stochastic, with
\[
 \|P-J_N/N\|\leq1-\delta,\qquad \max_{ij}P_{ij}=\alpha,
 \qquad \alpha(\log(2N))^5\to0 .
\]
A proportion \(1-o(1)\) of all balanced ordered cuts \(X,Y\)
have the following property.
The symmetric bipartite kernel \(Q\) with cross block
\(2P[X,Y]\) has an exact stochastic scaling
\[
 P^{X,Y}=\diag(u)Q\diag(u),\qquad
 \|u-\one\|_\infty=O_\delta(\sqrt{\alpha}L^{3/2}).
\]
It has fixed nontrivial gap at least \(\delta/8\) and largest
entry at most \(4\alpha\).
The same conclusion holds among cuts with any fixed
\(a\in X,b\in Y\), uniformly in distinct \(a,b\).
\end{lemma}
\begin{proof}
Put \(E=P-J_N/N\).
Since \(P\) is symmetric stochastic,
\[
 \sum_j E_{ij}^2=\sum_jP_{ij}^2-1/N\leq\alpha .
\]
Lemma~\ref{lem:complementaryrestriction} thus applies.
For a balanced cut, exactly
\[
 2P[X,Y]-J_{N/2}/(N/2)=2E[X,Y].
\]
Its norm is at most
\(\sqrt{(1-\delta)^2+O(\sqrt{\alpha L}+\alpha L)}\),
hence at most \(1-3\delta/4\) for sufficiently large \(N\).
This bounds the second singular value of the cross block
by the rank-one approximation inequality.

Scalar concentration of each weighted row sum, with variance
at most \(\alpha/4\), gives
\[
 \|Q\one-\one\|_\infty
             \leq C(\sqrt{\alpha L}+\alpha L).
\]
Union bounds and the same polynomial conditioning costs give
this simultaneously with the operator estimate.
They also apply if \(P\) has nonzero diagonal entries:
every cross row sum is an actual selected sum, and each
diagonal contribution is at most \(\alpha\).
Normalizing the degrees changes the symmetric operator by
\(O(\sqrt{\alpha L}+\alpha L)\).
The normalized bipartite operator has its exact eigenvalues
\(1,-1\), and the remaining spectrum stays a fixed distance
from both; simplicity follows from the second-singular-value
bound.
Now \(\varepsilon L^2=O(\sqrt\alpha L^{5/2})=o(1)\).
Apply Lemma~\ref{lem:scaling}.
The multiplier, maximum-entry and gap conclusions follow,
with constants uniform over the good cuts and endpoints.
\end{proof}

\begin{theorem}[Hamiltonian-weight universality at every order]
\label{thm:generalweightedentropy}
Fix \(0<\delta<1,C>0\), and let \(K_{\rm w}=3\cdot10^7\).
Let \(P=P_N\) be any symmetric nonnegative stochastic kernels,
with
\[
 \|P-J_N/N\|\leq1-\delta,\qquad
 \alpha=\max_{ij}P_{ij}\leq C(\log(2N))^{-K_{\rm w}}.
\]
Loops are permitted in the kernel; Hamilton paths and cycles
use distinct vertices and no loops.
For all sufficiently large \(N\), of either parity,
\[
 \begin{split}
 \log Z_{\rm H}(P)
       &=-N+O_{\delta,C}(N\alpha^{1/8}\log(2N)),\\
 \log Z_{\rm path}(P;a,b)
       &=-N+O_{\delta,C}(N\alpha^{1/8}\log(2N))
                         \quad(a\ne b).
 \end{split}
\]
The endpoint error and threshold are uniform over all \(a,b\).
Together with the bipartite theorems, this covers both a fixed
aperiodic spectral gap and the exact balanced parity case.
It makes no assertion uniform in a spectral gap tending to zero.
\end{theorem}
\begin{proof}
The permanent upper bounds, including path minors, are already
Lemma~\ref{lem:diffusepermanent} and Lemma~\ref{lem:pathminor}.
Write \(R_N=N\alpha^{1/8}\log(2N)\).
Since \(\alpha\geq1/N\), this absorbs \(O(\log N)\).

Suppose first that \(N\) is even.
For every good balanced cut apply the bipartite weighted
theorem to \(P^{X,Y}\).
For a Hamilton cycle crossing that cut, the exact relation is
\[
 \prod_{ij\in E(C)}P^{X,Y}_{ij}
       =2^N\left(\prod_vu_v^2\right)
                         \prod_{ij\in E(C)}P_{ij}.
\]
The logarithm of the multiplier is
\(O_\delta(N\sqrt\alpha L^{3/2})=o(R_N)\).
The weighted theorem is uniform on these cuts, since their
gap is at least \(\delta/8\) and their largest entry at most
\(4\alpha\). Each good cut therefore contributes at least
\[
 2^{-N}\exp(-N-O_{\delta,C}(R_N))
\]
to the Hamilton sum in the original kernel.

Every even Hamilton cycle has exactly two balanced ordered
alternating cuts. Hence the exact double-counting identity is
\[
 2Z_{\rm H}(P)
       =\sum_{\substack{X\subseteq[1,N]\\|X|=N/2}}
                        Z_{\rm H}(P\hbox{ restricted to }X,X^c).
\]
There are \(\binom N{N/2}\) cuts and a proportion \(1-o(1)\)
are good. The factor \(2^{-N}\) cancels their exponential
number, leaving only \(O(\log N)\) in the logarithm.
This proves the cycle lower bound at even orders.

For a fixed oriented Hamilton path from \(a\) to \(b\),
there is exactly one balanced cut with \(a\in X,b\in Y\)
which it crosses throughout.
The exact scaling factor for its weight is instead
\[
 2^{N-1}\frac{\prod_vu_v^2}{u_a u_b}.
\]
Apply Theorem~\ref{thm:bipartitepathentropy} on each good cut
with these endpoints.
There are \(\binom{N-2}{N/2-1}\) such cuts, almost all good
uniformly in \(a,b\). Their exponential number cancels
\(2^{-(N-1)}\), again up to a polynomial factor.
This proves the even-order path lower bound, including all
endpoint and scaling factors.

Here is the reduction for odd orders.
Remove an arbitrary vertex \(v\), leaving the principal
matrix \(Q\) of order \(n=N-1\).
Its row defect is at most \(\alpha\).
Interlacing gives
\[
 \lambda_2(Q)\leq1-\delta,\qquad
 \lambda_{\min}(Q)\geq-1+\delta .
\]
Its normalized matrix differs from \(Q\) by \(O(\alpha)\),
has exact eigenvalue 1, and all other eigenvalues have modulus
at most \(1-\delta/2\), for large \(N\).
Lemma~\ref{lem:aperiodicrepair} therefore supplies
\[
 \widetilde P=\diag(u)Q\diag(u),\quad
 \|u-\one\|_\infty=O_\delta(\alpha L),\quad
 \max\widetilde P_{ij}\leq2\alpha
\]
with a fixed positive aperiodic gap.
The even-order path conclusion just proved applies to it.
Transferring a path back to \(Q\) uses the exact factor
\(\prod u_i^2/(u_a u_b)\), costing \(O_\delta(N\alpha L)=o(R_N)\).

Choose two distinct off-diagonal neighbours \(a,b\) of \(v\)
with \(P_{va},P_{vb}\geq1/(2N)\).
Such neighbours exist for large \(N\): after excluding the
diagonal and the largest off-diagonal entry, the remaining
row mass is at least \(1-2\alpha\).
Close a spanning \(a\)-to-\(b\) path of \(Q\) through \(v\).
This is an injective construction of Hamilton cycles, and
\[
 Z_{\rm H}(P)\geq
 P_{va}P_{vb}\,Z_{\rm path}(Q;a,b)
       \geq (4N^2)^{-1}\exp(-n-O_{\delta,C}(R_N)).
\]
Its additional logarithmic cost is \(O(\log N)\).
This proves the odd cycle bound without assuming a bipartition.

Finally, for prescribed endpoints \(a,b\) at odd order, remove
\(a\) and repair its principal matrix in the same way.
There is a neighbour \(c\notin\{a,b\}\) with
\(P_{ac}\geq1/(2N)\), since excluding the two specified entries
leaves mass at least \(1-2\alpha\).
Prepend \(a,c\) to a spanning path from \(c\) to \(b\) on the
remaining even-order matrix.
The construction is injective, its scaling cost is
\(O_\delta(N\alpha L)\), and its extra edge costs only
\(O(\log N)\).
This proves the final lower bound uniformly in the endpoints.
No parity case or diagonal mass has been omitted.
\end{proof}

\begin{corollary}[Sparse graph entropy without a bipartite restriction]
\label{cor:generalsinkhornentropy}
Let simple graphs \(G_N\), of any order, admit positive
symmetric stochastic scalings
\(P_N=\diag(u)A_{G_N}\diag(u)\).
Suppose either the aperiodic hypotheses of
Theorem~\ref{thm:generalweightedentropy}, or the balanced
bipartite hypotheses of Theorem~\ref{thm:weightedentropy}, hold.
Then
\[
 \log h(G_N)
   =2\max_{\substack{x_e\geq0\\\sum_{e\ni v}x_e=1}}
                \sum_{e\in E(G_N)}x_e\log(1/x_e)
      -N+O_{\delta,C}(N\alpha^{1/8}\log(2N)).
\]
Moreover the logarithmic cost of imposing the one-cycle
condition on weighted permutations is at most the same
order, for all orders covered here.
\end{corollary}
\begin{proof}
The cycle-weight multiplier is always \(\prod_vu_v^2\).
The entropy-maximization calculation in
Corollary~\ref{cor:sinkhornentropy} uses only the row equations,
so applies to general graphs as well.
The full permanent and Hamiltonian permanent change by the
same multiplier. Their comparison follows from the two
weighted upper and lower estimates just proved.
\end{proof}

\section{A growing-modulus arithmetic application: coprime circles}
\label{sec:coprimecircles}
This application uses exact finite Fourier analysis, with no
prime number theorem, fixed-modulus assumption or prime-tuple input.
Let \(X_n\) have vertex set \(\mathbb Z/n\mathbb Z\), with
distinct vertices adjacent when their difference is coprime to \(n\).
These are the unitary Cayley graphs.
Their Hamiltonicity and stronger path properties are known;
see Klotz--Sander~\cite{KS} and Ba\v{s}i\'c~\cite{Basic}.
The conclusion here concerns the full exponential count.

\begin{theorem}[Coprime-circle entropy at every modulus]
\label{thm:coprimecircles}
As \(n\to\infty\) through all integers,
\[
 \log h(X_n)
   =n(\log\varphi(n)-1)
       +O\left(n\varphi(n)^{-1/8}\log(2n)\right).
\]
In particular \(h(X_n)=(\varphi(n)/e)^n\exp(o(n))\).
The constant is absolute.
For odd \(n\), the same logarithmic estimate holds for the
number of spanning paths joining any prescribed distinct endpoints.
For even \(n\), it holds for any prescribed opposite-parity endpoints.
\end{theorem}
\begin{proof}
The degree is \(d=\varphi(n)\), and \(P=A_{X_n}/d\) is symmetric
stochastic with maximum entry \(1/d\).
The character \(\e(kx/n)\) has eigenvalue
\[
 \lambda_k(A_{X_n})
   =\sum_{\substack{a\bmod n\\(a,n)=1}}\e(ka/n)
   =\varphi(n)\frac{\mu(t)}{\varphi(t)},
       \qquad t=n/(n,k).
\]
This is the standard Ramanujan-sum spectrum; we retain its sign.
The only \(t\) with \(\varphi(t)=1\) are \(1,2\).
The eigenvalue at \(t=1\) is the simple constant eigenvalue.
When \(n\) is even, \(t=2\) supplies the unique parity
eigenvalue \(-d\).
Every remaining eigenvalue has modulus at most \(d/2\).
For odd \(n\) there is no \(t=2\), so
\[
 \|P-J_n/n\|\leq1/2.
\]
For even \(n\), all units are odd and the parity parts are equal;
the same bound holds on the complement of the constant
and parity vectors.
These are exact statements for each modulus, not estimates
uniformly transferred from a fixed progression.

We also verify diffuseness without invoking an ineffective
totient estimate.
For a prime-power factor \(p^a\), its factor in
\(\varphi(n)^2/n\) is \(p^{a-2}(p-1)^2\).
It is at least 1 for \(p\geq3\), and the sole possible factor
below 1 is \(1/2\) at \(p=2,a=1\).
Thus \(\varphi(n)\geq\sqrt{n/2}\).
In particular \(1/d\) satisfies every fixed inverse-logarithmic
diffuseness threshold for sufficiently large \(n\).

Use Theorem~\ref{thm:generalweightedentropy} for odd \(n\), and
Theorems~\ref{thm:weightedentropy} and
\ref{thm:bipartitepathentropy} for even \(n\), with a fixed
gap such as \(\delta=1/3\).
Every Hamilton cycle has weight \(d^{-n}\).
Every spanning path has weight \(d^{-(n-1)}\).
The cycle assertion follows directly.
For paths the difference of one \(\log d=O(\log n)\) is
absorbed by the displayed error, uniformly over endpoints.
The totient lower bound makes the error \(o(n)\).
All constants may be chosen absolute because the spectral
gap was uniform over every modulus.
\end{proof}

\paragraph{Maintained spectral input and count conclusions.}
The exact coprime-difference graph, without a prime-modulus or
squarefree-modulus restriction, is formalized in
\texttt{HamiltonianCoprimeFourier}.
Its normalized kernel is nonnegative, symmetric and doubly
stochastic, with maximum entry at most \(1/\varphi(n)\).
\texttt{HamiltonianCoprimeConductor} proves the exact signed
unit-group reduction before any estimate is applied.
\texttt{HamiltonianCoprimeSpectrum} then proves, for every
residue frequency \(k\), the normalized eigenvalue
\[
 \frac{c_n(k)}{\varphi(n)}
    =\frac{\mu(t)}{\varphi(t)},\qquad
 t=\frac n{\gcd(n,k)},
\]
and exactly \(\varphi(t)\) frequencies of each conductor
\(t\mid n\). Conductor one is exactly the constant mode;
conductor two has normalized value \(-1\) and can occur only
at even modulus. Conductors at most two are exactly the
frequencies annihilated by multiplication by two.
Every remaining mode has modulus at most \(1/2\).
\texttt{HamiltonianCoprimeSymbol} proves that the full signed
symbol inverts to the actual physical adjacency kernel.
Using the previously proved real-to-complex norm transfer,
\texttt{HamiltonianCoprimeGap} proves the actual centered
real operator bound \(\|P-J_n/n\|\leq1/2\) for every odd
modulus. Thus the odd spectral hypothesis is discharged,
rather than supplied as an additional assumption.
\texttt{HamiltonianCoprimeEvenGap} also proves the native even
physical half-gap after removing precisely the constant and
parity modes. It identifies the unique conductor-two frequency
with \(n/2\) and its exact inverse-transform parity kernel.
No nonsquarefree modulus or high mode is omitted.
\texttt{HamiltonianCoprimeCrossBlock} now identifies the native
even kernel with its exact bipartite lift and proves stochasticity,
the entry cap and the physical centered half-gap on that cross block.
\texttt{HamiltonianCycleTransport} preserves the full single-cycle
sum under the same vertex equivalence.
\texttt{HamiltonianCoprimeEvenCount} now applies the full
bipartite counting theorem and restores actual arithmetic
weights and reversal classes, uniformly over every even modulus.
The full uniform totient density estimate is compiled in
\texttt{HamiltonianTotientDiffusion}.
\texttt{HamiltonianTotientSublog} supplies its quantified
squared-sublogarithmic consequence.
\texttt{HamiltonianCoprimeAperiodicCount} now applies the
actual odd-modulus spectral gap and all counting hypotheses:
the positive determinant reference and relative error
\(17n^{-1/10}\) hold at every sufficiently large odd modulus.
The full statement and retained normalization are recorded in
Section~\ref{sec:relativecoprime}.
\texttt{HamiltonianCoprimeDeterminant} also proves the exact
full characteristic determinant as the conductor-divisor product
recorded in Section~\ref{sec:relativecoprime}.
Its proof preserves the principal compensation and every
signed mode, including the possible even parity mode.
The positive square-root identification with the explicit
arithmetic constant is now compiled in
\texttt{HamiltonianCoprimeConstant}.
\texttt{HamiltonianCoprimeUnorientedCount} also proves the
actual odd unoriented factorial count, with relative error
at most \(17n^{-1/10}\) uniformly above one threshold.
The full signed logarithmic series and its exact leading
coefficient \(1\) are now compiled in
\texttt{HamiltonianCoprimeConstantSeries}.
\texttt{HamiltonianCoprimeEulerCoefficients} now identifies
every higher coefficient with its exact prime Euler product,
and the full prime-product logarithmic series is compiled.
\texttt{HamiltonianCoprimeConstantBounds} and
\texttt{HamiltonianCoprimeConstantTail} now prove the direct
distinct-prime dependence, uniform geometric remainder and
absolute bounds \(e\leq\mathfrak C(n)\leq e^{11/9}\) over
every odd modulus greater than one.
\texttt{HamiltonianCoprimeConstantLimit} now proves the actual
constant's monotonicity under prime-support inclusion, genuine
multipliability over all odd primes, and the complete primorial
limit with its absolutely convergent infinite-prime series.
Proposition~\ref{prop:coprimeconstant} is fully formalized.
\texttt{HamiltonianCoprimeStirlingCount} now proves the final
odd normalization with its positive displayed reference and
explicit relative error \(23n^{-1/10}\), keeping both the
factorial-count and Stirling errors.
The actual even count is also compiled.
\texttt{HamiltonianCoprimeRelativeCount} combines both parity
branches under one absolute threshold and proves the full
relative arithmetic enumeration with the actual displayed
references and common error at most \(23n^{-1/10}\).
Theorem~\ref{thm:relativecoprime} is fully formalized.
The actual cycle and prescribed-endpoint path entropy conclusions
are now both compiled as follows.

\paragraph{Maintained actual cycle entropy and normalized limit.}
\texttt{HamiltonianCoprimeLogReference} retains the full
parity-dependent logarithmic correction:
\[
 \log D_n-n(\log\varphi(n)-1)
 =\begin{cases}
 \log\pi,&n\ {\rm even},\\
 \log\mathfrak C(n)+
       \tfrac12(\log\pi-\log2-\log n),&n\ {\rm odd}.
 \end{cases}
\]
Its uniform constant bound therefore controls this entire
correction by \(O(\log(2n))\).

\texttt{HamiltonianCoprimeCycleEntropy} consumes the completed
actual all-modulus relative theorem. Above one absolute
threshold it proves \(h(X_n)>0\) and the stronger estimate
\[
 \left|\log h(X_n)-n(\log\varphi(n)-1)\right|
                 \leq K\log(2n),\qquad
 K=4+|\log\pi|+|\log2|.
\]
The relative ratio is eventually within \(1/2\) of one,
so its logarithm contributes at most one; no logarithm of
an unproved positive count is used.
The same owner proves the paper's cycle-error comparison,
using \(n\varphi(n)^{-1/8}\geq1\), and proves the exact
normalized limit formulation
\[
 \forall\varepsilon>0\ \exists N_\varepsilon\
 \forall n\geq N_\varepsilon:\quad
 \left|\frac{\log h(X_n)}n-(\log\varphi(n)-1)\right|
                                      \leq\varepsilon.
\]
All thirteen new declarations pass their focused foundational
audit. Both the cycle assertion and its stated entropy
consequence in Theorem~\ref{thm:coprimecircles} are complete.
\paragraph{Maintained actual integer path entropy.}
The module \texttt{HamiltonianPathCountNormalization} defines the
finite set of prescribed one-orbit path permutations. It proves
its exact membership conditions and identifies its cardinality
\(H_n(a,b)\) with the path sum for the zero--one adjacency.
Its scalar-weight theorem retains precisely \(n-1\) path edges.
The module \texttt{HamiltonianCoprimePathUpper} therefore proves
\[
 Z_{\rm path}(P_n;a,b)=\frac{H_n(a,b)}{\varphi(n)^{n-1}},
 \qquad
 Z_{\rm path}(P_n;a,b)
 \leq\exp\!\left(-(n-1)+16\frac n{\varphi(n)}\log(2n)\right).
\]
The lower count in Section~\ref{sec:relativecoprime} is positive
before any logarithm is taken.
The module \texttt{HamiltonianCoprimePathEntropy} combines the
complete signed lower count with this actual upper bound and
normalization. For one absolute \(N_0\geq512\), every \(n\geq N_0\)
and every compatible distinct pair satisfy \(H_n(a,b)>0\) and
the stronger compiled estimate
\[
 \left|\log H_n(a,b)-n(\log\varphi(n)-1)\right|
     \leq200\,\frac n{\varphi(n)}\log(2n).
\]
The same owner proves the stated weaker bound
\(200n\varphi(n)^{-1/8}\log(2n)\).
It also proves that, for even \(n\), oddness of
\((-(a-b))\bmod n\) is equivalent to opposite parity of the
actual residue representatives. Thus the final theorem has
exactly the prescribed endpoint domain in
Theorem~\ref{thm:coprimecircles}, with no remaining spectral,
count-normalization or endpoint hypothesis.
All ten new membership, count, edge-scaling, upper-count and
entropy declarations pass their combined focused foundational
axiom audit. Every clause of Theorem~\ref{thm:coprimecircles}
is now formalized; the other general weighted, linkage and
layered headline statements retain their own obligations.

\begin{remark}[Normalization, prior existence results, and scope]
When \(n\) is prime, \(X_n\) is complete and
\(h(X_n)=(n-1)!/2\).
When \(n=2^a\), it is complete balanced bipartite, and
\(h(X_n)=(n/2)!(n/2-1)!/2\).
Both exact formulas agree with the normalization above;
cyclic rotations and reversal have been identified, while
translations of the vertex labels have not.

Existence of a Hamilton cycle is not a new conclusion for
this family. Each unit step already gives a Hamilton cycle,
and its opposite gives the same one.
Nor do we claim novelty for the character spectrum or
Hamilton-connectedness.
The quantitative count is supplied to demonstrate a second
arithmetic consumer of the general entropy theorem, now with
the modulus itself growing and with both parity cases covered.
Priority for this counting conclusion requires comparison
beyond the existence and spectral results cited here.
\end{remark}

\section{Prime-residue pressure and the cost of an atypical profile}
\label{sec:residuepressure}
The weighted theorem has an arithmetic consequence beyond the
unweighted total count.
Keep a fixed even modulus \(q\), a nonempty reduced residue set
\(R\) with \(g_R=2\), and write \(k=|R|\).
For each unoriented Hamilton cycle \(C\) of \(G_N(q,R)\), put
\[
 n_r(C)=|\{ij\in E(C):i+j\equiv r\pmod q\}|,\qquad r\in R.
\]
These are edge counts, including repeated prime values if they
occur on different edges. Every cycle has exactly \(N\) edges,
so \(\sum_{r\in R}n_r(C)=N\).
For fixed weights \(w_r>0\), define
\[
 W_N(w)=\sum_{C\ {\rm Hamilton}}\prod_{r\in R}w_r^{\,n_r(C)},
 \qquad
 M(w)=\frac{\sum_{r\in R}w_r}{\varphi(q)}.
\]
The weights, modulus and allowed set are fixed before \(N\)
tends to infinity. No uniformity for growing moduli or
\(N\)-dependent weights is asserted.

\begin{theorem}[The full positive residue-pressure formula]
\label{thm:residuepressure}
Along even \(N\),
\[
 \log W_N(w)
 =N\left(\log\frac{N}{e\varphi(q)\log N}
                      +\log\sum_{r\in R}w_r\right)+o_w(N).
\]
In particular, if \(C_N\) is a uniform unoriented Hamilton cycle
of \(G_N(q,R)\) and \(X_N=(n_r(C_N)/N)_{r\in R}\), then for every
fixed \(\theta\in\mathbb R^R\),
\[
 \lim_{N\to\infty,\ 2\mid N}\frac1N
      \log\mathbb E\exp\bigl(N\langle\theta,X_N\rangle\bigr)
       =\log\left(\frac1k\sum_{r\in R}e^{\theta_r}\right).
\]
\end{theorem}

\begin{proof}
The positive-weight version of the signed Fourier model is
obtained by replacing the uniform average \(\rho_R\) by
\[
 \rho_w(b)=
       \frac{\sum_{r\in R}w_r\e(br/q)}{\sum_{r\in R}w_r}.
\]
For completeness, each matrix entry is a finite linear
combination of the same fixed-class prime convolution
operators used in Lemma~\ref{lem:filtermodel}.
The Chinese remainder computation there consequently gives
the exact reflection coefficient
\[
 w(\gamma)=\rho_w(b)\frac{\mu(c)}{\varphi(c)},\qquad
 \gamma=b/q+a/c,\quad (c,q)=1,
\]
with the same squarefree and primitive-conductor conditions
as in (38), and zero in every other case.
All phases are retained before taking eigenvalues.
A positive weighted average of unit complex numbers has
modulus one precisely when all the phases agree.
Its set of unit-modulus characters is therefore unchanged:
\(b=0,q/2\), because \(g_R=2\).
The nontrivial gap is the positive fixed number
\[
 \delta_w=
 1-\max\left\{\frac1{p_q-1},
       \max_{\substack{b\bmod q\\b\ne0,\ q/2}}|\rho_w(b)|
          \right\}>0.
\]
The finite-mask operator approximation in (39) is linear
and has only a fixed weight-dependent constant.
It therefore proves the corresponding spectral approximation,
not merely a pointwise residue heuristic.

Consider the true-prime weighted kernel
\[
 Q_w(i,j)=
 \frac{w_{\,i+j\bmod q}\log(i+j)}{M(w)N}
      \one_{\substack{i\ne j,\ i+j\ {\rm prime}\\
                       i+j\bmod q\in R}}.
\]
The fixed-progression prime number theorem gives its row
defect \(O_{q,w,B}((\log N)^{-B})\) for every fixed \(B\).
Proper prime powers and the diagonal are removed with the
same row-norm bound as in the filtered model.
Thus its normalized spectrum has only the constant and
parity exceptions and gap \(\delta_w-o(1)\).
Lemma~\ref{lem:scaling}, with any fixed \(B>4\), produces
\[
 P_w=\diag(u)Q_w\diag(u),\qquad
 P_w\one=\one,\quad
 \|u-\one\|_\infty=o(1).
\]
This kernel has a fixed nontrivial gap, the original true-prime
support and maximum entry \(O_{q,w}(\log N/N)\).
It satisfies Theorem~\ref{thm:weightedentropy}, giving
\[
 \log Z_{\rm H}(P_w)=-N+o_w(N).
\]
Every Hamilton-cycle weight changes by the common factor
\(\prod_vu_v^2=\exp(o_w(N))\), so the same logarithmic formula
holds for \(Z_{\rm H}(Q_w)\).

One must still remove the nonconstant logarithmic prime weight.
This cannot be justified by a uniform entrywise \(1+o(1)\)
comparison at the smallest sums.
Put \(D_w=M(w)N/\log N\). Exactly,
\[
 \prod_{ij\in E(C)}Q_w(i,j)
 =D_w^{-N}\left(\prod_{r\in R}w_r^{n_r(C)}\right)
             \prod_{ij\in E(C)}\frac{\log(i+j)}{\log N}.
\]
The logarithm of the last product is \(o(N)\), uniformly
over all such cycles.
Indeed, for sums at least \(N/(\log N)^4\), each logarithm
has modulus \(O(\log\log N/\log N)\).
Every smaller sum has both endpoints in the first
\(N/(\log N)^4\) vertices; a cycle has at most twice that
many edges incident with this set.
For each of these edges the logarithm has modulus
\(O(\log\log N)\), since the smallest permitted sum is at least 3.
The complete exceptional contribution is
\(O(N\log\log N/(\log N)^4)=o(N)\).
Consequently
\[
 Z_{\rm H}(Q_w)=D_w^{-N}W_N(w)\exp(o_w(N)).
\]
This proves the partition-function formula.
Divide by its specialization \(w_r=1\), and then put
\(w_r=e^{\theta_r}\), to obtain the moment-generating formula.
\end{proof}

Let \(\Delta_R=\{x\in[0,1]^R:\sum_rx_r=1\}\).
Use \(0\log0=0\) in
\[
 I_R(x)=\sum_{r\in R}x_r\log(kx_r),\qquad x\in\Delta_R.
\]
Outside this simplex set \(I_R=+\infty\).

\begin{theorem}[A full empirical residue large-deviation law]
\label{thm:residueldp}
The random vectors \(X_N\) satisfy the large-deviation principle
with speed \(N\) and rate \(I_R\).
Precisely, for every relatively open \(O\subseteq\Delta_R\)
and relatively closed \(F\subseteq\Delta_R\),
\[
 \begin{split}
 \liminf_{N\to\infty,\ 2\mid N}
       N^{-1}\log\mathbb P(X_N\in O)&\geq-\inf_{x\in O}I_R(x),\\
 \limsup_{N\to\infty,\ 2\mid N}
       N^{-1}\log\mathbb P(X_N\in F)&\leq-\inf_{x\in F}I_R(x).
 \end{split}
\]
Here \(\log0=-\infty\).
In particular, for each \(x\in\Delta_R\), if
\(B_\epsilon(x)=\{y\in\Delta_R:\|y-x\|_\infty<\epsilon\}\),
then both
\[
 \lim_{\epsilon\downarrow0}\liminf_N
       N^{-1}\log\mathbb P(X_N\in B_\epsilon(x)),
 \qquad
 \lim_{\epsilon\downarrow0}\limsup_N
       N^{-1}\log\mathbb P(X_N\in B_\epsilon(x))
\]
equal \(-I_R(x)\), with the inner limits along even \(N\).
\end{theorem}

\begin{proof}
We give the exponential-tilt proof to expose the quantifiers.
Write
\(\psi(\theta)=\log(k^{-1}\sum_re^{\theta_r})\).
Markov's inequality and Theorem~\ref{thm:residuepressure}
give, for each fixed \(\theta,x,\epsilon\),
\[
 \limsup_N N^{-1}\log\mathbb P(X_N\in B_\epsilon(x))
 \leq \psi(\theta)-\langle\theta,x\rangle
                                      +\epsilon\|\theta\|_1.
\]
If all \(x_r>0\), take \(\theta_r=\log(kx_r)\), for which
\(\psi(\theta)=0\) and \(\langle\theta,x\rangle=I_R(x)\).
If some coordinates vanish, set their \(\theta_r=-M\)
and use the same positive-coordinate values, then let
\(M\to\infty\) after taking the \(N\) and \(\epsilon\) limits.
This proves the local upper bound.
A finite cover of the compact simplex by the resulting
neighbourhoods gives the upper bound for every closed \(F\).

For the lower bound first suppose all \(x_r>0\).
Tilt the law of cycles by \(\prod_rx_r^{n_r(C)}\), and denote
the resulting law by \(\nu_{N,x}\).
Its limiting logarithmic moment-generating function is
\[
 \log\sum_rx_re^{\theta_r},
\]
by applying the partition-function formula at the two fixed
positive weight vectors \(x\) and \((x_re^{\theta_r})_r\).
For each coordinate, a sufficiently small fixed positive or
negative tilt and Chernoff's inequality show
\[
 \nu_{N,x}(B_\epsilon(x))=1-\exp(-\Omega_{x,\epsilon}(N)).
\]
For a constant coordinate the corresponding deviation event
is empty. There are only finitely many coordinates.
On \(B_\epsilon(x)\),
\[
 \left|\sum_r(X_{N,r}-x_r)\log x_r\right|
                        \leq\epsilon\sum_r|\log x_r|.
\]
The exact Radon--Nikodym identity is
\[
 \mathbb P(X_N\in B_\epsilon(x))
 =\frac{W_N(x)}{W_N(\one)}
   \mathbb E_{\nu_{N,x}}\left[
       \exp\left(-N\sum_r X_{N,r}\log x_r\right)
                        \one_{B_\epsilon(x)}(X_N)\right].
\]
The pressure ratio has logarithm \(-N\log k+o_x(N)\).
Consequently the lower limiting exponent is at least
\(-I_R(x)-\epsilon\sum_r|\log x_r|\).
Send \(\epsilon\) to zero to obtain the local lower bound.

Every boundary point of the simplex is a limit of interior
points, and \(I_R\) is continuous on the whole simplex.
Apply the interior lower bound in a smaller ball about an
interior point contained in the prescribed neighbourhood.
This proves the boundary local lower bound as well.
For an open \(O\), choose any such ball contained in \(O\),
and optimize its centre. This proves the stated principle,
without requiring a theorem uniform in the tilt vector.
\end{proof}

\begin{corollary}[Exponential concentration and profile counts]
\label{cor:primeprofilecount}
For every fixed \(\epsilon>0\) there is \(c_{q,R,\epsilon}>0\)
such that a uniform prime circle in \(G_N(q,R)\) satisfies
\[
 \mathbb P\left\{\max_{r\in R}
          \left|\frac{n_r(C_N)}N-\frac1k\right|\geq\epsilon\right\}
           \leq\exp(-c_{q,R,\epsilon}N)
\]
for all sufficiently large even \(N\), unless the event is empty.
For a prescribed approximate profile \(x\), its fraction of
all prime circles has the exact exponential cost \(I_R(x)\)
in the iterated neighbourhood-limit sense of the theorem.
Thus the total count and the empirical-profile count are
determined at exponential precision.
\end{corollary}
\begin{proof}
The rate is nonnegative and vanishes only at the uniform
vector. Its infimum on the indicated closed set is strictly
positive, whenever that set is nonempty.
The large-deviation upper bound gives the displayed estimate
with any fixed smaller positive exponent.
The local formula gives the remaining assertion.
\end{proof}

\begin{remark}[The edges are constrained, not independent]
Each prime circle obeys the exact identity
\[
 \sum_{r\in R}r\,n_r(C)
       \equiv\sum_{ij\in E(C)}(i+j)
       =2\sum_{i=1}^N i=N(N+1)\pmod q.
\]
The large-deviation law agrees with independent uniform
residue sampling only at the exponential profile scale.
It does not assert independence, a central limit theorem,
or positive counts for every exact integer profile.
For example \(q=8,R=\{1,3\}\) satisfies \(g_R=2\).
The number \(n_3(C)\) always obeys
\[
 n_3(C)\equiv N^2/2\pmod4.
\]
Nevertheless its normalized proportion has binary rate
\[
 I(x)=x\log(2x)+(1-x)\log(2(1-x)).
\]
An exact profile using only one class has no Hamilton cycle
for large \(N\), by the obstruction theorem.
A fixed positive-width neighbourhood of that boundary
profile has the finite exponential rate \(I(0)=I(1)=\log2\).
The order of limits is essential here.
\end{remark}

\section{Relative permanent estimates with a growing entry bound}
\label{sec:relativepermanents}
The preceding entropy bounds do not determine a relative
\(1+o(1)\) count.
For that purpose the entry parameter in a permanent asymptotic
must be tracked; a theorem uniform only for a fixed bound
cannot simply be applied when that bound grows.
Throughout this section \(J_N\) is the all-ones matrix.

\begin{theorem}[Polylogarithmic entry bounds in the permanent expansion]
\label{thm:growingpermanent}
Fix \(\delta\in(0,1)\) and \(K>0\).
Let \(D\) be real nonnegative doubly stochastic of order \(N\),
and write
\[
 B=D-J_N/N,\qquad
 \|B\|\leq1-\delta,\qquad
 \max_{ij}D_{ij}\leq A/N,\qquad 2\leq A\leq(\log N)^K .
\]
Then, uniformly over these matrices for sufficiently large \(N\),
\[
 \operatorname{per}D
 =\frac{N!}{N^N}\det(I-B^\top B)^{-1/2}
                       \bigl(1+O_{\delta,K}(N^{-1/8})\bigr).
\]
Zero entries are allowed.
\end{theorem}

\begin{proof}
We track the finite estimates behind the fixed-entry theorem
of Divoux--Wang~\cite[Section 3]{DW}, based on McCullagh~\cite{McC}.
This is a parameter-uniformity argument, not an application
of their fixed-\(A\) asymptotic statement.

Stochasticity gives the useful stronger bounds
\[
 \max_i\sum_jB_{ij}^2,\quad
 \max_j\sum_iB_{ij}^2\leq A/N,\qquad
 \|B\|_F^2\leq A.
\]
Choose fixed \(0<a<1<b\) and \(R>1\) so that
\(bR(1-\delta/2)<1\).
Let \(R_B(s)\) be the rook polynomial
\[
 R_B(s)=\sum_{k=0}^N s^k
       \sum_{\substack{I,J\subseteq[1,N]\\|I|=|J|=k}}
                         \operatorname{per}B[I,J].
\]
The exact rook/Gamma identity, with
\(T\sim\Gamma(N+1,1)\), is
\[
 \frac{N^N\operatorname{per}D}{N!}
                         =\mathbb E R_B(N/T).
\]
Define
\[
 F_s(t)=R_{tB}(s)
             \det(I-s^2t^2B^\top B)^{1/2},
 \qquad s\in[a,b],
\]
with the analytic square root equal to one at \(t=0\).
We show that \(F_s(1)=1+O_\delta(A^2N^{-1/6})\),
uniformly in \(s\).

Here are the finite Gaussian estimates and their parameter
dependence explicitly.
With
\[
 \Omega_{s,t}=\begin{pmatrix}0&stB\\stB^\top&0\end{pmatrix},
 \quad
 \Gamma_{s,t}=\Omega_{s,t}(I+\Omega_{s,t})^{-1},
\]
the normalized Gaussian identity in~\cite[Lemmas 3.5--3.7]{DW}
writes \(F_s(t)\) as the Gaussian expectation of
\(\prod_i r(Y_i)\), where
\[
 r(z)=(1+z)\exp(-z+z^2/2)=1+O(z^3)
\]
and \(Y\) has bilinear covariance \(\Gamma_{s,t}\).
The zero row and column sums of \(B\) cancel the linear
Gaussian term exactly, before taking any bound.

The polynomial transfer underlying this identity has an independent
maintained proof. For any finite coordinate set, choose complex numbers
\(z_j,w_j\) with \(|z_j|<1\) and
\(w_j^2=\gamma_j=z_j/(1+z_j)\), and put
\(b_j=(1-\gamma_j)/2\). The complete quotient gives
\[
 (1-2\Re\gamma_j)|1+z_j|^2=1-|z_j|^2,
 \qquad \Re b_j>1/4.
\]
Every polynomial times the resulting complex Gaussian density is
integrable, and each scalar normalizer is nonzero. Define the actual
normalized integral
\[
 \mathcal G_z(P)=
 \frac{\displaystyle\int_{\mathbb R^J}
      P((w_jx_j)_j)\exp(-\sum_jb_jx_j^2)\,dx}
      {\displaystyle\prod_j\int_{\mathbb R}\exp(-b_jx^2)\,dx}.
\]
Scalar integration by parts and finite-dimensional Fubini prove
\[
 \mathcal G_z(X^d)=\prod_j W_{z_j}(d_j),\quad
 W_z(2e)=(2e-1)!!z^e,\quad W_z(2e+1)=0,\quad W_z(0)=1,
\]
and, for every polynomial \(P\),
\(\mathcal G_z(X_jP)=z_j\mathcal G_z(\partial_jP)\).
For \(L_a=\sum_ja_jX_j\), set
\(C_z(a,b)=\sum_jz_ja_jb_j\). The complete recursion is
\[
 \mathcal G_z\left(L_a\prod_{\ell=1}^mL_{a_\ell}\right)
 =\sum_{k=1}^m C_z(a,a_k)\,
    \mathcal G_z\left(\prod_{\substack{1\leq\ell\leq m\\\ell\ne k}}
                                      L_{a_\ell}\right).
\]
In particular, if every covariance among \(2e\) listed forms has
modulus at most \(c\geq0\), the actual normalized integral satisfies
\[
 \left|\mathcal G_z\left(\prod_{\ell=1}^{2e}L_{a_\ell}\right)\right|
                         \leq(2e-1)!!c^e .
\]
All these statements, including integrability, all repeated-vertex
contractions and the uniform bound in \(e\), have maintained Lean
proofs. The bound is used on the diagram remainder after the signed
low-degree terms have been retained.

The restriction \(|z_j|<1\) belongs to the quadratic transfer, not
to the existence of a Gaussian realization. For arbitrary complex
\(v_j\), choose \(u_j^2=v_j\). Standard Gaussian integration gives
\[
 \frac{\displaystyle\int_{\mathbb R^J}
      P((u_jx_j)_j)e^{-\sum_jx_j^2/2}\,dx}
      {Z_0^{|J|}}=\mathcal G_v(P),\qquad
 Z_0=\int_{\mathbb R}e^{-x^2/2}\,dx=\sqrt{2\pi}.
\]
This unrestricted realization, including the equality of the real and
complex scalar normalizations, is formalized. Thus an indefinite
or large compensated covariance is not excluded by a spurious
positive-semidefinite or norm restriction.

The exact rook identification is also formalized. Let \(L_i,R_j\),
\(1\leq i,j\leq N\), be linear forms whose within-left and within-right
covariances vanish. For every \(S\subseteq\{1,\ldots,N\}\), the maintained
support identity is
\[
 (N-|S|)!\,\mathcal G_v\!\left(
       \prod_{i\in S}L_i\prod_{j=1}^{N}(1+R_j)\right)
 =\sum_{\sigma\in\mathfrak S_N}
       \prod_{i\in S}\mathcal G_v(L_iR_{\sigma(i)}).
\]
For nonempty \(S\), delete one selected vertex \(i\) and contract it
with each right vertex \(j\). The support loses one element and
\(N\) decreases by one, so the factorial is identical on every branch.
The full permutation decomposition at column \(i\) has exactly these
branches. For empty \(S\), the Gaussian expectation is one and
the permutation count is \(N!\). Summing every support and dividing
by its nonzero factorial gives the complete signed affine-product
rook formula, without a discarded contraction or an assumed independence.

For a real symmetric matrix \(M=U\diag(\lambda_j)U^\top\), let
\(Y_i=\sum_jU_{ij}\sqrt{\lambda_j}X_j\), where the square roots
are complex and the \(X_j\) are independent real standard normals.
The maintained spectral proof derives
\[
 \mathcal G_\lambda(L_iL_k)=M_{ik},\qquad
 W:=\sum_i|Y_i|^2=\sum_j|\lambda_j|X_j^2.
\]
If \(\one^\top M=0\), it also proves \(\sum_iY_i=0\) pointwise.
Both identities retain every cross term until orthogonality cancels it.
For the actual \(M=\Omega\) from a doubly stochastic \(D\),
symmetry and this zero-sum condition are discharged directly.

The scalar Gaussian integral and finite-dimensional Fubini give, for
\(\tau\geq0\), \(\|M\|\leq\rho\) and \(\tau\rho<1\),
\[
 \E e^{\tau W/2}
 \leq\exp\!\left(\frac{\tau}{1-\tau\rho}\sum_j|\lambda_j|\right).
\]
The exact Frobenius identity
\(\sum_j\lambda_j^2=\|M\|_F^2\) is formalized.
For \(M=\Omega\), stochasticity gives
\(\|M\|_F^2=2\|B\|_F^2\leq2(A-1)\), hence
\[
 \E e^{\tau W/2}
 \leq\exp\!\left(\frac{2\tau}{1-\tau\rho}
                                      \sqrt{N(A-1)}\right).
\]
Finally the exact zero-sum product bound, with \(c=\sqrt t\),
proves the actual collective integral estimate
\[
 \boxed{\quad
 \E\prod_{i=1}^{2N}|1+\sqrt t\,Y_i|
 \leq\exp\!\left(\frac{2|t|}{1-|t|\rho}
                                     \sqrt{N(A-1)}\right),
 \quad |t|\rho<1.\quad}
\]
Every hypothesis in this stochastic consumer, including integrability,
is discharged in maintained Lean. This is the Cauchy/tail input;
the signed low-degree corrections are retained separately.
The exact Gaussian-to-rook consumer now transfers this bound directly:
\[
 \boxed{\quad |R_B(t)|\leq
 \exp\!\left(\frac{2|t|}{1-|t|\rho}\sqrt{N(A-1)}\right),
 \qquad |t|\rho<1.\quad}
\]
The maintained complex polynomial agrees with the original real rook
polynomial at every real parameter. The compensated resolvent
identification and the complete diagram, Cauchy and outside-Gamma
consumers remain formalization obligations.

The block-power expansion gives, for a fixed
\(\theta=b(1-\delta/2)<1\),
\[
 \max_{i,j}|[t^\ell](\Gamma_{s,t})_{ij}|
       \leq(2b^2 A/N)\theta^{\ell-1},
 \qquad \ell\geq1.
\]
Indeed each power of length at least two is bounded using
the two row/column Euclidean norms above and the intervening
operator power. The length-one term is at most \(bA/N\).
Thus the covariance constant grows linearly in \(A\),
not as an unspecified function. The actual real matrix geometric
series and this bound for every degree are formalized; the complex
analytic identification used here is a further consumer.

For clarity the finite Wick coefficient estimate used here is:
if \(F_s(t)=1+\sum_{k\geq1}f_k(s)t^k\), then
\[
 |f_k(s)|\leq
 \sum_{e=1}^k\ \sum_{1\leq v\leq\lfloor2e/3\rfloor}
 (C_\delta A)^e(e/N)^{e-v}
                  \binom{k-1}{e-1}\theta^{k-e}.
\]
This is the explicit estimate in~\cite[Lemma 3.7]{DW}.
Its constant is an absolute constant times the covariance
constant: Wick pairings have \(e\) pairs, the nonconstant
vertices have degree at least three, hence \(v\leq2e/3\),
and their choices supply \(N^v\) against \(N^{-e}\).
There is no hidden fixed-entry asymptotic in this finite bound.

Take \(H=\lceil C_\delta A^3\sqrt N\rceil\), increasing its
constant below if necessary. For large \(N\), \(H<N\).
Sum over \(k\leq H\), use
\(\sum_{k\geq e}\binom{k-1}{e-1}\theta^{k-e}
=(1-\theta)^{-e}\), and retain every degree choice.
Since \(e/N\leq H/N\) and \(e-v\geq e/3\),
\[
 \sum_{1\leq k\leq H}|f_k(s)|
 \leq\sum_{e\geq1}e
       \bigl(C_\delta A^2N^{-1/6}\bigr)^e
 =O_\delta(A^2N^{-1/6}).
\]
The base tends to zero throughout the prescribed range of \(A\).

The same Gaussian product estimate on \(|t|=R\) gives
\[
 \max_{|t|=R}|F_s(t)|
 \leq\exp\bigl(C_\delta\sqrt{NA}+C_\delta A\bigr).
\]
To see its dependence, the Gaussian exponential moment is
bounded by a product \(\prod_\ell(1-bR\sigma_\ell)^{-1}\);
its logarithm is at most
\(C_\delta\sum_\ell\sigma_\ell\leq C_\delta\sqrt{NA}\).
The additional determinant factor costs at most
\(C_\delta\sum_\ell\sigma_\ell^2\leq C_\delta A\).
Cauchy's coefficient bound and the chosen \(H\) now give
\[
 \sum_{k>H}|f_k(s)|
       \leq\exp(-c_\delta A^3\sqrt N).
\]
Consequently the asserted relative rook estimate follows.

It remains to average the determinant factor over \(T\);
replacing it by an absolute bounded constant would lose
uniformity. Put
\(\Delta(s)=\det(I-s^2B^\top B)^{-1/2}\).
On \([a,b]\),
\[
 0\leq\log\Delta(s)\leq C_\delta A,\qquad
 |(\log\Delta)'(s)|\leq C_\delta A .
\]
On \(|N/T-1|\leq N^{-1/4}\), this gives
\(\Delta(N/T)/\Delta(1)=1+O_\delta(AN^{-1/4})\).
The remaining Gamma mass inside \([a,b]\) is at most
\(\exp(-c_\delta\sqrt N)\), while the determinant ratio is
at most \(\exp(C_\delta A)\).

For completeness the outside-\([a,b]\) Gaussian/Gamma tail
has absolute bound
\[
 \mathbb E\bigl[\one_{\{N/T\notin[a,b]\}}
                         |R_B(N/T)|\bigr]
 \leq\exp(-c_\delta N+C_\delta\sqrt{NA}).
\]
Here is the collective cancellation underlying the change of
variable. In a real standard-Gaussian representation let
\(Y\in\mathbb C^{2N}\), \(\sum_iY_i=0\), and
\(W=\sum_i|Y_i|^2\). For every complex \(c\),
\[
 \sum_i|1+cY_i|^2=2N+|c|^2W,\qquad
 \prod_i|1+cY_i|\leq\exp(|c|^2W/2).
\]
The linear term vanishes before either bound is taken. Arithmetic
and geometric means, with \(c=\sqrt{N/t}\), also give
\[
 \prod_i|1+\sqrt{N/t}Y_i|
       \leq(1+W/(2t))^N\qquad(t>0).
\]
Multiplication by the actual Gamma density cancels its full power:
\[
 \frac{e^{-t}t^N}{N!}\prod_i|1+\sqrt{N/t}Y_i|
 \leq\frac{e^{W/2}e^{-(t+W/2)}(t+W/2)^N}{N!}.
\]
These deterministic inequalities have maintained Lean proofs.
Set \(w=W/2\) and shift \(u=t+w\) in the remaining integral.
On \(w\leq\eta N\), a sufficiently small fixed \(\eta>0\)
leaves a fixed Gamma-tail separation. On \(w>\eta N\), use
\[
 \E[e^w\one_{\{w>\eta N\}}]
 \leq e^{-(\lambda-1)\eta N}\E e^{\lambda w},
 \qquad 1<\lambda<(1-\delta)^{-1}.
\]
Diagonalizing the actual block matrix gives
\(W=\sum_\ell|\lambda_\ell|Z_\ell^2\) for independent
real standard normals. The paired eigenvalues yield
\(\E e^{\lambda W/2}
 =\prod_\ell(1-\lambda\sigma_\ell)^{-1}
 \leq\exp(C_\delta\sum_\ell\sigma_\ell)\).
This proves the stated tail with its Frobenius dependence retained,
as in~\cite[Section 3.3]{DW}. The Gaussian representation,
integration and tail consumer are still formalization obligations.
They vanish uniformly for polylogarithmic \(A\).

As \(\Delta(1)\geq1\), the total relative error is bounded by
\[
 C_\delta(A^2N^{-1/6}+AN^{-1/4})
 +\exp(C_\delta A-c_\delta\sqrt N)
 +\exp(-c_\delta N+C_\delta\sqrt{NA}).
\]
For fixed \(K\) this is \(O_{\delta,K}(N^{-1/8})\), as claimed.
\end{proof}

\begin{lemma}[Uniform compensated minors with mismatched supports]
\label{lem:relativecompensatedminors}
Under Theorem~\ref{thm:growingpermanent}, uniformly for arbitrary
row and column sets \(I,J\) with
\(|I|=|J|=r\leq(\log(2N))^3\),
\[
 \frac{\operatorname{per}D[I^c,J^c]}{\operatorname{per}D}
                  =1+O_{\delta,K}(N^{-1/9}).
\]
The removed supports need not coincide.
\end{lemma}
\begin{proof}
The case \(r=0\) is immediate. Otherwise put
\(n=N-r\), \(Q=D[I^c,J^c]\), and
\(s_{I,J}=\sum_{i\in I,j\in J}D_{ij}\).
Each remaining row or column loses at most
\(\varepsilon=rA/N\), and its total lost mass is exactly
\(r-s_{I,J}\). Also \(s_{I,J}\leq r^2A/N\).

The second singular value of \(Q\), after degree normalization,
stays a fixed distance from 1.
Indeed restrict \(B\) and compare \(J_n/N\) with \(J_n/n\);
their rank-one difference has norm \(r/N\).
The row and column normalization changes the operator by
\(O(\varepsilon)\).
The symmetric bipartite form of \(Q\) thus satisfies
Lemma~\ref{lem:scaling}, since
\(\varepsilon(\log N)^2\to0\).
It gives positive row and column multipliers \(u,v\) with
\[
 \widetilde D=\diag(u)Q\diag(v),\quad
 \max(|u_i-1|,|v_j-1|)
                     \leq C_\delta\varepsilon\log N.
\]

The total scaling cost requires compensation, not just
the maximum multiplier. In the symmetric bipartite form
write \(z=(u-\one,v-\one)\), and denote its row degrees by \(d_i\).
Sum the exact scaling equations
\(Q_{\rm bip}z+d\mathbin{\odot}z+
 z\mathbin{\odot}(Q_{\rm bip}z)=\one-d\)
over all \(2n\) coordinates:
\[
 \sum_i d_i z_i
       =r-s_{I,J}-\tfrac12 z^\top Q_{\rm bip}z .
\]
Using \(1-d_i\leq\varepsilon\), the maximum bound on \(z\),
and the full quadratic term, we obtain
\[
 \sum_i\log u_i+\sum_j\log v_j
 =r-s_{I,J}
        +O_\delta(r^2A^2(\log N)^2/N).
\]
This retains the overlap mass \(s_{I,J}\) and the nonlinear
compensation term.

Put \(\widetilde B=\widetilde D-J_n/n\).
After padding its row and column sets by zero, its Frobenius
distance from \(B\) is at most
\[
 C_\delta\left(\sqrt{rA/N}+r/N+
                     rA^{3/2}\log N/N\right).
\]
The first term removes the actual rows and columns, the
second changes the constant block, and the third includes
both scaling multipliers.
Since \(\|B\|_F\leq\sqrt A\), the trace-norm distance of the
squared singular-value matrices is at most
\[
 C_\delta\left(A\sqrt{r/N}+r\sqrt A/N+
                                      rA^2\log N/N\right).
\]
Both squared matrices have norm bounded away from one.
The trace-norm log-determinant inequality therefore shows
that their determinant factors differ by
\(\exp(O_{\delta,K}(N^{-1/4}))\).
All displayed errors are powers of \(\log N\) times \(N^{-1/2}\)
or \(N^{-1}\), uniformly over the supports.

Apply Theorem~\ref{thm:growingpermanent} to \(D\) and
\(\widetilde D\). The latter has a fixed gap and entry bound
at most \(3A/n\), fitting a possibly larger fixed
polylogarithmic exponent. Exactly,
\[
 \operatorname{per}Q=
                 \operatorname{per}\widetilde D/(\prod_i u_i\prod_jv_j).
\]
Stirling's formula gives
\[
 \log\frac{n!/n^n}{N!/N^N}=r+O(r^2/N).
\]
Its leading \(r\) cancels the full leading scaling cost.
The overlap term is \(O(r^2A/N)\); all the remaining scaling,
determinant and factorial errors are \(o(N^{-1/8})\).
The two relative permanent errors are \(O(N^{-1/8})\).
This proves the stated uniform \(O(N^{-1/9})\) comparison.
\end{proof}

\section{Signed Hamiltonian-permanent comparison and parity cancellation}
\label{sec:relativehamilton}
Write \(\operatorname{cyc}D\) for the sum of products over
directed cyclic permutations. For a symmetric kernel of order
at least three it equals \(2Z_{\rm H}(D)\).
The exact identity of Liu~\cite{Liu} is
\[
 N\operatorname{cyc}D
 =\sum_{\emptyset\ne U\subseteq[1,N]}
       |U|(-1)^{|U|-1}\det D[U]\,
                                  \operatorname{per}D[U^c].
\]
We use this signed identity before estimating errors.
Its principal terms must not be replaced by their absolute mass.

The full identity now has the maintained Lean proof
\texttt{hamiltonianLiuIdentity}, over every commutative coefficient ring.
For clarity, the proof expands the actual determinant and permanent and
regroups their products by the resulting full permutation \(\sigma\).
The selected vertex set must be invariant under \(\sigma\), so it is
uniquely a union of complete orbits, including singleton fixed points.
For a selected orbit family \(\mathcal C\), the restricted sign satisfies
\[
 (-1)^{|\bigcup\mathcal C|}
       \operatorname{sgn}(\sigma|_{\bigcup\mathcal C})
                         =(-1)^{|\mathcal C|}.
\]
Consequently the full coefficient of the permutation product is
\[
 -\sum_{\mathcal C\subseteq\operatorname{Orb}(\sigma)}
       (-1)^{|\mathcal C|}|\bigcup\mathcal C|
 =N\,\mathbf1_{\{|\operatorname{Orb}(\sigma)|=1\}}.
\]
Indeed each orbit length is multiplied by the alternating subset sum
over all other orbits. This sum is zero unless there are no other
orbits. The empty selected set contributes zero, including in
dimension zero. The complementary-permutation bijection and its
matrix-product identity are part of the formal proof, so no fixed-point,
overlap or completion term is omitted. This formalizes the established
identity; it is not a new enumeration claim.

\paragraph{Formalized weighted row-cap permanent bound.}
This bound is used only for a separately retained tail.
Let \(M\) be nonnegative of order \(n\geq2\), with row sums at
most \(1\) and entries at most \(c>0\).
Write \(w_\sigma=\prod_iM_{i,\sigma(i)}\) and
\(P=\sum_\sigma w_\sigma\).
If \(P=0\), every upper bound below is immediate.
For each fixed row order, weighted sequential revelation gives
a subprobability \(q_\sigma\), with its actual available-column
denominators. If \(C_\sigma\) is the sum of their logarithms,
the inequality \(\log x\leq x-1\) gives
\[
 \log P\leq P^{-1}\sum_\sigma w_\sigma C_\sigma.
\]
Zero matching weights contribute zero; no positive-entry
approximation is needed.

Fix a matching of positive weight and average over all row orders.
When row \(i\) is revealed with \(k\) other rows remaining,
the latter are a uniform subset of the other \(n-1\) rows.
The available mass therefore has conditional mean
\[
 M_{i,\sigma(i)}
       +\frac{k}{n-1}\sum_{j\ne i}M_{i,\sigma(j)}
                         \leq c+\frac{k}{n-1}.
\]
Applying the logarithmic mean inequality at each stage, then
averaging with the full matching weights, proves
\[
 \operatorname{per}M\leq
 \exp\!\left(\sum_{k=0}^{n-1}\log\left(c+\frac{k}{n-1}\right)\right).
\]
For \(m=n-1\), the logarithm is at most
\[
 \log c-m+\log m+1+cm(1+\log m),
\]
by \(\log(1+u)\leq u\), the harmonic bound and the Stirling
upper bound for the factorial logarithm.
Consequently, if \(A\geq1\), \(N\geq1\), \(n\leq N\) and
\(c=A/N\), then
\[
 \operatorname{per}M\leq \exp(-n+16A\log(2N)).
\]
The same conclusion for \(n=0,1\) follows directly.
The weighted argument, explicit constant, exceptional orders and
transfer to every actual independent-support rectangular minor
have maintained Lean proofs. In particular, for a doubly stochastic
matrix \(D\) with maximum entry \(A/N\),
\[
 \operatorname{per}D[I^c,J^c]
                       \leq \exp(r-N+16A\log(2N)),\qquad |J|=r.
\]
The denominator bound \(N!/N^N\geq e^{-N}\), together with
van der Waerden, gives the ordinary-proof ratio estimate
\[
 \frac{\operatorname{per}D[I^c,J^c]}{\operatorname{per}D}
                       \leq\exp(r+16A\log(2N)).
\]
For the admissible stochastic expander kernels at sufficiently large
order, the maintained proof instead consumes the already formalized
normalized positivity data.  It proves
\(\operatorname{per}D\geq e^{-N}/2\), hence the all-rank bound
\[
 0\leq\frac{\operatorname{per}D[I^c,J^c]}{\operatorname{per}D}
                       \leq2\exp(r+16A\log(2N)).
\]
The factor \(2\) is retained throughout the full compensated bipartite
tail.  Its tilted Cauchy--Binet summation, including the constant and
first-order subtractions, now has a maintained proof and gives
\(\exp(-\log(2N)^3\log\log(2N)/4)\) above one uniform threshold
when \(A\leq C_0\log(2N)\).  The proof and its explicit scalar budget
are given in the signed comparison section.

\begin{theorem}[A relative comparison with a growing density parameter]
\label{thm:relativeaperiodic}
Fix \(\delta\in(0,1)\). Let \(D_N\) be real nonnegative
doubly stochastic matrices, and put
\[
 B=D_N-J_N/N,\qquad \|B\|\leq1-\delta,\qquad
 \max D_{ij}\leq A_N/N,\qquad A_N\geq2.
\]
If \(A_N^2=o(\log N)\), then
\[
 \operatorname{cyc}D_N
 =(1+o(1))\frac{(N-1)!}{N^N}
       \frac{\det(I-B)}{\sqrt{\det(I-B^\top B)}}.
\]
The relative error is
\(O_\delta(N^{-1/9}\exp(C_\delta A_N^2))\), and hence
tends to zero along every sequence with the stated entry
growth and fixed gap. No positive-entry lower bound is assumed.
\end{theorem}
\begin{proof}
For all sufficiently large \(N\), \(A_N\leq\log N\).
The permanent and compensated-minor estimates of
Section~\ref{sec:relativepermanents} therefore apply
with a fixed polylogarithmic exponent.
Put \(L=\log(2N)\), and split Liu's sum at \(r=L^3\),
with integer rounding at this cutoff.

For \(r\leq L^3\), the minor ratio is
\(1+O_\delta(N^{-1/9})\), uniformly in \(U\).
Hadamard's inequality gives
\[
 \sum_{|U|=r}|\det D[U]|
   \leq\binom Nr\left(\frac{A\sqrt r}{N}\right)^r
   \leq(eA/\sqrt r)^r .
\]
Consequently
\[
 \sum_{r\geq1}r\sum_{|U|=r}|\det D[U]|
                         \leq\exp(C A^2).
\]
For example split at a sufficiently large constant times
\(A^2\); before that cutoff each logarithmic summand is
\(O(A^2)\), and afterwards the series decreases geometrically.
The small-minor error is therefore at most
\(N^{-1/9}\exp(C A^2)\).

For \(r>L^3\), the separately retained permanent-ratio bound
above yields total absolute tail at most
\[
 \exp(C A L)\sum_{r>L^3}
                    r(e^2A/\sqrt r)^r
       \leq\exp(-c L^3\log L),
\]
for large \(N\), since \(A\leq\log N\).
The same estimate without the permanent ratio controls the
tail of the determinant-only sum.
Thus the entire signed leading expression can be evaluated:
\[
 \frac{N\operatorname{cyc}D}{\operatorname{per}D}
 =\sum_{\emptyset\ne U}|U|(-1)^{|U|-1}\det D[U]
    +O_\delta\left(N^{-1/9}e^{CA^2}
                         +e^{-cL^3\log L}\right)
 =\det(I-B)+\text{the stated error}.
\]
For the last equality,
\(\det(I-tD)=(1-t)\det(I-tB)\), using both stochastic row
and column equations; differentiate at \(t=1\).
No determinant contribution has been dropped termwise.

We check the relative normalization.
The convergent trace expansion gives
\[
 \log\det(I-B)
       =-\operatorname{tr}B-\sum_{j\geq2}
                               \operatorname{tr}(B^j)/j .
\]
Here \(|\operatorname{tr}B|\leq A+1\), and
\[
 |\operatorname{tr}(B^j)|
       \leq\|B\|^{j-2}\|B\|_F^2
       \leq(1-\delta)^{j-2}A .
\]
Therefore \(\det(I-B)\geq\exp(-C_\delta A)>0\).
The normalized error is at most
\[
 O_\delta\left(
 N^{-1/9}\exp(C A^2+C_\delta A)
              +\exp(-cL^3\log L+C_\delta A)\right)=o(1).
\]
Apply Theorem~\ref{thm:growingpermanent} to the remaining full
permanent. Its determinant and factorial factors give the
displayed formula.
This extends the fixed-entry conclusion
of~\cite[Proposition 5.1]{DW} by an explicit growing-parameter
argument; it does not claim the underlying identity as new.
\end{proof}

\begin{theorem}[Relative bipartite universality]
\label{thm:relativebipartite}
Let \(T=T_m\) be real nonnegative doubly stochastic, with
\[
 B=T-J_m/m,\qquad \|B\|\leq1-\delta,\qquad
 \max T_{ij}\leq A_m/m,\qquad 2\leq A_m=o(\log m),
\]
for fixed \(\delta\in(0,1)\).
For the symmetric bipartite kernel
\(P=\begin{pmatrix}0&T\\T^\top&0\end{pmatrix}\),
\[
 Z_{\rm H}(P)
 =(1+o(1))\frac{(m!)^2}{2m\,m^{2m}}
 =(1+o(1))\pi e^{-2m}.
\]
The relative error is
\(O_\delta(m^{-1/9}\exp(C_\delta A_m))\).
In particular the leading relative constant is independent
of the nontrivial singular spectrum of \(T\).
The exact parity eigenvalue is treated separately, not inserted
into a formula requiring norm strictly below one.
\end{theorem}
\begin{proof}
For \(m\geq2\), Liu's identity for \(P\) has contributions
only when the selected set contains equally many vertices
from the two parts.
Writing those subsets as \(I,J\), each of size \(r\), exactly
\[
 \det P[I\cup J]=(-1)^r\det T[I,J]^2,\qquad
 \operatorname{per}P[(I\cup J)^c]
                         =\operatorname{per}T[I^c,J^c]^2.
\]
Consequently the complete signed expression is
\[
 2m Z_{\rm H}(P)
 =\sum_{r=1}^m(-1)^{r+1}r
     \sum_{\substack{|I|=|J|=r}}
       \det T[I,J]^2\,\operatorname{per}T[I^c,J^c]^2.
\]
This retains every mismatched row/column support and every
complementary permanent.
Its directed form is now the maintained theorem
\texttt{hamiltonianBipartiteCycleIdentity}:
\(m\operatorname{cyc}P\) equals the complete signed sum on the right.
The displayed unoriented normalization follows from
\(\operatorname{cyc}P=2Z_{\rm H}(P)\) for symmetric order \(2m\geq4\).
The formal proof transports both selected supports and both complements
explicitly; their equal-cardinality enumerations are independent.
Unbalanced determinant terms are proved zero, and the balanced
determinant and permanent squares are actual matrix identities.
The empty supports remain in the code with their zero cardinality
coefficient. No minor-ratio estimate is part of this exact identity.

The compensated-minor lemma gives a uniform
\(1+O_\delta(m^{-1/9})\) ratio for \(r\leq(\log(2m))^3\).
By Cauchy--Binet,
\[
 \sum_{|I|=|J|=r}\det T[I,J]^2=e_r(TT^\top),
\]
where \(e_r\) is the elementary symmetric polynomial in the
eigenvalues.
The total absolute mass for the small-minor error is now
bounded more sharply than the general determinant bound:
\[
 \sum_r r e_r(TT^\top)
 \leq\operatorname{tr}(TT^\top)\exp(\operatorname{tr}(TT^\top))
 \leq A e^A.
\]
The last inequality uses stochasticity,
\(\sum_{ij}T_{ij}^2\leq m\max T_{ij}\leq A\).

The retained large-minor tail is also explicit.
Hadamard and the two subset counts give
\[
 \sum_{|I|=|J|=r}\det T[I,J]^2
          \leq(e^2A^2/r)^r .
\]
Together with the squared permanent-ratio bound, its normalized
tail for \(r>(\log(2m))^3\) is at most
\[
 e^{C A\log(2m)}
       \sum_{r>(\log(2m))^3}r(e^4A^2/r)^r
          \leq e^{-c(\log(2m))^3\log\log(2m)} .
\]
Here \(A\leq\log m\) eventually.
The determinant-only tail satisfies the same smaller bound.

Now evaluate the signed leading sum, rather than bounding
it by \(Ae^A\).
Since \(TT^\top=J_m/m+BB^\top\),
\[
 \det(I-xTT^\top)=(1-x)\det(I-xBB^\top).
\]
Its derivative at \(x=1\) proves the exact cancellation
\[
 \sum_r(-1)^{r+1}r e_r(TT^\top)=\det(I-BB^\top).
\]
Thus
\[
 \frac{2m Z_{\rm H}(P)}{(\operatorname{per}T)^2}
 =\det(I-BB^\top)
      +O_\delta(m^{-1/9}Ae^A)
      +e^{-c(\log(2m))^3\log\log(2m)}.
\]
The determinant is at least \(e^{-C_\delta A}\).
The relative error therefore tends to zero because
\(A=o(\log m)\).
Finally
\[
 (\operatorname{per}T)^2
  =(1+o(1))\left(\frac{m!}{m^m}\right)^2
                               \det(I-B^\top B)^{-1}.
\]
The two determinants are equal.
Their entire spectral contribution cancels, giving the
factorial formula and then its Stirling form.
The argument shows why an absolute estimate of the signed
leading minors would miss the relative constant.
\end{proof}

\begin{corollary}[Relative counts from entropy-normalized graph weights]
\label{cor:relativegraphentropy}
For a simple balanced bipartite graph with positive stochastic
scaling \(P=\diag(u)A_G\diag(u)\), assume the hypotheses of
Theorem~\ref{thm:relativebipartite}.
Then
\[
 h(G)=(1+o(1))\pi
       \exp\left(-N-2\sum_v\log u_v\right).
\]
For a simple graph of either parity with stochastic scaling
satisfying Theorem~\ref{thm:relativeaperiodic}, its relative
count is
\[
 h(G)=(1+o(1))\frac{(N-1)!}{2N^N}
  \exp\left(-2\sum_v\log u_v\right)
           \sqrt{\frac{\det(I-B)}{\det(I+B)}} ,
 \qquad B=P-J_N/N .
\]
Both formulas allow zero support. The second uses symmetry
to simplify the determinant ratio.
\end{corollary}
\begin{proof}
Every cycle weight is multiplied by \(\prod_vu_v^2\).
For symmetric \(B\),
\(I-B^\top B=(I-B)(I+B)\), and both factors are positive
definite on their nontrivial spaces under the stated gap.
Substitute the two relative theorems, dividing directed
cycle counts by two only in the aperiodic formula.
\end{proof}

\begin{remark}[The density ranges are part of the result]
The relative aperiodic estimate here requires \(A^2=o(\log N)\);
the bipartite estimate permits \(A=o(\log N)\).
They are stronger precision statements on a narrower domain
than the earlier diffuse-kernel entropy theorem.
These estimates alone do not apply to the prime-sum
normalization where \(N\max P_{ij}\) is of order \(\log N\).
The signed first-correction argument in
Section~\ref{sec:signedjetcomparison} closes that boundary
under its explicit determinant-amplification condition.
They do cover all unitary Cayley graphs, whose corresponding
parameter is \(O(\log\log N)\), as proved next.
\end{remark}

\section{Full relative asymptotics for coprime circles}
\label{sec:relativecoprime}
The growing-entry comparison closes the arithmetic count at
relative precision. This is stronger than an exponential-base
estimate and treats every modulus, with the parity effect explicit.

\begin{lemma}[An elementary sufficiently uniform totient bound]
\label{lem:totientdiffusion}
For all \(n\geq3\),
\[
 n/\varphi(n)\ll\log\log(3n).
\]
\end{lemma}
\begin{proof}
We recall an elementary argument, not claiming this classical
estimate as new.
Chebyshev's upper bound \(\vartheta(x)\ll x\) and partial
summation give
\(\sum_{p\leq x}(\log p)/p\ll1+\log x\).
Take \(s=1+1/\log x\).
Expanding the two logarithmic Euler products, and using
\(1-e^{-u}\leq u\), gives
\[
 \log\prod_{p\leq x}
       \frac{1-p^{-s}}{1-p^{-1}}
 \leq(s-1)\sum_{p\leq x}\frac{\log p}{p-1}\ll1 .
\]
Thus
\[
 \prod_{p\leq x}(1-1/p)^{-1}
                  \ll\zeta(s)\ll\log x,
\]
using \(\zeta(s)\leq1+1/(s-1)\).
Set \(y=\log(3n)\).
There are at most \((\log n)/\log y\) prime divisors of
\(n\) exceeding \(y\), whose contribution to the logarithm
of the product is at most
\(2\log n/(y\log y)=O(1)\).
The small-prime contribution is \(O(\log y)\).
This proves the bound for large \(n\); a fixed constant
covers the remaining finite range.
No prime number theorem or RH hypothesis is used.
\end{proof}

\paragraph{Maintained full-domain density input.}
Lemma~\ref{lem:totientdiffusion} is now formalized for every
integer \(n\geq3\), with one positive absolute constant.
\texttt{HamiltonianMertensProduct} consumes the existing
cached classical Mertens product asymptotic and derives a
positive reciprocal bound.
\texttt{HamiltonianTotientEuler} proves the actual totient
Euler product and applies the previously proved large-prime
product estimate to the same prime divisors.
\texttt{HamiltonianTotientCutoff} retains both supported
products and proves, for every real \(y\geq4\),
\[
 \frac n{\varphi(n)}
 \leq
 \frac{\exp(2\log n/(y\log y))}
      {\prod_{p\leq y}(1-1/p)}.
\]
\texttt{HamiltonianTotientDiffusion} chooses
\(y=\log(3n)\). Above the fixed Mertens and logarithmic
cutoffs its tail factor is at most \(e^2\).
One constant covers all smaller \(n\geq3\) using
\(\log\log(3n)\geq\log\log3>0\) and \(n/\varphi(n)\leq n\).
The maintained theorem has the exact quantifiers
\(\exists C>0\ \forall n\geq3\).
The Mertens input is unconditional and its new consumer has
passed a foundational axiom audit; no dependency rebuild or
numerical range test is used.
This closes the density lemma.

\paragraph{Maintained quantitative odd-modulus count.}
\texttt{HamiltonianTotientSublog} proves, for every \(b>0\),
one threshold above which
\((n/\varphi(n))^2\leq b\log n\), uniformly over every modulus.
For every fixed \(K>0\), \texttt{HamiltonianCoprimeCountBudget}
supplies the simultaneous entry cap \(n/\varphi(n)\leq\log n\)
and the full loss budget
\(K(n/\varphi(n))^2\leq(\log n)/90\).
Thus the complete error
\(17n^{-1/9}\exp(K(n/\varphi(n))^2)\)
is at most \(17n^{-1/10}\).

\texttt{HamiltonianCoprimeAperiodicCount} applies these bounds
and the proved physical half-gap to the actual native
coprime kernel \(P_n=A_{X_n}/\varphi(n)\).
Writing \(B_n=P_n-J_n/n\) and
\[
 R_n=\frac{(n-1)!}{n^n}\det(I-B_n)
                      \det(I-B_nB_n^*)^{-1/2},
\]
the compiled theorem states that one \(N_0\geq256\) satisfies,
for every odd \(n\geq N_0\),
\[
 R_n>0,\qquad
 \left|\frac{Z_{\rm dir}(P_n)}{R_n}-1\right|
                      \leq17n^{-1/10}.
\]
Here \(Z_{\rm dir}\) is the actual single-cycle permutation sum.
The norm transfer between the native residue vertices and
the finite counting coordinates uses the proved norm-preserving
reindexing theorem explicitly.
Every new declaration passes the focused foundational audit.
This closes the actual odd stochastic count, including its
rate and positive reference.

\paragraph{Maintained full odd arithmetic normalization.}
\texttt{HamiltonianCoprimeConstant} proves positivity on
every supported divisor and identifies the paper's exact
real-power product \(\mathfrak C(n)\) with the positive
square root of the determinant quotient. It is therefore
the actual determinant--Gram factor in \(R_n\).
\texttt{HamiltonianCoprimeUnorientedCount} restores the
edge weight and the reversal factor exactly.
Writing
\[
 H_n=\frac{(n-1)!}{2n^n}\varphi(n)^n\mathfrak C(n),
\]
its compiled theorem gives one \(N_0\geq256\) such that,
for every odd \(n\geq N_0\),
\[
 H_n>0,\qquad
 \left|\frac{h(X_n)}{H_n}-1\right|\leq17n^{-1/10}.
\]
This is the displayed factorial normalization of the actual
unoriented arithmetic graph, with no restriction on its
distinct prime factors or prime-power exponents.
All fourteen new declarations pass the focused foundational
audit. The constant's infinite-series analysis and uniform
bounds, the final Stirling normalization and the even-modulus
count remain formalization obligations.

\paragraph{Maintained exact arithmetic determinant.}
\texttt{HamiltonianFourierDeterminant} proves that the determinant
of the complete cyclic difference operator is the product of its
full Fourier symbol. Both frame determinants cancel the exact
grid normalization.
\texttt{HamiltonianCoprimeDeterminant} applies this to the
actual compensated coprime kernel, for every positive modulus
and every real scalar \(z\):
\[
 \det(I-zB_n)
   =\prod_{t\mid n}F_z(t)^{\varphi(t)},\qquad
 F_z(1)=1,\qquad
 F_z(t)=1-z\frac{\mu(t)}{\varphi(t)}\quad(t>1).
\]
The constant mode is compensated before the determinant is taken.
The same formula retains the parity mode at even modulus and all
nonsquarefree divisors. The frequency product is regrouped using
the proved exact conductor multiplicities on the actual divisor
support. All ten new declarations pass the focused foundational
audit. Thus both determinants in the symmetric quotient are
explicit finite arithmetic products. The positive square-root
identification with the displayed \(\mathfrak C(n)\) and the
actual odd unoriented factorial count are now compiled as
recorded above. Its infinite series and bounds, final Stirling
normalization and the even count remain separate obligations.


For odd \(n\geq3\), define the positive arithmetic factor
\[
 {\mathfrak C}(n)=
 \prod_{\substack{t\mid n\\t>1}}
 \left(\frac{\varphi(t)-\mu(t)}
            {\varphi(t)+\mu(t)}\right)^{\varphi(t)/2}.
\]
Every denominator is positive, since an odd \(t>1\) has
\(\varphi(t)\geq2\).
Factors with \(\mu(t)=0\) equal one.

\begin{theorem}[Relative enumeration through every modulus]
\label{thm:relativecoprime}
As \(n\to\infty\), the number of unoriented coprime circles is
\[
 h(X_n)=
 \begin{cases}
 (1+O(n^{-1/10}))\,\pi\,(\varphi(n)/e)^n,&n\ {\rm even},\\[2mm]
 (1+O(n^{-1/10}))\,{\mathfrak C}(n)
       \sqrt{\dfrac{\pi}{2n}}\,(\varphi(n)/e)^n,
                                      &n\ {\rm odd}.
 \end{cases}
\]
The implied constant is absolute for sufficiently large \(n\).
The estimate holds along all moduli of the indicated parity,
without restrictions on their prime factors.
For odd moduli the equivalent factorial form is
\[
 h(X_n)=(1+O(n^{-1/10}))\,
       \frac{(n-1)!}{2n^n}\,\varphi(n)^n\,{\mathfrak C}(n).
\]
\end{theorem}
\begin{proof}
Use the exact spectrum proved in Theorem~\ref{thm:coprimecircles}.
The normalized adjacency \(P=A_{X_n}/\varphi(n)\) has a
fixed gap after removing the constant vector and, only in
the even case, its exact parity vector.
The largest entry is \(1/\varphi(n)\).
By Lemma~\ref{lem:totientdiffusion}, its dimension times this
entry is \(O(\log\log n)\).

If \(n=2m\), its cross block \(T\) is doubly stochastic,
has nontrivial singular norm at most \(1/2\), and satisfies
\[
 m\max T_{ij}\leq n/\varphi(n)=O(\log\log n)=o(\log m).
\]
Theorem~\ref{thm:relativebipartite} therefore gives
\(Z_{\rm H}(P)=(1+o(1))\pi e^{-n}\).
Every Hamilton cycle has weight \(\varphi(n)^{-n}\),
which proves the even formula.

If \(n\) is odd, \(B=P-J_n/n\) has norm at most \(1/2\), and
\(A=O(\log\log n)\) satisfies \(A^2=o(\log n)\).
Apply Theorem~\ref{thm:relativeaperiodic}.
For each divisor \(t>1\) of \(n\), the eigenvalue
\(\mu(t)/\varphi(t)\) has multiplicity \(\varphi(t)\).
The determinant ratio for symmetric \(B\) is exactly
\[
 \frac{\det(I-B)}{\sqrt{\det(I-B^2)}}
   =\prod_{\substack{t\mid n\\t>1}}
       \left(\frac{1-\mu(t)/\varphi(t)}
                  {1+\mu(t)/\varphi(t)}\right)^{\varphi(t)/2}
   ={\mathfrak C}(n).
\]
Divide the directed count by two and restore the common
edge-weight factor \(\varphi(n)^n\).
This proves the factorial formula, and Stirling proves the
displayed odd-order normalization.
Choose, for example, the common gap parameter \(\delta=1/3\).
The relative comparison errors are at most
\(O(n^{-1/9}\exp(C(\log\log(3n))^2))\), uniformly over
all moduli. Since \((\log\log(3n))^2=o(\log n)\), the
exponential factor is at most \(n^{1/90}\) eventually.
The permanent error \(O(n^{-1/8})\) and Stirling error
\(O(n^{-1})\) are smaller. Thus the asserted common
\(O(n^{-1/10})\) error follows, with an absolute constant;
no numerical threshold is asserted.
There is no unproved prime-sampling estimate in this consumer.
\end{proof}

\begin{proposition}[The prime-factor dependence of the odd constant]
\label{prop:coprimeconstant}
The factor \({\mathfrak C}(n)\) depends only on the distinct
prime divisors of odd \(n\), and
\[
 \log{\mathfrak C}(n)
 =1+\sum_{j\geq1}\frac{
       1-\prod_{p\mid n}(1-(p-1)^{-2j})}{2j+1}.
\]
The series converges absolutely and uniformly over all odd moduli.
In particular
\[
 e\leq{\mathfrak C}(n)\leq e^{11/9}.
\]
It increases when new distinct prime factors are added.
Along odd primorials it converges to the finite constant
obtained by taking the product over all odd primes in the
series above.
\end{proposition}
\begin{proof}
Only squarefree divisors contribute.
Expand \(\tfrac12\log((1-x)/(1+x))
=-\sum_{j\geq0}x^{2j+1}/(2j+1)\), valid for
\(|x|\leq1/2\).
The \(j=0\) sum is
\(-\sum_{t\mid n,t>1}\mu(t)=1\).
For \(j\geq1\), the multiplicative divisor sum is
\[
 \sum_{t\mid n}\frac{\mu(t)}{\varphi(t)^{2j}}
       =\prod_{p\mid n}(1-(p-1)^{-2j}),
\]
giving the displayed series.

Its numerator is nonnegative and at most
\[
 \sum_{\substack{p\ {\rm odd}\\p\ {\rm prime}}}(p-1)^{-2j}
   \leq\sum_{\ell\geq1}(2\ell)^{-2j}
   =4^{-j}\zeta(2j)\leq2\cdot4^{-j}.
\]
Here the elementary integral bound gives \(\zeta(2j)\leq2\).
This majorant proves uniform absolute convergence.
Moreover
\[
 0\leq\log{\mathfrak C}(n)-1
   \leq\frac23\sum_{j\geq1}4^{-j}=2/9.
\]
Adding a prime decreases each finite product inside the
series, and therefore increases the constant.
The same majorant justifies the primorial limit and the
absolutely convergent infinite-prime expression.
\end{proof}

\begin{remark}[Exact calibration families and what has been settled]
For odd prime powers \(n=p^a\),
\[
 {\mathfrak C}(n)
       =\left(\frac p{p-2}\right)^{(p-1)/2};
\]
in particular \({\mathfrak C}(3^a)=3\).
For prime \(n\), the formula agrees asymptotically with
the exact complete-graph count \((n-1)!/2\).
For powers of two, it agrees with the exact complete
bipartite count \((n/2)!(n/2-1)!/2\).
These are symbolic consistency checks on infinite families,
not finite experimental certificates.

The result determines the full relative asymptotic, including
its parity-dependent polynomial factor and arithmetic constant.
It does not merely count the already known constant-step
Hamilton cycles or reprove their existence.
The infinite-prime interaction appears in an explicit,
uniformly convergent expression for the odd constant.
Priority remains subject to primary-literature comparison.
\end{remark}

\section{Cyclic layers: a signed comparison beyond the bipartite case}
\label{sec:periodiccomparison}
The bipartite cancellation is one case of a finite cyclic-layer
identity. Here the layers are part of the observable; they are
not discarded by replacing the kernel with a norm envelope.

\begin{theorem}[Relative comparison for a fixed number of cyclic layers]
\label{thm:periodiccomparison}
Fix an integer \(k\geq2\) and \(\delta\in(0,1)\).
For \(i\) modulo \(k\), let \(T_i\) be real nonnegative doubly
stochastic matrices of order \(m\), with
\[
 B_i=T_i-J_m/m,\qquad \|B_i\|\leq1-\delta,\qquad
 \max T_i\leq A/m,\qquad 2\leq A=o(\log m).
\]
Let \(P\) have block \(T_i\) from layer \(i\) to layer \(i+1\)
and all other blocks zero. Put \(C=B_1B_2\cdots B_k\).
Then
\[
 \operatorname{cyc}P
 =\frac{(m!)^k}{m\,m^{km}}\,
   \frac{\det(I-C)}
        {\prod_{i=1}^k\sqrt{\det(I-B_i^\top B_i)}}\,
   \left(1+O_{\delta,k}
       \left(m^{-1/9}\exp(C_{\delta,k}A)\right)\right).
\]
The positive determinant factors are retained in full.
No commutativity, common eigenbasis, or matching row and
column supports is assumed.
\end{theorem}
\begin{proof}
In Liu's exact identity a selected principal submatrix has
zero determinant unless its selected cardinality is the
same, say \(r\), in every layer.
For selected sets \(I_1,\ldots,I_k\), each of size \(r\),
the determinant is exactly
\[
 (-1)^{(k-1)r}\prod_{i=1}^k
                    \det T_i[I_i,I_{i+1}],\qquad I_{k+1}=I_1.
\]
The complementary permanent is the product of the \(k\)
corresponding complementary permanents.
Combining the sign with \((-1)^{kr-1}\) and the coefficient
\(kr\), and dividing by \(k\), gives
\[
 m\operatorname{cyc}P
 =\sum_{r=1}^m(-1)^{r-1}r
  \sum_{\substack{I_1,\ldots,I_k\subseteq[1,m]\\|I_i|=r}}
  \prod_{i=1}^k
   \left(\det T_i[I_i,I_{i+1}]\,
        \operatorname{per}T_i[I_i^c,I_{i+1}^c]\right).
 \tag{20}
\]
This identity includes every compensation term on mismatched
supports.

Normalize by \(\prod_i\operatorname{per}T_i\).
For \(r\leq(\log(2m))^3\), each complementary-permanent
ratio is \(1+O_\delta(m^{-1/9})\), by
Lemma~\ref{lem:relativecompensatedminors}.
Their product has the same error with a constant depending on \(k\).
To bound only that error, let \(M_i\) be the matrix indexed
by \(r\)-subsets, with entries
\(|\det T_i[I,J]|\). Cauchy--Binet gives
\[
 \|M_i\|_F^2=e_r(T_iT_i^\top)\leq A^r/r!.
\]
Repeated Hilbert--Schmidt inequalities, for \(k\geq2\), imply
\[
 \sum_{I_1,\ldots,I_k}\prod_i|\det T_i[I_i,I_{i+1}]|
   =\operatorname{tr}(M_1\cdots M_k)
   \leq\prod_i\|M_i\|_F
   \leq(A^r/r!)^{k/2}.
\]
In particular the total error mass is at most
\[
 \sum_{r\geq1}r(A^r/r!)^{k/2}
               \leq(Ae^A)^{k/2}\leq\exp(C_kA).
\]
For the middle inequality, use \(k/2\geq1\) and
\(\sum_r rA^r/r!=Ae^A\).
Absolute values have been used only after separating the
minor-ratio error.

For the retained tail \(r>(\log(2m))^3\), the full product
of permanent ratios is at most \(\exp(kr+C_kA\log(2m))\).
The same Hilbert--Schmidt bound and \(r!\geq(r/e)^r\)
therefore give a total tail at most
\[
 e^{C_kA\log(2m)}
 \sum_{r>(\log(2m))^3}r
          (e^3A/r)^{kr/2}
 \leq \exp(-c_k(\log(2m))^3\log\log(2m)).
\]
The factor \(e^3\) retains the full \(e^{kr}\) permanent-ratio
cost. The determinant-only tail is smaller.

The retained signed main term is evaluated by the
Cauchy--Binet formula for products:
\[
 \sum_{I_1,\ldots,I_k}\prod_i\det T_i[I_i,I_{i+1}]
                       =e_r(T_1\cdots T_k).
\]
Since \(T_1\cdots T_k=J_m/m+C\),
\[
 \det(I-xT_1\cdots T_k)=(1-x)\det(I-xC).
\]
Differentiating at \(x=1\) now gives
\[
 \frac{m\operatorname{cyc}P}{\prod_i\operatorname{per}T_i}
 =\det(I-C)+
 O_{\delta,k}\left(m^{-1/9}e^{C_kA}
       +e^{-c_k(\log(2m))^3\log\log(2m)}\right).
\]
This is the signed step; bounding the main term by the
Hilbert--Schmidt mass would not prove the formula.

Finally put \(\rho=1-\delta\).
For \(j\geq1\), the two Frobenius factors and the remaining
operator factors give
\[
 |\operatorname{tr}(C^j)|\leq A\rho^{kj-2}.
\]
The logarithmic determinant series is absolutely convergent,
so \(|\log\det(I-C)|\leq C_{\delta,k}A\).
Its determinant is positive by continuity from \(I\) and
\(\|C\|<1\). Dividing by it costs at most \(e^{C_{\delta,k}A}\).
Apply the growing-entry permanent theorem to every \(T_i\).
All full determinants remain in the displayed result,
and all relative errors fit the stated bound.
\end{proof}

\begin{corollary}[Why two layers are special]
\label{cor:twolayersspecial}
If \(k=2\) and \(T_2=T_1^\top\), the spectral factor in
Theorem~\ref{thm:periodiccomparison} is exactly one.
For \(k\geq3\), if all cross blocks equal a symmetric matrix
\(T=J_m/m+B\), that factor is
\[
 \frac{\det(I-B^k)}{\det(I-B^2)^{k/2}}.
\]
Thus higher layers retain a genuine joint spectral contribution.
\end{corollary}
\begin{proof}
For two layers \(C=B_1B_1^\top\); the numerator cancels
the product of the two square-root determinants.
For equal symmetric blocks \(C=B^k\), while each denominator
is \(\sqrt{\det(I-B^2)}\).
These are exact finite identities, preceding any limiting estimate.
\end{proof}

\begin{remark}
The new comparison range is \(A=o(\log m)\), with fixed \(k\).
There is no claim of uniformity when the number of layers grows.
The proof neither invokes an aperiodic hypothesis on the full
\(km\)-vertex matrix nor removes its \(k\) exact unit-circle modes.
They are encoded by the cyclic block structure from the outset.
\end{remark}

\section{An arithmetic consumer with all cyclic layers and parity obstructions}
\label{sec:layeredcoprime}
For fixed \(k\geq2\), form the directed graph on
\((\mathbb Z/k\mathbb Z)\times(\mathbb Z/n\mathbb Z)\)
with arcs
\[
 (i,x)\longrightarrow(i+1,y)
             \quad\Longleftrightarrow\quad \gcd(y-x,n)=1 .
\]
Let \(H_k(n)\) count directed Hamilton cycles modulo cyclic
rotation, retaining their orientation. This convention is
important: for \(k=2\) opposite orientations are both counted.
The graph is \(\varphi(n)\)-regular on \(kn\) vertices.

For odd \(s\), including \(s=1\), put
\(\lambda_d=\mu(d)/\varphi(d)\) and define
\[
 {\cal D}_k(s)=
 \prod_{\substack{d\mid s\\d>1}}
 \left(\frac{1-\lambda_d^k}
              {(1-\lambda_d^2)^{k/2}}\right)^{\varphi(d)},
 \qquad
 {\cal E}_k(s)=
 \prod_{\substack{d\mid s\\d>1}}
 \left(\frac{1-\lambda_d^{2k}}
              {(1-\lambda_d^2)^k}\right)^{\varphi(d)} .
\]
Empty products equal one; factors with \(\mu(d)=0\) equal one.
All displayed factors are positive.

\begin{theorem}[Complete relative counts for layered coprime circles]
\label{thm:layeredcoprime}
Fix \(k\geq2\). Uniformly as \(n\to\infty\):
\begin{enumerate}
\item If \(n\) is odd, then
\[
 H_k(n)=
 \frac{(n!)^k}{n\,n^{kn}}\varphi(n)^{kn}{\cal D}_k(n)
                  \left(1+O_k(n^{-1/10})\right).
\]
Equivalently its leading prefactor is
\((2\pi)^{k/2}n^{k/2-1}{\cal D}_k(n)\)
multiplying \((\varphi(n)/e)^{kn}\).
\item If both \(n\) and \(k\) are even, then \(H_k(n)=0\) exactly.
\item If \(n=2m\) is even and \(k\) is odd, write
\(n=2^a s\) with \(s\) odd. Then
\[
 H_k(n)=
 \frac{(m!)^{2k}}{m\,m^{2km}}\varphi(n)^{kn}{\cal E}_k(s)
                   \left(1+O_k(n^{-1/10})\right).
\]
Equivalently its leading prefactor is
\((2\pi)^k(n/2)^{k-1}{\cal E}_k(s)\)
multiplying \((\varphi(n)/e)^{kn}\).
\end{enumerate}
There is no restriction on the number of distinct prime
divisors, on their sizes, or on the two-adic valuation.
The constants may depend on the fixed \(k\).
\end{theorem}
\begin{proof}
For odd \(n\), every normalized cross block is the same
symmetric doubly stochastic unitary kernel \(T\).
Its nontrivial eigenvalues are \(\lambda_d\), each with
multiplicity \(\varphi(d)\), for \(d\mid n\), \(d>1\).
They have absolute value at most \(1/2\).
The entry parameter is \(A=O(\log\log(3n))\).
Theorem~\ref{thm:periodiccomparison} applies with \(m=n\);
its full determinant quotient is precisely \({\cal D}_k(n)\).
Restoring the common edge-weight factor
\(\varphi(n)^{kn}\) proves the first formula.

When \(n\) and \(k\) are both even, every arc preserves
\(x+i\pmod2\). Both invariant classes are nonempty.
No directed Hamilton cycle can cover both, proving exact
vanishing, rather than a failure of an asymptotic estimate.

Now suppose \(n=2m\) and \(k\) is odd.
Every arc flips the parity of the coordinate \(x\).
The map on layer/parity states
\((i,\epsilon)\mapsto(i+1,1-\epsilon)\) is a single cycle
of length \(2k\). Thus the actual full graph has \(2k\)
cyclic layers, each of size \(m\).
Their cross blocks alternate between \(U\) and \(U^\top\),
where \(U\) is the normalized even-to-odd unitary block.

By the Chinese remainder identification modulo \(2^a\)
and \(s\), row and column permutations identify \(U\) with
the tensor product of the unitary kernel on
\(\mathbb Z/s\mathbb Z\) and the constant stochastic
matrix of order \(2^{a-1}\).
For \(s=1\) the former kernel is the \(1\)-by-\(1\) matrix \((1)\).
Consequently the squared nontrivial singular values of \(U\)
are \(\lambda_d^2\) with multiplicity \(\varphi(d)\),
for \(d\mid s\), \(d>1\), together with zero values.
Their maximum is at most \(1/4\), uniformly in \(a,s\).

Write \(B=U-J_m/m\). The cyclic product of the \(2k\)
centered cross blocks is \((BB^\top)^k\).
The full numerator is \(\det(I-(BB^\top)^k)\);
the product of square-root denominators is
\(\det(I-BB^\top)^k\).
Their quotient is \({\cal E}_k(s)\).
The density parameter satisfies
\[
 m\max U_{ij}=\frac{m}{\varphi(n)}
            =\frac{s}{\varphi(s)}
               =O(\log\log(3n)).
\]
Apply Theorem~\ref{thm:periodiccomparison} with \(2k\) layers
and layer size \(m\), and restore \(\varphi(n)^{kn}\).
Its error is at most
\(O_k(n^{-1/9}\exp(C_k\log\log(3n)))\), hence
\(O_k(n^{-1/10})\).
The same bound applies in the odd case.
Finally Stirling's formula proves the two equivalent
prefactor expressions. All parity modes, orientations,
and layer-size factorials have been included.
\end{proof}

\begin{proposition}[The uncancelled arithmetic interaction for higher layers]
\label{prop:layeredconstants}
For odd \(s\), define, for integers \(b\geq2\),
\[
 F_b(s)=\prod_{p\mid s}
              \left(1+\frac{(-1)^b}{(p-1)^{b-1}}\right)-1 .
\]
Then \(F_2(s)=s/\varphi(s)-1\), and exactly
\[
 \log{\cal D}_k(s)
 =\sum_{j\geq1}
       \left(\frac{k}{2j}F_{2j}(s)-\frac1jF_{kj}(s)\right),
 \qquad
 \log{\cal E}_k(s)
 =\sum_{j\geq1}
       \left(\frac{k}{j}F_{2j}(s)-\frac1jF_{2kj}(s)\right).
\]
In particular \({\cal D}_2(s)=1\) exactly.
For every fixed \(k\geq3\),
\[
 \log{\cal D}_k(s)
      =\frac k2\left(\frac{s}{\varphi(s)}-1\right)+O_k(1),
 \qquad
 \log{\cal E}_k(s)
      =k\left(\frac{s}{\varphi(s)}-1\right)+O_k(1),
\]
uniformly over all odd \(s\).
After subtracting those respective leading terms, the
remaining series converge absolutely and uniformly.
The corresponding normalized constants have finite positive
limits along odd primorials.
\end{proposition}
\begin{proof}
For \(b\geq2\), the exact multiplicative identity is
\[
 \sum_{\substack{d\mid s\\d>1}}\varphi(d)\lambda_d^b
       =\sum_{\substack{d\mid s\\d>1}}
                   \frac{\mu(d)^b}{\varphi(d)^{b-1}}
       =F_b(s).
\]
Expand each logarithm in the definitions of
\({\cal D}_k,{\cal E}_k\), retaining the numerator and
denominator series with their signs. Since every
nontrivial \(|\lambda_d|\leq1/2\), these expansions
converge for each \(s\), and give the displayed formulas.
For \(k=2\) they cancel term by term.

Uniformity after the leading term is also explicit.
For \(b\geq3\),
\[
 |F_b(s)|
 \leq\exp\left(\sum_{p\ {\rm odd}}(p-1)^{-(b-1)}\right)-1
 \leq\exp\left(2^{-(b-1)}\zeta(b-1)\right)-1
 \leq C\,2^{-b}.
\]
The constant \(C\) is absolute, as \(\zeta(b-1)\leq2\).
This gives a summable bound for every remaining term.
In \({\cal D}_k\), \(k\geq3\), the only \(b=2\) term
is \(kF_2/2\); in \({\cal E}_k\) it is \(kF_2\).
All other terms have \(b\geq3\).
Each infinite-prime product with \(b\geq3\) converges
absolutely along primorials. Dominated convergence now
proves the asserted normalized limits; exponentiation
makes them positive and finite.
\end{proof}

\begin{remark}[The mechanism and its limits]
For a single symmetric pair of cross blocks the quadratic
spectral interaction cancels exactly.
For three or more layers it survives as the explicit
\(s/\varphi(s)-1\) contribution, with a uniformly controlled
higher-order prime interaction.
Thus the method can distinguish exact cancellation from
a nonzero arithmetic constant, rather than just relocating
an unestimated signed term.
It proves relative counts on the entire modulus domain in
this theorem. It makes no assertion about integer prime
gaps, RH, or a zero-free exponent.
The fixed-layer restriction and the stated diffuseness
range remain essential hypotheses of the comparison used.
\end{remark}

\section{A permanent expansion with a second-order uniform remainder}
\label{sec:secondpermanent}
The earlier \(N^{-1/8}\) remainder was sufficient for
\(A=o(\log N)\). At prime density it is necessary to retain
the first correction before summing the signed minors.

For a stochastic matrix \(D=J_N/N+B\), define
\[
 \Omega=\begin{pmatrix}0&B\\B^\top&0\end{pmatrix},\quad
 \Gamma=\Omega(I+\Omega)^{-1},\quad
 \Delta(s)=\det(I-s^2B^\top B)^{-1/2},
\]
and
\[
 \Psi(\Gamma)=
 -\frac34\sum_i\Gamma_{ii}^2
 +\frac12\sum_{i,j}\Gamma_{ii}\Gamma_{jj}\Gamma_{ij}
 +\frac13\sum_{i,j}\Gamma_{ij}^3,\qquad
 {\cal Q}_N(D)=\frac{\Delta''(1)}{2(N-1)\Delta(1)}
                                     +\Psi(\Gamma).
\]
The indices in \(\Psi\) range over the full \(2N\)-coordinate
block matrix, including all diagonal terms.

\begin{theorem}[First correction with a uniform two-power remainder]
\label{thm:secondpermanent}
Under Theorem~\ref{thm:growingpermanent}, for sufficiently large
\(N\) depending on \(\delta,K\),
\[
 \operatorname{per}D
 =\frac{N!}{N^N}\Delta(1)
   \left(1+{\cal Q}_N(D)
             +O_\delta((1+A)^{40}/N^2)\right),
 \qquad |{\cal Q}_N(D)|\leq C_\delta(1+A)^2/N .
\]
In particular the first correction is explicit and is not
hidden in a relative \(O(N^{-1/8})\) error.
\end{theorem}
\begin{proof}
Use the exact normalized Gaussian rook representation in
Section~\ref{sec:relativepermanents}. Its vertex factor has
expansion
\[
 r(z)=1+\tfrac13z^3-\tfrac14z^4+\tfrac15z^5
                              -\tfrac19z^6+O(z^7).
\]
Write \(\Gamma_s=s\Omega(I+s\Omega)^{-1}\).
The diagram with one degree-four vertex contributes
\(-3\sum_i(\Gamma_s)_{ii}^2/4\).
Two distinct degree-three vertices contribute
\[
 \sum_{i<j}
 \left((\Gamma_s)_{ii}(\Gamma_s)_{jj}(\Gamma_s)_{ij}
                              +\tfrac23(\Gamma_s)_{ij}^3\right).
\]
These coefficients follow from the three pairings of four
Gaussian variables and the \(9+6\) pairings of two triples.
Replacing \(i<j\) by half the ordered sum changes the result
by \(O_\delta(A^3/N^2)\), since
\(\max|(\Gamma_s)_{ii}|\leq C_\delta A/N\).
The one-vertex degree-six term also has that smaller order.

Here is a bound on all remaining diagrams, including their
infinite power-series tails.
In the finite Wick estimate already displayed, \(e\) is the
number of pairs and \(v\leq2e/3\).
For \(4\leq e\leq11\), one has \(e-v\geq2\);
their total is \(O_\delta((1+A)^{11}/N^2)\), after summing
the geometric coefficient convolution.
For \(e\geq12\), use the same cutoff
\(H=\lceil C_\delta A^3\sqrt N\rceil\) and sum
\[
 \sum_{e\geq12}e(C_\delta A^2N^{-1/6})^e
                    =O_\delta((1+A)^{24}/N^2).
\]
The cases \(e=1,2,3\) are exactly the terms retained above,
the zero odd Gaussian moments, and the smaller degree-six term.
The Cauchy tail past \(H\) is exponentially small.
This also applies to the retained degree-four and two-triple
contributions themselves: their covariance power series are
analytic on the same fixed circle \(|t|=R\), where the
resolvent bounds above bound their suprema by a polynomial
in \(1+A\). Hence extending their truncated coefficient sums
to their full values introduces only an exponentially small tail.
Thus, uniformly on the fixed interval \(s\in[a,b]\),
\[
 R_B(s)=\Delta(s)
        \left(1+\Psi(\Gamma_s)
                  +O_\delta((1+A)^{40}/N^2)\right).
\]
No coefficientwise asymptotic with an unspecified entry-bound
constant is used.

The row and column squared norms give
\[
 \|\Gamma_s\|\leq C_\delta,\quad
 \|\Gamma_s\|_F^2\leq C_\delta A,\quad
 \max_{ij}|(\Gamma_s)_{ij}|\leq C_\delta A/N.
\]
The same bounds hold for any fixed number of \(s\)-derivatives.
For \(g_i=(\Gamma_s)_{ii}\), the middle term in \(\Psi\) is
\(g^\top\Gamma_s g/2\), bounded by \(C_\delta A^2/N\).
The cubic entry sum is bounded by
\(\max|\Gamma_s|\|\Gamma_s\|_F^2\).
This proves \(|\Psi(\Gamma_s)|\leq C_\delta A^2/N\),
also for the derivatives needed below.

It remains to perform the Gamma average with its first
correction retained. Put \(S=N/T\), \(T\sim\Gamma(N+1,1)\),
\(I=[a,b]\), and
\(C_N=\{|S-1|\leq N^{-1/4}\}\subseteq\{S\in I\}\).
All expectations of \(\Delta(S)\) or
\(\Delta(S)\Psi(\Gamma_S)\) below are restricted to \(S\in I\).
No extension of these local analytic factors outside \(I\)
is needed.

For integers \(0\leq j\leq8\) and \(N\geq16\), direct
integration of the Gamma density gives
\[
 \mathbb E S^j=\prod_{\ell=0}^{j-1}\frac{N}{N-\ell}
 \leq2^j .
\]
In particular the exact full centered moments are
\[
 \mu_1=0,\qquad \mu_2=\frac1{N-1},\qquad
 \mu_3=\frac4{(N-1)(N-2)},\qquad
 \mu_4=\frac{3(N+6)}{(N-1)(N-2)(N-3)} .
\]
More explicitly, if \(0<\varepsilon\leq1/2\) and
\(N\varepsilon\geq4\), both actual Gamma tails give
\[
 \mathbb P(|S-1|>\varepsilon)
 \leq2\exp(-(N+1)\varepsilon^2/36).
\]
For \(4k\leq N\), the raw moments imply
\(\mathbb E|S-1|^{2k}\leq2^{2k}+1\).
Put \(u=\exp(-(N+1)\varepsilon^2/72)\) and
\(E_\varepsilon=\{|S-1|>\varepsilon\}\).
The elementary quadratic inequality
\[
 u|S-1|^k\one_{E_\varepsilon}
 \leq u^2|S-1|^{2k}+\one_{E_\varepsilon}
\]
therefore proves the explicit bound
\[
 \mathbb E[|S-1|^k\one_{E_\varepsilon}]
 \leq(2^{2k}+3)\exp(-(N+1)\varepsilon^2/72).
\]
Taking \(\varepsilon=N^{-1/4}\) gives the required
\(C_k\exp(-\sqrt N/72)\) bound for \(0\leq k\leq4\).
Thus each truncated signed moment on \(C_N\) equals
\(\mu_k+O(\exp(-c'\sqrt N))\), including the first
and third moments. In particular the truncated first moment
is not set to zero without its error.

For clarity, the precise local averaging estimate is as follows.
Let \(C_\varepsilon=\{|S-1|\leq\varepsilon\}\), \(M\geq0\),
and let \(f\) be measurable with
\[
 |f(s)-a-b(s-1)-c(s-1)^2-d(s-1)^3|
 \leq M(s-1)^4\quad (|s-1|\leq\varepsilon).
\]
For \(N\geq16\) and the preceding range of \(\varepsilon\),
the actual normalized Euler integral satisfies
\[
 \begin{split}
 &\left|\mathbb E[\one_{C_\varepsilon}f(S)]
       -a-\frac{c}{N-1}-\frac{4d}{(N-1)(N-2)}\right|\\
 &\quad\leq M\mu_4+
 e^{-(N+1)\varepsilon^2/72}
       (4|a|+7|b|+19|c|+67|d|).
 \end{split}
\]
Indeed the full cubic has exactly the displayed signed average;
only its discarded tail is bounded absolutely.
The local remainder is bounded by \(M\mu_4\).
These bounds also prove integrability on the window, without
extending the local approximation to other values of \(S\).

On \(C_N\), fixed-order derivatives of
\(\Delta(s)/\Delta(1)\) are bounded by polynomials in \(1+A\).
Indeed its logarithmic derivatives have that property and
\(\Delta(s)/\Delta(1)\leq\exp(C_\delta A N^{-1/4})\)
is bounded for sufficiently large \(N\).
Taylor's formula to third order, with a fourth-derivative
remainder and the truncated moments just proved, gives
\[
 \mathbb E\bigl[\one_{\{S\in I\}}\Delta(S)\bigr]
 =\Delta(1)\left(
  1+\frac{\Delta''(1)}{2(N-1)\Delta(1)}
                  +O_\delta((1+A)^8/N^2)\right).
\]
Here the cubic term uses the signed \(\mu_3=O(N^{-2})\);
the fourth-order remainder uses \(\mu_4=O(N^{-2})\).
The contribution of \(I\setminus C_N\) is at most
\(\Delta(1)\exp(C_\delta A-c'\sqrt N)\), which is
smaller than any fixed inverse power of \(N\) throughout
\(A\leq(\log N)^K\).
All implicit constants in these averages depend only on \(\delta\);
\(K\) enters only the sufficiently-large-\(N\) threshold.

For \(f(s)=\Delta(s)\Psi(\Gamma_s)\), the previous covariance
and derivative bounds give
\[
 \sup_{C_N}\frac{|f''(s)|}{\Delta(1)}
       \leq C_\delta(1+A)^6/N .
\]
Taylor's formula through the linear term, the exponentially
small truncated first moment, and
\(\mathbb E(S-1)^2=1/(N-1)\) therefore give
\[
 \mathbb E\bigl[\one_{\{S\in I\}}\Delta(S)\Psi(\Gamma_S)\bigr]
 =\Delta(1)\Psi(\Gamma)
       +O_\delta\bigl(\Delta(1)(1+A)^{12}/N^2\bigr).
\]
The portion in \(I\setminus C_N\) is again exponentially
small. The uniform rook-expansion remainder averages to
\(O_\delta(\Delta(1)(1+A)^{40}/N^2)\).
On \(S\notin I\) we retain the original polynomial \(R_B(S)\);
the preceding outside-interval estimate gives
\[
 \left|\mathbb E\bigl[\one_{\{S\notin I\}}R_B(S)\bigr]\right|
 \leq\exp(-c_\delta N+C_\delta\sqrt{NA}).
\]
Since \(\Delta(1)\geq1\), this too is negligible relative
to the claimed error. Finally the logarithmic derivative
formulas give
\(\Delta''(1)/\Delta(1)=O_\delta((1+A)^2)\).
Combining the restricted averages and the original polynomial
tail proves both conclusions.
\end{proof}

\section{The full compensated-minor correction}
\label{sec:compensatedjet}
Fix \(\delta\in(0,1)\) and \(K_0>0\).
Let \(D=J_N/N+B\) satisfy Theorem~\ref{thm:secondpermanent}
with entry parameter \(A\leq(\log N)^{K_0}\).
Use its block matrices \(\Omega,\Gamma\), and set
\[
 P=\begin{pmatrix}0&D\\D^\top&0\end{pmatrix},\quad
 c=\frac{(\one,\one)}{\sqrt{2N}},\quad
 s=\frac{(\one,-\one)}{\sqrt{2N}},\quad
 K=(I+\Omega)^{-1}-\tfrac12cc^\top-ss^\top .
\]
Thus \(K\) is the inverse of \(I+P\) on \(s^\perp\),
extended by zero on \(s\).
Define \(h_i=-\Gamma_{ii}\) and
\[
 v=(\tfrac12I-K)h .
\]
The two halves of \(h\) have equal sums.
All these matrices are real, and
\[
 \max|\Gamma_{ij}|\leq C_\delta A/N,\quad
 \|K\|_{\infty\to\infty}\leq C_\delta(1+A),\quad
 \|v\|_\infty\leq C_\delta(1+A)^2/N .
\]

\begin{theorem}[A uniform signed first-order minor expansion]
\label{thm:minorjet}
For arbitrary row and column sets \(I,J\) with
\(|I|=|J|=r\leq(\log(2N))^3\), put
\(u=(\one_I,\one_J)\), retaining their separate coordinates.
There is a fixed finite exponent \(M=M(\delta,K_0)\) such that,
uniformly in these sets,
\[
 \left(\frac{\operatorname{per}D[I^c,J^c]}
                   {\operatorname{per}D}\right)^2
 =1+\frac{r(r-1)}N+u^\top\Gamma u+2v^\top u
                         +O_{\delta,K_0}\left(
                              \frac{\log(2N)^M}{N^2}\right).
\]
The two-support quadratic term and the compensation field
are part of the expansion. They are not assumed to vanish.
\end{theorem}
\begin{proof}
Write \(n=N-r\), \(Q=D[I^c,J^c]\), and
\(s_{IJ}=\sum_{i\in I,j\in J}D_{ij}\).
Let \(\ell=Pu=(D\one_J,D^\top\one_I)\).
The remaining row defects are the corresponding coordinates
of \(\ell\), each at most \(rA/N\), and their sum is
\(2(r-s_{IJ})\).
Let \(\widetilde D=\diag(1+z_A)Q\diag(1+z_B)\)
be its exact stochastic repair, with \(z=0\) on removed
coordinates and gauge \(s^\top z=0\).
The established repair gives
\(\|z\|_\infty\leq C_\delta rA\log(2N)/N\).

We give the local error accounting needed for a two-power
remainder. On retained coordinates the exact row equation
is
\[
 (I+P)z=\ell+\ell\mathbin{\odot}z
                         -z\mathbin{\odot}(Pz).
\]
On removed coordinates the discrepancy is \((Pz)-\ell\).
Its magnitude is at most a polynomial in \(r,A,\log(2N)\)
divided by \(N\).
Every off-diagonal entry of \(K\) is \(O_\delta((1+A)/N)\).
Multiplying that discrepancy, supported on \(2r\) coordinates,
by \(K\), and the retained quadratic discrepancy by the
row-sum bound for \(K\), proves
\[
 z=K\ell+
 O_\delta\left(
  \frac{\operatorname{poly}(r,1+A,\log(2N))}{N^2}\right)
 \quad\hbox{on retained coordinates}.
\]
Here and below \(\operatorname{poly}\) denotes a fixed finite
polynomial with nonnegative coefficients depending only on
the displayed fixed parameters. No degree grows with \(N,r,A\).
On removed coordinates \(K\ell\) can instead be bounded by
\(C_\delta rA\log(2N)/N\), which is multiplied by an
additional \(O_\delta(A/N)\) factor whenever used below.

Sum the exact nonlinear scaling equations and expand
\(\log(1+z)\) through degree two. With the actual retained
row degrees \(d_i\), they give
\[
 \sum d_i z_i=r-s_{IJ}-\tfrac12z^\top Q_{\rm bip}z,
\]
and hence the full scaling cost is
\[
 {\cal S}:=\sum\log(1+z_A)+\sum\log(1+z_B)
 =r-s_{IJ}+\tfrac12\ell^\top K\ell
       +O_\delta\left(
       \frac{\operatorname{poly}(r,1+A,\log(2N))}{N^2}\right).
\]
Indeed the cubic Taylor cost is at most \(2N\|z\|_\infty^3\);
the omitted-coordinate costs are at most
\(2r(rA/N)\|K\ell\|_\infty\).
Both have the stated scale. The retained-coordinate approximation and
\(\|\ell\|_1=2r\) control the remaining substitutions.
This step retains the quadratic compensation, not just
the maximum scaling multiplier.

Next we compute the determinant cost locally.
Put \(M=I+\Omega\), and let \(S\) be the \(2r\) removed
coordinates. Jacobi's identity gives
\[
 \frac{\det M[S^c]}{\det M}
     =\det(M^{-1}[S])
     =\det(I-\Gamma[S]).
\]
Since \(\|\Gamma[S]\|\leq C_\delta rA/N=o(1)\),
its contribution to the logarithm of the inverse-square-root
determinant is
\[
 \tfrac12\operatorname{tr}\Gamma[S]
       +O_\delta(r^2A^2/N^2)=-\tfrac12 h^\top u+
                                      O_\delta(r^2A^2/N^2).
\]
The remaining changes are scaling and the change
\(J_n/N\) to \(J_n/n\).
The linear scaling contribution to the log determinant is
\(h^\top K\ell\).
For detail the off-diagonal block of \(M^{-1}\) is
\(-B(I-B^\top B)^{-1}\); its row and column sums vanish.
Contracting it with \(\diag(z_A)D+D\diag(z_B)\)
therefore gives \(h^\top z\), with the constant \(J_N/N\)
parts cancelling exactly.

All other determinant terms have a two-power remainder.
The Schur-complement change from \(M^{-1}[S^c]\) to
\((M[S^c])^{-1}\) has trace norm at most \(C_\delta rA/N\).
It is multiplied by a perturbation of operator norm
\(O_\delta(rA\log(2N)/N)\).
The squared Frobenius norm of the scaling perturbation is
\(O_\delta(r^2(1+A)^3\log(2N)^2/N^2)\).
Removed-coordinate contractions contain \(rN\) terms,
each with one scaling factor and two \(O_\delta(A/N)\)
matrix entries.
Finally the constant-block change contracts to zero on the
full supports; its remaining restricted contraction contains
only the removed row/column overlap and is of the same
smaller scale. The exact quadratic log-determinant remainder
is bounded by these squared Frobenius norms, since the
matrices retain a fixed positive spectral lower bound.
Consequently, writing \(\widetilde B=\widetilde D-J_n/n\),
\[
 \log\frac{\det(I-\widetilde B^\top\widetilde B)^{-1/2}}
               {\det(I-B^\top B)^{-1/2}}
 =-\tfrac12h^\top u+h^\top K\ell
 +O_\delta\left(
       \frac{\operatorname{poly}(r,1+A,\log(2N))}{N^2}\right).
\]

The first permanent correction must also be compared,
rather than independently discarded at order \(N^{-1}\).
The same Schur formula and scaling estimates imply, on
retained coordinates,
\[
 |\widetilde\Gamma_{ij}-\Gamma_{ij}|
 \leq\frac{\operatorname{poly}(r,1+A,\log(2N))}{N^2}.
\]
To check this entry bound, both full inverse matrices have
off-diagonal entries \(O_\delta(A/N)\) and absolute row
sums \(O_\delta(1+A)\). The entrywise inverse perturbation
identity therefore preserves the displayed scale.
Deleted coordinates contribute at most \(2r\) diagonal
terms or \(4rN\) off-diagonal terms.

Apply this estimate to every displayed term of \(\Psi\):
a diagonal square has difference bounded by its size
\(O(A/N)\) times the entry difference; the cubic sums have
two such small factors.
After summation their differences are
\(\operatorname{poly}(r,1+A,\log(2N))/N^2\).
For the Gamma averaging term, set
\[
 S_1=\operatorname{tr}\frac{B^\top B}{I-B^\top B},\qquad
 S_2=\operatorname{tr}
                      \left(\frac{B^\top B}{I-B^\top B}\right)^2.
\]
Then \(\Delta''(1)/\Delta(1)=S_1^2+S_1+2S_2\).
Here \(S_1\) is minus half the trace of \(\Gamma\), and
\(S_2\) is half the sum of squared entries of its two
diagonal blocks.
The preceding local bounds change each by
\(\operatorname{poly}(r,1+A,\log(2N))/N\).
The additional factor \(1/(N-1)\), including its change to
\(1/(n-1)\), proves
\[
 {\cal Q}_n(\widetilde D)-{\cal Q}_N(D)
       =O_\delta\left(
       \frac{\operatorname{poly}(r,1+A,\log(2N))}{N^2}\right).
\]
This is why the explicit first correction in the full
permanent theorem was needed.

Now use Theorem~\ref{thm:secondpermanent} in numerator
and denominator, together with the exact scaling cost.
Stirling gives
\[
 \log\frac{n!/n^n}{N!/N^N}
       =r-\frac r{2N}+O((r+r^2)/N^2).
\]
The scaling, determinant, and permanent-correction estimates yield
\[
 \log\frac{\operatorname{per}Q}{\operatorname{per}D}
 =s_{IJ}-\tfrac12\ell^\top K\ell-\frac r{2N}
           +(\tfrac12h-Kh)^\top u
           +O_{\delta,K_0}(\log(2N)^M/N^2).
\]
Finally \(u\perp s\), \(\ell=Pu\), and on \(s^\perp\)
\(PKP=P-I+K\). Since \(u^\top Pu=2s_{IJ}\),
\(u^\top u=2r\), and
\[
 u^\top Ku=2r-u^\top\Gamma u-r^2/N,
\]
the scalar and overlap terms simplify to
\[
 \frac{r(r-1)}{2N}+\tfrac12u^\top\Gamma u+v^\top u.
\]
This quantity is a fixed polynomial in \(r,A,\log(2N)\)
divided by \(N\). Exponentiating twice changes it only by
a two-power remainder and proves the displayed expansion.
Every polynomial above becomes a fixed power of \(\log(2N)\)
on \(r\leq\log(2N)^3\), \(A\leq(\log N)^{K_0}\).
\end{proof}


\begin{remark}[Independent check against every cofactor recursion]
Let \(W(I,J)=r(r-1)/(2N)+u^\top\Gamma u/2+v^\top u\)
be the first-order logarithmic correction, and put \(y=Kh\).
For \(i\notin I\), \(j\notin J\), its exact increment is
\[
 \begin{split}
 W(I\cup\{i\},J\cup\{j\})-W(I,J)
 ={}&r/N+(\Gamma u)_{A_i}+(\Gamma u)_{B_j}\\
    &+\Gamma_{A_i,B_j}-y_{A_i}-y_{B_j}.
 \end{split}
\]
Extend this polynomial expression temporarily to every \(j\),
only for the algebraic summation.
The exact matrix identities
\[
 (I+P)\Gamma=\Omega,\qquad (I+P)y=h,\qquad
 \sum_jD_{ij}\Gamma_{A_i,B_j}=h_{A_i}
\]
show that its \(D_{ij}\)-weighted sum over all \(j\) is
\[
 r/N+\bigl((I+P)\Gamma u\bigr)_{A_i}
             =r/N+(B\one_J)_i=(D\one_J)_i .
\]
Deleting the \(r\) forbidden \(j\)'s changes the sum by
only \(\operatorname{poly}(r,1+A,\log(2N))/N^2\).
Multiplying \(1+W(I,J)\) by the actual remaining row
mass \(1-(D\one_J)_i\) proves
\[
 \sum_{j\notin J}D_{ij}
       \bigl(1+W(I\cup\{i\},J\cup\{j\})\bigr)
 =1+W(I,J)+O_{\delta,K_0}(\log(2N)^M/N^2).
\]
This is precisely the first-order form of the exact
permanent cofactor recursion, for every support pair,
not just its row average. The column check is its transpose.
For \(D=J_N/N\), the exact ratio is
\(\prod_{\ell=0}^{r-1}(1-\ell/N)^{-1}\);
its square has first correction \(r(r-1)/N\), as above.
These are symbolic checks on the full quantified domain.
They do not replace the proof of the remainder.
\end{remark}

\section{Signed compensation at logarithmic density}
\label{sec:signedjetcomparison}
The decisive improvement is to sum the complete first minor
correction with its signs before bounding the two-power
remainder. In particular a polynomial correction in the
selected size is not charged its absolute determinant mass.

\begin{lemma}[Signed occupancy moments at the parity root]
\label{lem:signedoccupancy}
Let \(T=J_N/N+B\) be real nonnegative doubly stochastic with
\(\|B\|\leq\rho\), \(0\leq\rho<1\), \(\max T\leq A/N\), \(A\geq2\).
For \(|I|=|J|=r\), let \(u=(\one_I,\one_J)\), and define
\[
 {\cal L}(f)=\sum_{r=1}^N(-1)^{r+1}r
       \sum_{|I|=|J|=r}\det T[I,J]^2 f(I,J).
\]
Put \(d(t)=\det(I-tBB^\top)\), \(d=d(1)>0\).
Uniformly over coordinates \(a,b\) of \(u\),
\[
 {\cal L}(1)=d,\quad
 \frac{{\cal L}(r(r-1))}{d}
                =\frac{4d'(1)+3d''(1)}d=O_\rho((1+A)^2),
\]
\[
 \frac{|{\cal L}(u_a)|}{d}\leq C_\rho(1+A)/N,\qquad
 \frac{|{\cal L}(u_a u_b)|}{d}
          \leq C_\rho(1+A)^2/N^2\quad(a\ne b).
\]
These are signed identities and bounds; no positivity of
individual occupancy moments is asserted.
\end{lemma}
\begin{proof}
Introduce separate row variables \(x_a\) on both parts,
let \(W=\diag(x)\), and put
\[
 F(t,x)=\det\left(I+\sqrt t\,W
                   \begin{pmatrix}0&T\\T^\top&0\end{pmatrix}\right),
 \qquad g(t,x)=-t\,\partial_tF(t,x).
\]
Cauchy--Binet shows that \(g(1,x)\) is precisely the signed
sum defining \({\cal L}\), with monomial \(x^u\).
At \(x=\one\),
\[
 F(t,\one)=(1-t)d(t),\qquad
 g(t,\one)=t d(t)-t(1-t)d'(t).
\]
Thus \(g(1,\one)=d\) and
\(g''(1,\one)=4d'(1)+3d''(1)\), giving the first identities.
The normalized derivatives of \(d\) through order two are
\(O_\rho((1+A)^2)\), by the convergent resolvent trace series.

For completeness we compute the coordinate derivatives at
the singular parity root. For \(t<1\), set
\[
 \Lambda(t)=\sqrt t\,P(I+\sqrt t\,P)^{-1},
       \qquad P=\begin{pmatrix}0&T\\T^\top&0\end{pmatrix}.
\]
With the unit constant and parity vectors \(c,s\) used above,
\[
 \Lambda(t)=\Gamma_{\sqrt t}
      +\frac{\sqrt t}{1+\sqrt t}cc^\top
      -\frac{\sqrt t}{1-\sqrt t}ss^\top .
\]
Write \(\tau=1-t\). Its expansion is
\[
 \Lambda(t)=-\frac2\tau ss^\top+H_0+\tau H_1+O(\tau^2),
 \qquad H_0=\Gamma+\tfrac12cc^\top+\tfrac32ss^\top .
\]
For clarity the coefficient of \(\tau\) is explicitly
\[
 H_1=-\tfrac12(\Gamma-\Gamma^2)
                  -\tfrac18cc^\top+\tfrac18ss^\top .
\]
Indeed \(\partial_s\Gamma_s|_{s=1}=\Gamma-\Gamma^2\).
The full block covariance satisfies the concrete bound
\[
 (1-\rho)|\Gamma_{ab}|\leq(2-\rho)A/N .
\]
To see this directly, \(\|\Omega\|\leq\|B\|\leq\rho\);
every row and column of \(\Omega\) has squared norm at most
\(A/N\). Writing \(R=(I+\Omega)^{-1}\), the identity
\(\Gamma=\Omega-\Omega R\Omega\) and
\(\|R\|\leq(1-\rho)^{-1}\) give the displayed estimate.
Also every row and column of \(\Gamma\) has norm at most
\(\sqrt{A/N}/(1-\rho)\). Consequently
\[
 |(\Gamma^2)_{ab}|\leq A/[N(1-\rho)^2] .
\]
Thus every entry of \(H_0,H_1\) is
\(O_\rho((1+A)/N)\), with the complete derivative and both
constant/parity contributions retained.

Let \(\partial_a\) here denote \(x_a\partial/\partial x_a\).
Multiaffinity in the rows gives, at \(x=\one\),
\[
 \partial_aF=F\Lambda_{aa},\qquad
 \partial_a\partial_bF
      =F(\Lambda_{aa}\Lambda_{bb}-\Lambda_{ab}^2)
                       \quad(a\ne b).
\]
The potential double pole in the second expression cancels
because the pole matrix \(ss^\top\) has rank one.
Define
\(q_{ab}(H)=s_a^2H_{bb}+s_b^2H_{aa}-2s_as_bH_{ab}\).
Taking the coefficient of \(\tau\), since \(g=-t\partial_tF\),
gives exactly
\[
 \frac{{\cal L}(u_a)}d=(H_0)_{aa}
                                  +2\frac{d'(1)}d s_a^2 ,
\]
and, for \(a\ne b\),
\[
 \frac{{\cal L}(u_a u_b)}d
 = (H_0)_{aa}(H_0)_{bb}-(H_0)_{ab}^2
       -2q_{ab}(H_1)+2\frac{d'(1)}d q_{ab}(H_0).
\]
Here \(s_a^2=1/(2N)\),
\(q_{ab}(H_i)=O_\rho((1+A)/N^2)\), and
\(d'(1)/d=O_\rho(1+A)\).
The stated coordinate bounds follow.
For a repeated coordinate use \(u_a^2=u_a\).
This evaluates every overlap type, including a row with
a column and two coordinates in the same part.
\end{proof}

\paragraph{Formalized finite occupancy input.}
The actual marked polynomials on the full vertex set \(V\) are
\[
 f_a(z)=\sum_{\substack{S\subseteq V\\a\in S}}\det P[S]z^{|S|},
 \qquad
 f_{ab}(z)=\sum_{\substack{S\subseteq V\\a,b\in S}}\det P[S]z^{|S|}.
\]
\texttt{HamiltonianCoordinateMinorPolynomial} proves these full
support expansions and their Euler readouts. The pair polynomial
includes the complete two-row inclusion--exclusion term, even when
the coordinates coincide. For distinct coordinates, writing \(L(z)=zP(I+zP)^{-1}\),
the actual row-replacement coefficient is
\(L_{aa}L_{bb}-L_{ab}L_{ba}\);
its off-diagonal subtraction is retained.

Put \(\Pi=ss^\top\) and, with \(z=1-\tau\), set
\[
 M(\tau)=I+(1-\tau)P+\Pi,\quad
 D(\tau)=\det M(\tau),\quad A(\tau)=I-M(\tau)^{-1}.
\]
The maintained proof derives that \(M(0)\) is invertible and
\(M(0)^{-1}=K+\Pi\) from the original stochastic and gap hypotheses.
For \(\tau\ne0,-1\), when \(M(\tau)\) is invertible, the complete
physical resolvent is exactly
\[
 L(1-\tau)=A(\tau)-\frac{\Pi}{\tau(1+\tau)} .
\]
Define, without suppressing either ordered cross term,
\[
 Q_{ab}(A)=\Pi_{aa}A_{bb}+\Pi_{bb}A_{aa}
                       -\Pi_{ab}A_{ba}-\Pi_{ba}A_{ab}.
\]
The rank-one identity
\(\Pi_{aa}\Pi_{bb}-\Pi_{ab}\Pi_{ba}=0\)
cancels the apparent squared pole. Consequently the actual marked
polynomials satisfy the exact cleared formulas
\[
 \begin{split}
 f_a(1-\tau)
   &=D(\tau)\left(\frac{\tau}{1+\tau}A_{aa}(\tau)
                         -\frac{\Pi_{aa}}{(1+\tau)^2}\right),\\
 f_{ab}(1-\tau)
   &=D(\tau)\left(\frac{\tau}{1+\tau}
         (A_{aa}A_{bb}-A_{ab}A_{ba})(\tau)
                         -\frac{Q_{ab}(A(\tau))}{(1+\tau)^2}\right)
                         \quad(a\ne b).
 \end{split}
\]
These identities, with precisely the displayed nonsingularity
conditions, are formalized in \texttt{HamiltonianPhysicalOccupancy},
using \texttt{HamiltonianParityRootRegularization},
\texttt{HamiltonianParityResolventSplit} and
\texttt{HamiltonianTwoRowDeterminant}.
The signed Euler readouts are sums over every principal support;
no independent coordinate model is introduced.
The regularized root jets now have maintained Lean proofs:
\[
 \begin{split}
 D(0)&=2d,\\
 D'(0)&=D(0)\bigl(\tfrac12-\operatorname{tr}\Gamma\bigr),\\
 A(0)&=\Gamma+\tfrac q4\one\one^\top,\\
 A'(0)&=-\Gamma+\Gamma^2-\tfrac q8\one\one^\top+\Pi,
       \qquad q=1/N.
 \end{split}
\]
The normalization is an exact constant-direction determinant
factorization. The derivative identity retains and cancels both
constant-direction cross products using stochastic balance.
The cleared formulas extend by their actual continuous limits;
the singular physical inverse is never evaluated at the root.

Precisely, define the full-principal-support functional
\[
 {\cal L}_{\rm pr}(f)=
   -\sum_{S\subseteq V}\frac{|S|}{2}\det P[S]\,f(S).
\]
The maintained polynomial Euler readouts and root derivatives prove
\[
 \frac{{\cal L}_{\rm pr}(u_a)}d
 =A_{aa}(0)+\bigl(\tfrac32+\operatorname{tr}\Gamma\bigr)\Pi_{aa},
\]
and, for \(a\ne b\),
\[
 \begin{split}
 \frac{{\cal L}_{\rm pr}(u_au_b)}d
 ={}&A_{aa}(0)A_{bb}(0)-A_{ab}(0)A_{ba}(0)
       -Q_{ab}(A'(0))\\
 &+\bigl(\tfrac32+\operatorname{tr}\Gamma\bigr)Q_{ab}(A(0)).
 \end{split}
\]
The Lean statements are the multiplied identities and require no
division by \(d\). Under the original nonnegative stochastic,
entry-cap and gap hypotheses, for every \(A\geq1\), they also prove
\[
 (1-\rho)|{\cal L}_{\rm pr}(u_a)|\leq5|d|Aq,\qquad
 (1-\rho)^2|{\cal L}_{\rm pr}(u_au_b)|
                         \leq50|d|A^2q^2\quad(a\ne b).
\]
In particular the covariance-square coefficient is independently
bounded by
\((1-\rho)^2|(\Gamma^2)_{ab}|\leq Aq\).
These are actual signed moments, rather than bounds on unsigned
minor mass. The maintained owners are
\texttt{HamiltonianParityRootJets},
\texttt{HamiltonianParityRootNormalization},
\texttt{HamiltonianClearedRootDerivative},
\texttt{HamiltonianRootOccupancyMoments},
\texttt{HamiltonianOccupancyEntryBounds},
\texttt{HamiltonianRootOccupancyBounds} and
\texttt{HamiltonianSignedOccupancyBounds}.
The arbitrary weighted balanced two-support readout and complete
joint first-order correction now also have maintained proofs.
The exact analytic size-derivative readout and the complete
remainder/tail and headline assembly remain unfinished.

\paragraph{Formalized complete signed main correction.}
For every real function on the finite vertex supports,
\texttt{HamiltonianSignedBipartiteFunctional} proves the actual identity
\[
 {\cal L}_{\rm pr}(f)=
 -\sum_{J\subseteq[1,N]}(-1)^{|J|}|J|
     \sum_{\substack{I\subseteq[1,N]\\|I|=|J|}}
          \det T[I,J]^2 f(I\sqcup J).
\]
The row and column supports are independent. Unbalanced determinants
are exactly zero, and the empty supports have their actual zero
rank coefficient. In particular the constant readout is \(d\).

Put \(h_a=-\Gamma_{aa}\) and \(v=(\tfrac12 I-K)h\).
The actual field satisfies
\((1-\rho)^2|v_a|\leq19A^2/N\).
\texttt{HamiltonianMinorFieldBudget} matches
\[
 \alpha(I,J)=r(r-1)/N+u^\top\Gamma u+2v^\top u
\]
to the precise first-order term of the maintained uniform minor
theorem. The finite signed functional is linear. The exact identities
\(u_a^2=u_a\) and
\[
 r(r-1)=\tfrac14\left(\sum_{\substack{a,b\in V\\a\ne b}}u_au_b
                                  -\sum_{a\in V}u_a\right)
\]
retain all overlaps and both orders of every off-diagonal pair.
With \(p=A/(1-\rho)\), the compiled full signed consumer gives
\[
 \left|{\cal L}(\alpha)\right|\leq900|d|p^3/N .
\]
Its proof controls the rank, diagonal, ordered off-diagonal and
full field contributions by
\[
 \frac{|d|}{N}
       \bigl(53p^2+20p^2+400p^3+380p^3\bigr)
             \leq\frac{900|d|p^3}{N},\qquad p\geq1.
\]
These bounds are taken after the actual signed coordinate
moments have been evaluated; no unsigned determinant mass enters
the main correction. The full-support finite identities are in
\texttt{HamiltonianSignedMinorFunctional}, and the actual stochastic
and balanced-support bounds are in
\texttt{HamiltonianJointSignedCorrection}. The separate complete
small-support residual now has a maintained quantitative proof,
as follows.

\paragraph{Formalized full small-support residual.}
Put \(R=\log(2N)^3\), \(r=|J|\), and
\[
 E_R(I,J)=\mathbf1_{r\leq R}
 \left[
  \left(\frac{\operatorname{per}T[I^c,J^c]}{\operatorname{per}T}\right)^2
                       -1-\alpha(I,J)\right].
\]
The actual retained error is the signed packet
\[
 \mathcal E_R=-\sum_{J\subseteq[1,N]}(-1)^r r
       \sum_{\substack{I\subseteq[1,N]\\|I|=r}}
                   \det T[I,J]^2 E_R(I,J).
\]
\texttt{HamiltonianAbsoluteMinorMass} evaluates its weighted mass
by the actual positive Gram spectrum and proves
\[
 \sum_J |J|\sum_{|I|=|J|}\det T[I,J]^2
                         \leq2A\det(I+B^\top B).
\]
The constant-direction factor is exactly \(2\).
\texttt{HamiltonianSignedMinorResidual} consumes the full uniform
two-power minor theorem. For every fixed \(0<\delta\leq1\)
and \(K\geq0\), one threshold covers every finite order,
every admissible kernel with \(A\leq(\log N)^K\), and all
independent deleted supports. It proves
\[
 |\mathcal E_R|\leq
 2A C_\delta^{\rm minor}
       \frac{\log(2N)^{12+24\lceil K\rceil}}{N^2}
                                  \det(I+B^\top B).
\]
The unsigned mass is used only for this residual. The previously
evaluated full first-order term retains its signed determinant
normalization.

\paragraph{Formalized reveal input to the large-tail estimate.}
For a finite family of injections separated by their complete row
histories, the sequential products
\[
 q_\sigma=\prod_{\text{revealed rows }i}
 \frac{M_{i,\sigma(i)}}{\sum_{j\ {\rm still\ available}}M_{ij}}
\]
have total mass at most \(1\).
\texttt{HamiltonianWeightedReveal} proves this with each actual
denominator and every branch retained. Define
\(w_\sigma=\prod_i M_{i,\sigma(i)}\) and let \(C_\sigma\)
be the sum of the logarithms of these denominators.
\texttt{HamiltonianWeightedRevealLog} proves, when the total
matching weight is positive,
\[
 \log\sum_\sigma w_\sigma
       \leq\frac{\sum_\sigma w_\sigma C_\sigma}
                      {\sum_\sigma w_\sigma}.
\]
It handles zero matching weights directly and derives the
inequality from \(\log x\leq x-1\).
The row-order average giving the global row-cap permanent bound,
its large-support complementary-permanent tail application, and
the final Hamiltonian comparison remain formalization obligations.

\begin{theorem}[Relative bipartite comparison at logarithmic density]
\label{thm:criticalrelative}
Fix \(0\leq\rho<1\) and \(C_0>0\).
Let \(T\) be real nonnegative doubly stochastic of order \(N\),
\[
 B=T-J_N/N,\quad \|B\|\leq\rho,\quad
 \max T\leq A/N,\quad 2\leq A\leq C_0\log(2N).
\]
Put
\[
 {\cal A}(T)=
        \frac{\det(I+BB^\top)}{\det(I-BB^\top)},\qquad
 P=\begin{pmatrix}0&T\\T^\top&0\end{pmatrix}.
\]
For some fixed finite \(M=M(\rho,C_0)\), uniformly in \(T\),
\[
 Z_{\rm H}(P)=\frac{(N!)^2}{2N\,N^{2N}}
 \left[1+O_{\rho,C_0}\left(
       \frac{\log(2N)^M}{N}
         +\frac{\log(2N)^M}{N^2}{\cal A}(T)\right)\right].
\]
In particular the full relative asymptotic holds whenever
\({\cal A}(T)\leq N^{2-\varepsilon}\) for fixed
\(\varepsilon>0\).
A sufficient explicit density condition is
\[
 A\leq a\log N,\qquad a\kappa_\rho<2,\qquad
 \kappa_\rho=\frac1{\rho^2}
                   \log\frac{1+\rho^2}{1-\rho^2};
\]
at \(\rho=0\) use the limiting value \(\kappa_0=2\).
\end{theorem}
\begin{proof}
Use the exact two-support Liu identity, normalized by
\((\operatorname{per}T)^2\).
On \(r\leq\log(2N)^3\), insert the complete expansion of
Theorem~\ref{thm:minorjet}.
The constant term is \(d=\det(I-BB^\top)\).
The \(r(r-1)/N\) term has normalized size \(O_\rho(A^2/N)\)
by Lemma~\ref{lem:signedoccupancy}.
The diagonal part of \(u^\top\Gamma u\) has normalized
size \(O_\rho(A^2/N)\), using all \(2N\) diagonal coordinates.
Its off-diagonal part has normalized size \(O_\rho(A^3/N)\):
there are at most \(4N^2\) ordered pairs, a coefficient
\(O_\rho(A/N)\), and a signed occupancy moment
\(O_\rho(A^2/N^2)\).
The field \(2v^\top u\) has the same normalized bound,
since \(\|v\|_\infty=O_\rho((1+A)^2/N)\).
These estimates are applied after the coupled signed moments
have been evaluated. The full main correction is consequently
\(O_\rho((1+A)^3/N)\), without an exponential determinant loss.

The moment identities cover all sizes.
Extending the truncated polynomial corrections to all \(r\)
costs only an exponentially small tail: their magnitude is
bounded by a polynomial in \(r,A,\log(2N)\), and the
Cauchy--Binet tail has \(r>\log(2N)^3\), \(A=O(\log N)\).
The same tail estimate with all complementary-permanent
ratios is
\(\exp(-c_{\rho,C_0}\log(2N)^3\log\log(2N))\),
as in the relative bipartite proof.

Only the second-order residual is now bounded by absolute
minor mass. Exactly, for \(Q=BB^\top\),
\[
 \sum_{r\geq1}r e_r(TT^\top)
 =\det(I+TT^\top)
       \operatorname{tr}\bigl(TT^\top(I+TT^\top)^{-1}\bigr)
 \leq2A\det(I+Q).
\]
After dividing by \(d\), its error is at most a fixed
logarithmic power times \(N^{-2}{\cal A}(T)\).
The determinant normalization remains explicit.

Thus
\[
 \frac{2N Z_{\rm H}(P)}{(\operatorname{per}T)^2}
 =d\left[1+O_{\rho,C_0}\left(
       \frac{\log(2N)^M}{N}
       +\frac{\log(2N)^M}{N^2}{\cal A}(T)\right)\right].
\]
The full permanent expansion has relative first correction
\(O_\rho(A^2/N)\).
Its squared reciprocal determinant cancels \(d\) exactly.
Multiplication gives the displayed relative formula; all product error terms are
absorbed by the displayed bound, since that permanent
relative correction tends to zero.

Finally the eigenvalues \(x\) of \(Q\) lie in \([0,\rho^2]\).
The function \(x^{-1}\log((1+x)/(1-x))\) is increasing there,
as its power series has nonnegative coefficients.
Also \(\operatorname{tr}Q\leq A-1\).
Consequently
\[
 \log{\cal A}(T)\leq\kappa_\rho\operatorname{tr}Q
                                      \leq\kappa_\rho A.
\]
The explicit sufficient condition leaves a positive power
between the two-power remainder and this amplification.
This proves the claimed relative range.
\end{proof}


\begin{corollary}[An inverse theorem for failure of relative universality]
\label{cor:relativeinverse}
Under Theorem~\ref{thm:criticalrelative}, fix
\(\eta>0\). If
\[
 \left|\frac{2N\,N^{2N}}{(N!)^2}Z_{\rm H}(P)-1\right|\geq\eta,
\]
then, for sufficiently large \(N\),
\[
 \log{\cal A}(T)\geq
            2\log N-O_{\rho,C_0,\eta}(\log\log(2N)),
\]
and consequently
\[
 \|T-J_N/N\|_F^2\geq
    \frac2{\kappa_\rho}\log N
              -O_{\rho,C_0,\eta}(\log\log(2N)).
\]
Thus a fixed relative counting defect requires a collective
nontrivial singular-value mass on the logarithmic scale.
A single fixed nontrivial eigenvalue cannot supply that
defect; all singular values are retained in \({\cal A}(T)\).
\end{corollary}
\begin{proof}
For large \(N\), the first error term
\(C_{\rho,C_0}\log(2N)^M/N\) is less than \(\eta/2\).
The second must therefore be at least \(\eta/2\).
It forces
\({\cal A}(T)\geq c_{\rho,C_0}\eta N^2/\log(2N)^M\).
Take logarithms, then use
\(\log{\cal A}(T)\leq\kappa_\rho\|B\|_F^2\).
This concerns a defect of the relative count, not just
Hamiltonian existence, which was already established in a
much wider diffuse range.
\end{proof}

\section{Relative prime-circle enumeration at the actual prime density}
\label{sec:relativeprimecircles}
The logarithmic-density theorem has an independent arithmetic
consumer on the unweighted prime-sum graph itself.
The row repair must preserve the coefficient of \(\log N\)
in the largest entry; a generic \(\log N\) inverse bound is
too coarse for that purpose.

\begin{lemma}[A diffuseness-sensitive scaling repair]
\label{lem:diffusionscaling}
Let \(Q\) be a symmetric nonnegative bipartite matrix on
\(n\) vertices with equally sized parts, row degrees
\(|d_i-1|\leq\varepsilon\), and \(\max Q_{ij}\leq A/n\),
\(A\geq2\).
Suppose its degree-normalized matrix has nontrivial gap
at least a fixed \(\delta\in(0,1)\).
If \(\varepsilon\log(2A)^2\) is sufficiently small, it has
a stochastic repair
\[
 \widehat P=\diag(1+z)Q\diag(1+z),\qquad
 \|z\|_\infty\leq C_\delta\varepsilon\log(2A).
\]
The repaired kernel differs from \(Q\) in operator norm
by \(O_\delta(\varepsilon\log(2A))\).
\end{lemma}
\begin{proof}
We refine the inverse estimate in Lemma~\ref{lem:scaling}.
For the reversible walk \(R=\diag(d)^{-1}Q\), remove its
constant and parity projections.
The absolute row sum of its \(j\)-th centered power is at
most a fixed constant.
After the first step the squared row norm is at most
\(C A/n\); the fixed spectral gap and Cauchy--Schwarz
therefore give the alternative bound
\(C_\delta\sqrt A(1-\delta)^{j-1}\).
The stationary degrees are between fixed constants, so
these estimates hold also in the stationary weighted norm.
Summing their minimum gives
\[
 \sum_{j\geq0}\min\{C,C_\delta\sqrt A
                                      (1-\delta)^{j-1}\}
                         \leq C_\delta\log(2A).
\]
The constant mode in \((I+R)^{-1}\) costs only \(1/2\);
the parity mode is removed by the same balanced gauge
as in the earlier scaling proof.
Multiplying by \(\diag(d)^{-1}\) proves that the inverse
of \(\diag(d)+Q\), on its compatible subspace, has
infinity norm at most \(C_\delta\log(2A)\).
Changing to the unweighted balanced gauge costs only a
fixed multiple of this bound.

Now solve the exact equation
\[
 (\diag(d)+Q)z=\one-d-z\mathbin{\odot}(Qz)
\]
in the balanced gauge.
The right side is compatible on both parts, including
its quadratic term.
On the ball
\(\|z\|_\infty\leq 2C_\delta\varepsilon\log(2A)\),
the inverse and quadratic Lipschitz bounds give a
contraction provided \(\varepsilon\log(2A)^2\) is small.
The fixed point is positive after adding one and gives
the exact row sums.
The operator perturbation follows from
\(\|\diag(z)Q+Q\diag(z)+\diag(z)Q\diag(z)\|
 \leq C\|z\|_\infty\).
\end{proof}

\paragraph{Formalized diffuseness estimates and positive inverse reference.}
The maintained owner \texttt{HamiltonianMixingResolvent} proves
an exact finite Neumann expansion, a full-tail smoothing estimate,
and its logarithmic cutoff.
\texttt{HamiltonianDiffuseRepairKernel} applies them to the actual
stochastic cross block. If \(T\) is nonnegative and doubly stochastic,
\(T_{ij}\leq a/m\), \(a\geq1\), and
\(\|T-J_m/m\|\leq\rho\in(0,1)\), its exact
constant/parity-corrected kernel \(K\) satisfies
\[
 (1-\rho)\max_i\sum_j|K_{ij}|
 \leq
 2(1-\rho)\left(\frac{\log(2a)}{-\log\rho}+3\right)+1.
\]
Thus the logarithmic cost is independently bounded, rather than
assumed as a hypothesis on the desired scaling.

For a nonnegative nearly stochastic cross block \(T\), let
\(r_i=\sum_jT_{ij}\), \(t_j=\sum_iT_{ij}\), with
\(|r_i-1|,|t_j-1|\leq\varepsilon\in(0,1/2]\).
Put \(c=1+2\varepsilon\) and
\(\tau=\sum_i(c-r_i)=\sum_j(c-t_j)>0\). The auxiliary matrix
\[
 T^*_{ij}=\frac{T_{ij}}c+
               \frac{(c-r_i)(c-t_j)}{c\tau}
\]
is nonnegative and exactly doubly stochastic.
\texttt{HamiltonianStochasticCompletion} proves these identities and
\[
 \max_i\sum_j|T_{ij}-T^*_{ij}|\leq6\varepsilon,\qquad
 \max_j\sum_i|T_{ij}-T^*_{ij}|\leq6\varepsilon,\qquad
 \|T-T^*\|\leq6\varepsilon .
\]
If \(T_{ij}\leq a/m\), it also proves
\[
 T^*_{ij}\leq\frac{a+9\varepsilon}{(1+2\varepsilon)m},
 \qquad
 \|T^*-J_m/m\|\leq\rho+6\varepsilon
 \quad\hbox{when }\|T-J_m/m\|\leq\rho.
\]
The real and complex mass-to-operator estimate is formalized in
\texttt{HamiltonianSchurBudget}.
All thirty-five new declarations in these four owners pass the
focused foundational-axiom audit.
The completion is only an inverse reference; it is not the final
kernel and supplies no assertion about preserving support.
\paragraph{Formalized full repair on the original support.}
The compatible inverse transfer and the nonlinear row equations
are now proved as well.
\texttt{HamiltonianCompatiblePerturbation} retains the exact parity
projection under the small perturbation.
\texttt{HamiltonianDegreeJacobian} proves its compatibility for
the actual diagonal \(\diag(d)\), before normalization.
\texttt{HamiltonianDiffuseJacobianInverse} consumes the positive
stochastic reference above and its entire \(6\varepsilon\) error.
\texttt{HamiltonianDegreeScalingFixedPoint} retains
\(-z\mathbin{\odot}(Qz)\) in the full quadratic contraction.
\texttt{HamiltonianDegreeScaling} uses the already proved exact
conservation law to remove the final parity residual.

The complete finite output is compiled in
\texttt{HamiltonianDiffuseScaling}. For the nearly stochastic
cross block above, set
\[
 g=2\left(\frac{\log(2(a+9\varepsilon))}{-\log\rho}+3\right)
                       +\frac1{1-\rho},\qquad
 r=4\varepsilon g .
\]
For every \(\rho\in(0,1)\), under the explicit hypotheses
\[
 \|T-J_m/m\|\leq\rho-6\varepsilon,\qquad
                    \varepsilon g^2\leq1/256,
\]
there is one vector \(z\) with \(\|z\|_\infty\leq r\)
and every \(1+z_i>0\), for which
\[
 \widehat T_{ij}=(1+z_{A,i})T_{ij}(1+z_{B,j})
\]
is exactly doubly stochastic and nonnegative. The same witness satisfies
\[
 \widehat T_{ij}=0\ \Longleftrightarrow\ T_{ij}=0,\qquad
 \widehat T_{ij}\leq(1+r)^2a/m,\qquad
 \|\widehat T-T\|\leq(2r+r^2)(1+\varepsilon).
\]
All twenty-eight declarations in these six further owners pass
the focused foundational-axiom audit.
For fixed \(\rho\), the bound is
\(g=O_\rho(\log(2a))\), giving the required
\(\varepsilon\log(2a)^2\) smallness scale.
The native degree-normalized spectral consumer of
Lemma~\ref{lem:diffusionscaling} is now formalized in
\texttt{HamiltonianDiffuseSpectralScaling}, with the advertised
logarithmic cap cost. Here are explicit uniform constants.
Put \(\rho=1-\delta/2\) and
\[
 G_\delta=\frac4{-\log\rho}
           +\frac{6+(1-\rho)^{-1}}{\log4},\qquad
 c_\delta=\min\left\{\frac{\delta}{28}(\log4)^2,\,
                    \frac1{256G_\delta^2}\right\},
 \qquad C_\delta=20G_\delta.
\]
For the physical cap \(A/n\), take \(a=A/2\geq1\);
then \(L=\log(4a)=\log(2A)\).
The already proved scalar inverse estimate has a free scale
parameter: applying it at that scale equal to \(a\) gives
\[
 g(a+9\varepsilon,\rho)\leq G_\delta L .
\]
Its scalar hypotheses require only
\(a\leq(1+\varepsilon)a\), which follows from
\(\varepsilon\geq0\); no dimension bound is substituted.

The single smallness hypothesis
\(\varepsilon L^2\leq c_\delta\) first gives
\(\varepsilon\leq\delta/28<1/2\), and then
\(\varepsilon g^2\leq1/256\).
The actual normalized spectral packet, including the two simple
exceptional modes, yields
\[
 \|T-J/m\|\leq1-\delta+8\varepsilon
             \leq\rho-6\varepsilon.
\]
Thus the full nonlinear repair applies to the original matrix.
Its same witness has \(R=4\varepsilon g\leq1/64\), exact
stochastic rows and columns, positive multipliers, identical
support, and
\[
 \|z\|_\infty\leq C_\delta\varepsilon L,\qquad
 \|\widehat T-T\|\leq C_\delta\varepsilon L,\qquad
 \widehat T_{ij}\leq2T_{ij}.
\]
The exact identity
\(\widehat Q-Q=ZQ+QZ+ZQZ\) retains the full quadratic term.
The zero-defect case uses the identity scaling explicitly.
The theorem
\texttt{hamiltonian\_normalized\_spectral\_physical\_cap\_scaling}
has exactly the physical cap \(A/n\) and the cost \(\log(2A)\);
no additional \(\varepsilon\leq1/2\) hypothesis is assumed.
The actual constant-edge prime transfer, uniform degrees and
half-gap are now formalized in Section~\ref{sec:regularfactors}.
\paragraph{Formalized actual prime-density count.}
\texttt{HamiltonianMatrixGeometry} proves that a finite relabeling
preserves the Euclidean operator norm and that the bipartite lift
has exactly the norm of its cross block.
\texttt{HamiltonianPrimeCrossBlock} identifies the actual block
\[
 T_{m,ij}=\frac{\log(2m)}{2m}
     {\bf1}_{\{2i+2j+3\ {\rm prime}\}},
 \qquad 0\leq i,j<m .
\]
Its rows and columns are exactly the corresponding full prime
degrees. Its centered norm is exactly the full constant/parity
operator error, without a dimension factor.

\texttt{HamiltonianPrimeScalingWindow} proves the uniform
smallness window at this density: for every \(B,t>0\), for all
sufficiently large real \(L\), putting
\(\varepsilon=B/\sqrt L\), \(a=L/2\), and \(\rho=2/3\) gives
\[
 0<\varepsilon\leq1/100,\qquad
 \varepsilon g(a+9\varepsilon,\rho)^2\leq t .
\]
The inverse budget is bounded by \(K\log L\) with the explicit
positive constant
\[
 K=\frac{2(1+\log2)}{-\log(2/3)}+9.
\]
The already proved limit \((\log L)^2/\sqrt L\to0\)
closes the condition on an entire terminal interval.

\texttt{HamiltonianPrimeStochasticScaling} applies the full
nonlinear repair to the original block. For every \(\eta>0\),
for every sufficiently large \(m\), there is one positive scaling
with \(\|z\|_\infty\leq\eta\) whose repaired block
\(\widehat T\) is nonnegative, doubly stochastic, and has exactly
the original prime support. The same witness satisfies
\[
 \widehat T_{ij}\leq\frac{(3/5)\log(2m)}m,\qquad
 \|\widehat T-J_m/m\|\leq3/5 .
\]
Finally \texttt{HamiltonianPrimeRelativeCount} consumes these
actual arithmetic inputs in the complete signed counting theorem:
\[
 \left|\left(\frac{m^m}{m!}\right)^2
       m\,\operatorname{cyc}(\operatorname{bip}(\widehat T))-1\right|
                       \leq m^{-3/5}.
\]
All thirty declarations in these five maintained owners pass the
focused foundational-axiom audit.
\paragraph{Formalized entropy and the actual unoriented count.}
\texttt{HamiltonianEntropy} proves the finite supporting-line
inequality, sums every row and column potential exactly, and
proves optimality over every feasible weighting on the original
support, including weightings with zero entries.
\texttt{HamiltonianPrimeEntropy} applies this result to the actual
prime scaling and identifies its entropy exactly.

\texttt{HamiltonianCycleScaling} proves the common vertex-product
weight for every Hamilton cycle before summing the cycles.
\texttt{HamiltonianCycleOrientation} counts actual single-orbit
permutations and their reversal classes. For at least three
vertices, reversal has no fixed point, so the directed count is
exactly twice the unoriented count. Cyclic rotation is already
removed by the permutation representation.

\texttt{HamiltonianPrimeCircleEntropyCount} proves existence and
uniqueness of the finite maximum entropy value \({\cal H}_{2m}\),
and the factorial-normalized relative formula on an entire
terminal interval:
\[
 \left|\left(\frac{m^m}{m!}\right)^2
       (2m)e^{-2{\cal H}_{2m}}h(G_{2m})-1\right|
                                    \leq m^{-3/5}.
\]
The count is the actual finite cardinality of the reversal
classes. \texttt{HamiltonianPrimeCircleAsymptotic} verifies that
its edge relation, under the proved odd/even bijection, is exactly
the relation of distinct integer vertices with prime sum.

\texttt{HamiltonianStirlingNormalization} sums Robbins' exact
stepwise logarithmic bound and takes its established limit.
For every \(m\geq1\), it proves
\[
 \left|\frac{s_m^2}{\pi}-1\right|\leq\frac1{3m},
 \qquad s_m=\frac{m!}{\sqrt{2m}(m/e)^m}.
\]
The full final consumer therefore proves, for every sufficiently
large \(m\),
\[
 \left|\frac{h(G_{2m})}
              {\pi e^{2{\cal H}_{2m}-2m}}-1\right|
                                    \leq 2m^{-3/5}.
\]
All fifty-four declarations in these seven further maintained
owners pass the focused foundational-axiom audit.
Both formulas in Theorem~\ref{thm:relativeprimecircles} now have
complete maintained proofs. The sharp physical repair, full
prime-edge cuts and integral regular-factor criterion with its
actual adjacency loss are now formalized as well in
Section~\ref{sec:regularfactors}. Other advertised arithmetic
consumers and the paper's complete public integration remain
separate formalization obligations.

Let \(G_N\), for even \(N=2m\), have vertices \(1,\ldots,N\)
and edges at distinct prime sums.
For a proper fractional edge weighting \(x\), meaning
\(x_e\geq0\) and \(\sum_{e\ni v}x_e=1\) for every vertex,
write
\[
 {\cal H}_N=\max_x\left(-\sum_{e\in E(G_N)}x_e\log x_e\right).
\]
The optimization is over the actual finite prime-sum graph.
It incorporates its nonuniform degrees and every allowed
edge, with no averaging of the graph itself.

\begin{theorem}[Full relative prime-circle count from its exact entropy]
\label{thm:relativeprimecircles}
Uniformly through even \(N\to\infty\),
\[
 h(G_N)=
 \frac{(m!)^2}{2m\,m^{2m}}\,e^{2{\cal H}_N}
                  \left(1+O(N^{-3/5})\right)
       =\pi e^{\,2{\cal H}_N-N}
                  \left(1+O(N^{-3/5})\right).
\]
Here \(h\) counts unoriented Hamilton cycles modulo cyclic
rotation, as in the earlier prime-circle statements.
The implied constant is absolute eventually.
The formula is unconditional and uses neither RH nor a
prime-pair asymptotic. It is a relative count, rather than
an error of exponential size \(e^{o(N)}\).
\end{theorem}
\begin{proof}
Put \(Q_0=(\log N/N)A_{G_N}\).
Lemma~\ref{lem:unweightedtransfer} supplies the actual
uniform degree defect
\[
 \varepsilon_N=O(\log\log N/\log N)
\]
and operator distance \(o(1)\) from a stochastic kernel
with nontrivial ratio \(1/2+o(1)\).
Consequently the degree normalization of \(Q_0\) also
has nontrivial ratio \(1/2+o(1)\): diagonal normalization
is \(I+O(\varepsilon_N)\), and its constant/parity
vectors differ by \(O(\varepsilon_N)\) from the
normalized constant/parity vectors.
Its largest entry is at most \(\log N/N\).

Apply Lemma~\ref{lem:diffusionscaling} with \(A=\log N\).
The compatibility quantity
\(\varepsilon_N\log(2A)^2=O((\log\log N)^3/\log N)\)
tends to zero. We obtain
\[
 P=\diag(1+z)Q_0\diag(1+z),\qquad
 \|z\|_\infty=O((\log\log N)^2/\log N)=o(1).
\]
Its cross block \(T\), of order \(m=N/2\), is doubly
stochastic and has
\[
 \|T-J_m/m\|\leq1/2+o(1),\qquad
 m\max T_{ij}\leq(1/2+o(1))\log N.
\]
Thus for sufficiently large \(m\) we may take
\(\rho=3/5\), \(A=(3/5)\log(2m)\).
Here
\[
 A-1\leq(3/5)\log m
\]
because \((3/5)\log2\leq1\), so the sharp amplification
consumer of Theorem~\ref{thm:criticalrelative} applies.

The numerical margin is checked explicitly:
\[
 \kappa_{3/5}=\frac{25}{9}\log(17/8)<\frac94,\qquad
 {\cal A}(T)\leq m^{27/20}.
\]
For the strict logarithmic inequality, the first four
terms of \(e^{81/100}\) already exceed \(17/8\).
The two error terms in the comparison are therefore
\(O(\log(2m)^M(m^{-1}+m^{-13/20}))\),
which is \(O(m^{-3/5})\).
The full critical-density signed consumer is now estimated.

Finally set
\(w_v=(1+z_v)\sqrt{\log N/N}\), so that
\(P=\diag(w)A_{G_N}\diag(w)\).
Its positive stochastic edge weighting is the entropy
maximizer. Indeed for every proper \(x\), the supporting
hyperplane inequality for \(-x\log x\) shows optimality
of \(P\), because
\(\log P_{ij}=\log w_i+\log w_j\) on its support and
the vertex sums are fixed.
Strict concavity gives uniqueness of the edge weights;
the harmless bipartite scaling gauge does not change
\(\sum_v\log w_v\).
Hence
\[
 {\cal H}_N=-\sum_v\log w_v.
\]
Every Hamilton cycle has exactly the same \(P\)-weight
\(\prod_v w_v^2=e^{-2{\cal H}_N}\).
Multiplying the relative kernel count by
\(e^{2{\cal H}_N}\) proves the factorial formula.
Stirling gives
\((m!)^2/(2m\,m^{2m})=\pi e^{-N}(1+O(N^{-1}))\),
proving the final normalization.
All arithmetic inputs were already used in the proved
unweighted transfer; the improvement is the retained
signed compensation and its independently bounded remainder.
\end{proof}

\begin{remark}[What precision has changed]
The earlier elementary entropy evaluation gives
\(2{\cal H}_N=N\log N-N\log\log N+o(N)\).
Substituting only that coarse evaluation recovers only
an exponential-base formula.
The theorem above keeps the exact finite variational
quantity and proves a power-saving relative error around
it. It does not silently promote the coarse explicit
expression to a relative formula.
No effective numerical starting order, full all-orders
Filz theorem, or relative classification of every fixed
congruence filter is claimed here.
\end{remark}


\section{Provenance and verification status}
Birkhoff decomposition, integral maximum flow, matrix Freedman,
classical prime estimates, and the Brada\v{c}--Janzer
Hamiltonicity theorem are established inputs.
The cut proof of regular-factor extraction is supplied to expose
its exact loss, not claimed as a new maximum-flow principle.

Spectral control of matrix scaling is developed by
Kwok--Lau--Ramachandran~\cite{KLR}.
Degree-preserving spectral sparsification~\cite{CGPSS}
preserves weighted degrees in a reweighted sparse graph.
Dinitz--Krauthgamer--Wagner~\cite{DKW} obtain sparse regular
expanding subgraphs, with edge weights in \(\{1,2\}\),
from fixed-density regular expanders.
Our finite rounding conclusion instead has unit edge weights,
simplicity, exact prescribed degree and an explicit operator
error, for a diffuse weighted kernel.
Independent Birkhoff sampling does not ensure simplicity;
the quantified forbidden-edge repair addresses this distinction.
No claim is made that the local scaling principle itself is new.

The arithmetic advance proposed here is from the infinitely
many cases proved in~\cite{CFG} to all sufficiently large even
orders, together with near-optimal factors, Hamilton-cycle
packings, local resilience, and the complete fixed-congruence
classification.
The signed Fourier model keeps the actual residue phases;
density alone cannot distinguish the Hamiltonian and
disconnected cases.
The general counting theorem permits a fixed nontrivial
spectral ratio, whereas the earlier Krivelevich and
Byrne--Tait enumeration theorems~\cite{Krivelevich,BT} require
a ratio tending to zero at least logarithmically.
Our degree threshold \(D\geq(\log N)^{20,000,000}\)
is explicit and unoptimized; the recent spanning-path
theorem in~\cite{BJ} is an essential input.
Gap inheritance under random vertex restriction is a
known matrix-concentration mechanism; see
M\"uyesser~\cite{Muyesser} and its antecedents.
It is supplied here in the form needed by the counting
consumer, not claimed as a new sampling principle.
The weighted extension permits zero support, fixed spectral
ratio and sparse nonuniform degrees.
Cuckler--Kahn~\cite{CK} prove entropy counting in the Dirac
setting, including dense balanced bipartite analogues;
Barvinok's comparison~\cite{Barvinok} assumes all entries
positive with a lower bound relative to their maximum.
Neither stated hypothesis covers the sparse kernels here.
The entropy method of Christoph et al.~\cite{CycleFactors}
controls the expected number of cycles in a random regular
directed cycle-factor; it does not assert the weighted
one-cycle formula proved here.
The general-kernel extension does not condition independent
half-sets to be complementary. Its exact covariance identity
controls both restrictions together, with only polynomial-cost
balance conditioning.
The odd-order reduction retains endpoint scaling factors and
uses actual minor permanents rather than assuming row sums
survive deletion.
Unitary Cayley spectra and Hamiltonian path existence are
established facts~\cite{KS,Basic}; the added application concerns
quantitative counts through all growing moduli, now with a
power-saving relative error and the full arithmetic constant.

The determinant permanent approximation originates with
McCullagh~\cite{McC}. His moderate-deviation hypothesis requires
bounded moments of every order for \(N(D-J_N/N)\).
For the unitary kernels along primorials, those moments grow
as powers of \(N/\varphi(N)\), so that hypothesis does not
apply uniformly along all moduli.
Divoux--Wang~\cite{DW} prove fixed-entry permanent and
Hamiltonian formulas with a fixed singular gap.
Our relative argument uses their finite Gaussian, Wick and
Gamma estimates with explicit growing-parameter dependence.
It also separates the exact bipartite parity mode, which
violates an aperiodic norm hypothesis.
The signed bipartite cancellation and the all-modulus arithmetic
consumer, and the cyclic-layer comparison with its full
noncommuting product and arithmetic parity classification,
are the additional conclusions proposed here;
neither the Gaussian representation nor Liu's identity~\cite{Liu}
is claimed as new.
The extension to actual prime density requires a further argument:
the first Gaussian permanent correction is kept explicitly,
its change across complementary minors cancels to two-power
precision, and the remaining support-dependent quadratic
correction is evaluated by signed determinant derivatives.
Only its \(N^{-2}\) remainder pays the determinant amplification.
The resulting logarithmic-density theorem is distinct from
substituting a growing parameter into a fixed-bound asymptotic.
The exact-entropy normalization is retained in the relative
prime-circle count; its coarse explicit evaluation is not
asserted to have relative precision.
Our residue-pressure application determines the exponential
cost of every approximate residue profile, while retaining
the exact congruence identities of each prime circle.
No independence of its edges is assumed.
Targeted primary-literature searches have not located these
combined conclusions; that is preliminary priority evidence,
not a certificate of priority or career-level recognition.
The proofs have undergone the internal hypothesis and loss
checks recorded in STATUS.json. The maintained Lean development
now verifies finite cycle-subset cancellation, principal-minor
generating formulas, Cauchy--Binet and the full signed two-support
squared-minor cancellation, rank-one pole cancellation on every
nonempty support, and centered stochastic energy bounds.
These finite inputs are exported by \texttt{HamiltonianFinite}.
The \texttt{HamiltonianBlockCovariance} module also derives the
full covariance entry bound from stochasticity and the actual
centered spectral gap. The Gamma modules evaluate the normalized
Euler integral, prove exact signed inverse-Gamma central moments
and uniform raw-moment bounds, and derive both Chernoff tails
from the actual exponential moment. They also prove the actual
central-window transfer, signed moment truncation and an explicit
cubic averaging estimate with a fourth-order local remainder.
The final averaging consumer derives window integrability from
measurability and the local bound. The spectral-factor modules
also prove the rational derivative bounds, every fourth-derivative
cross term, the local factor ratio, endpoint derivative transfer,
and their actual Gamma-average consumer with explicit constants.
The matrix bridge now derives the Gram eigenvalues, their exact
energy sum and determinant product from the actual stochastic matrix,
and transfers the Gamma average on its proved window.
The exact permanent/partial-permutation rook/Gamma identity and
its integrability are also formalized. The full covariance correction
and the actual determinant second derivative give the explicit bound
\((1-\rho)^3|{\cal Q}_N(D)|\leq21A^2/N\).
The actual complex vertex coefficients obey
\((n+3)c_{n+3}=c_n-(n+2)c_{n+2}\) and \(|c_n|\leq1\)
for every degree. Finite Wick pairing bounds and the exact low-degree
identity retain the diagonal compensation
\(-\tfrac52\sum_i\Gamma_{ii}^3\); its actual stochastic-block
contribution is bounded by \(40A^3/((1-\rho)^3N^2)\).
The full covariance geometric series has a uniform bound in every
degree. Collective product inequalities also cancel the complete
linear sum before estimating the Gaussian/Gamma tails.
These are maintained Lean proofs, and do not assume the analytic
comparison theorem.
The complex Gaussian modules now prove polynomial integrability,
the full covariance-quotient half-plane identity, actual scalar moments,
finite-dimensional quadratic reweighting, and the polynomial Stein identity.
They derive every mixed Wick contraction from that actual integral,
with a uniform all-degree bound; signed polynomial coefficients are retained.
Arbitrary complex covariance also has a proved standard-Gaussian
realization. For the actual real symmetric block, the orthogonal
eigenbasis gives the full covariance, pointwise zero-sum vector,
and exact quadratic-energy identity. The actual stochastic Gaussian
exponential moment and collective affine-product bound have explicit
\(\exp(C_\delta\sqrt{NA})\) dependence and maintained proofs.
The real and complex Gaussian scalar normalizations are now reconciled.
The complete mixed-moment permanent recursion and the support-by-support
factorial-completion identity are formalized. They identify the actual
Gaussian affine product with the maintained rook polynomial, retaining
every signed coefficient and every \((N-k)!\) completion factor.
For nonnegative doubly stochastic \(D\), \(B=D-J/N\),
\(D_{ij}\leq A/N\), and \(\|B\|\leq\rho\), the direct maintained consumer is
\[
 |R_B(t)|\leq
 \exp\!\left(\frac{2|t|}{1-|t|\rho}\sqrt{N(A-1)}\right),
 \qquad |t|\rho<1.
\]
The complete original rook outside-window inverse-Gamma integral is
now formalized, with explicit constants, actual joint integrability,
the full Gaussian/Gamma Fubini transfer and exact cutoff conversion.
The local determinant factor is never extended outside its domain.
The compensated integral is now identified with the original rook polynomial,
including its actual complex normalizer. Bilinear Parseval retains the
complex squares and cancels every cross term. The positive real normalizer
is proved to be the determinant factor on the stochastic window; the actual
real compensated covariance equals the full resolvent.
The full mixed Wick readout, every finite support and degree profile,
complete deformation-coefficient convolutions and the finite high-diagram
\(N^{-2}\) bound now have maintained proofs. Trace cancellation also proves
the explicit complex circle envelope for the actual compensated rook function.
The complete moderate-degree sum, actual normalizer holomorphy and scaled
Cauchy coefficient/tail transfer now have maintained proofs.
The finite Gaussian jet has an explicit integrable dominating moment,
and its full polynomial state is matched to the actual compensated integral.
Every scalar degree assignment now has a proved unique support/composition
owner and an exact inverse. Actual scalar Taylor coefficients, their
integrability and their complete Gaussian Wick readout are formalized.
The positive-real jet error also proves exact Taylor agreement with the
actual compensated rook function through every prescribed finite order.
The finite jet's full homogeneous decomposition, all-order actual
rook/Wick Taylor identity and exact covariance formal-series adapter now
have maintained proofs. The actual observable therefore has the full signed
matrix-resolvent Wick coefficients used by the uniform coefficient bounds.
Hamiltonian-cycle and
coordinate-occupancy identities, uniform permanent and minor remainders,
covariance-correction derivatives and arithmetic specializations remain
formalization obligations. In particular,
none of the headline
counting or inverse theorems is yet Lean-certified.
Compilation of this manuscript does not provide mathematical or
formal certification, and independent referee verification remains
outstanding.

\begin{thebibliography}{99}
\bibitem{DW}
A. Divoux and S. Wang,
\emph{The Maximum Number of Hamiltonian Cycles in Tournaments},
2026. \url{https://arxiv.org/abs/2610.23755}.
\bibitem{McC}
P. McCullagh,
\emph{An asymptotic approximation for the permanent of a doubly
stochastic matrix},
Journal of Statistical Computation and Simulation 84 (2014), 404--414.
\url{https://doi.org/10.1080/00949655.2012.712122}.
\bibitem{Liu}
C. J. Liu,
\emph{Enumeration of Hamiltonian cycles and paths in a graph},
Proceedings of the American Mathematical Society 111 (1991), 289--296.
\bibitem{BJ}
D. Brada\v{c} and O. Janzer,
\emph{Hamiltonicity of regular sublinear expanders},
2026. \url{https://arxiv.org/abs/2605.15043}.
\bibitem{Tropp}
J. A. Tropp, \emph{Freedman's inequality for matrix martingales},
2011. \url{https://arxiv.org/abs/1101.3039}.
\bibitem{CFG}
H.-B. Chen, H.-L. Fu and J.-Y. Guo,
\emph{Hamiltonicity in Prime Sum Graphs},
Graphs and Combinatorics 37 (2021), 209--219.
\url{https://doi.org/10.1007/s00373-020-02241-1}.
\bibitem{SLZ}
W. Shang, B. Li and S. Zhang,
\emph{An infinite version of prime circles},
Graphs and Combinatorics 41 (2025), article 112.
\url{https://doi.org/10.1007/s00373-025-02976-9}.
\bibitem{Vaughan}
R. C. Vaughan, \emph{The Hardy--Littlewood Method},
second edition, Cambridge University Press, 1997.
\bibitem{KLR}
T. C. Kwok, L. C. Lau and A. Ramachandran,
\emph{Spectral Analysis of Matrix Scaling and Operator Scaling},
SIAM Journal on Computing 50 (2021), 1034--1102.
\url{https://arxiv.org/abs/1904.03213}.
\bibitem{DKW}
M. Dinitz, R. Krauthgamer and T. Wagner,
\emph{Towards Resistance Sparsifiers}, APPROX/RANDOM 2015.
\url{https://doi.org/10.4230/LIPIcs.APPROX-RANDOM.2015.738}.
\bibitem{CGPSS}
T. Chu, Y. Gao, R. Peng, S. Sachdeva, S. Sawlani and J. Wang,
\emph{Graph Sparsification, Spectral Sketches, and Faster
Resistance Computation via Short Cycle Decompositions},
SIAM Journal on Computing 52 (2023).
\url{https://arxiv.org/abs/1805.12051}.
\bibitem{Krivelevich}
M. Krivelevich,
\emph{On the Number of Hamilton Cycles in Pseudo-Random Graphs},
Electronic Journal of Combinatorics 19 (2012), P25.
\url{https://doi.org/10.37236/1177}.
\bibitem{BT}
J. Byrne and M. Tait,
\emph{Improved upper bounds on even-cycle creating Hamilton paths},
Discrete Mathematics 347 (2024), 114107.
\url{https://arxiv.org/abs/2304.02164}.
\bibitem{Muyesser}
A. M\"uyesser,
\emph{Hamiltonicity of mildly pseudorandom regular graphs},
2026. \url{https://arxiv.org/abs/2609.35766}.
\bibitem{CK}
B. Cuckler and J. Kahn, \emph{Hamiltonian cycles in Dirac graphs},
Combinatorica 29 (2009), 299--326.
\url{https://doi.org/10.1007/s00493-009-2360-2}.
\bibitem{Barvinok}
A. Barvinok, \emph{On Testing Hamiltonicity of Graphs}, 2014.
\url{https://sites.lsa.umich.edu/barvinok/wp-content/uploads/sites/1434/2025/05/hamil.pdf}.
\bibitem{CycleFactors}
M. Christoph, N. Dragani\'c, A. Gir\~ao, E. Hurley,
L. Michel and A. M\"uyesser,
\emph{Cycle-factors of regular graphs via entropy}, 2025,
revised 2026. \url{https://arxiv.org/abs/2507.19417}.
\bibitem{KS}
W. Klotz and T. Sander, \emph{Some properties of unitary Cayley graphs},
Electronic Journal of Combinatorics 14 (2007), R45.
\url{https://doi.org/10.37236/963}.
\bibitem{Basic}
M. Ba\v{s}i\'c,
\emph{Hamiltonian properties on a class of circulant interconnection networks},
Filomat 32 (2018), 71--85.
\url{https://doi.org/10.2298/FIL1801071B}.
\end{thebibliography}
\end{document}
